\documentclass[12pt]{article}
\usepackage[utf8]{inputenc}
\usepackage[T2A,T1]{fontenc}
\usepackage{amssymb}
\usepackage{amsmath}
\usepackage{graphicx}
\usepackage{tikz}

\usepackage[english,russian]{babel}
\usepackage[backend=biber, style=ieee]{biblatex}

\usetikzlibrary{arrows.meta, positioning, decorations.pathmorphing}

\DeclareMathOperator*{\flim}{flim} % Field limit

\textwidth 14cm
\textheight 21cm

\addbibresource{article_ru.bib}

\title{Несчётность множества континуум, шах и мат}
\author{Роман Олегович Коротков}

\begin{document}
 
\selectlanguage{russian}

\begin{flushleft}
{\bf УДК 510.223 }\\{\bf MSC 03E99 }
\end{flushleft}

{\let\newpage\relax\maketitle}

\begin{abstract}

В работе продолжается критический анализ границ применимости теоремы Кантора о несчётности континуума и связанного с ней диагонального аргумента, начатый в предыдущих статьях серии. Показывается, что классическое доказательство (включая обобщение Хессенберга) корректно лишь в рамках вполне упорядоченных структур при достаточно сильных предположениях об универсальности и тотальности бесконечности. Выделяются три принципиально различные алгебраические структуры описания бесконечных множеств - внутренняя (кардинальная), гибридная (ординальная) и внешняя (индексная, связанная с нумеросити, сюрреальными числами и гросс-объектами). Современная аксиоматическая интерпретация (ZFC, NBG) смешивает эти структуры, что приводит к нарушению изоморфизма с исходными конструкциями Кантора и к устойчивой условной непротиворечивости теории. Демонстрируются архитектурные недостатки современных аксиоматических систем - связанные с точками разрыва предикативного описания на предельных ординалах, ошибкой реификации и неоднозначностью отношения принадлежности, - из-за которых нарушается этот изоморфизм. Показано, что каждая из алгебраических структур требует собственного баланса предикативности и импредикативности, а их смешение в рамках единого понятия биекции ведёт к утрате онтологической поддержки операций (в частности, различия левого и правого сложения и умножения у Кантора). Выявляется тесная связь парадокса Рассела с теоремой Кантора, указывающая на единый структурный конфликт в основе обоих явлений. Обсуждаются феномен невыразимых объектов, роль импредикативности в замыкании алгебраических структур и связь с нестандартным анализом, сюрреальными числами и гросс-объектами. Устанавливается фундаментальная несовместимость глобальной замкнутости алгебры и тотальной упорядоченности её объектов. Восстановление исходной онтологии Кантора позволяет согласовать кардинальную, ординальную и индексную алгебры, уточнить статус континуума (как объекта, не являющегося множеством в строгом смысле) и наметить путь к устранению ряда парадоксов и методологических коллизий оснований математики.

{\it Ключевые слова: теория множеств; теорема Кантора; диагональный аргумент; потенциальная и актуальная бесконечность; предикативность и импредикативность; парадокс Рассела; кардинальная алгебра; ординальная алгебра; индексная алгебра; аксиома регулярности; трансфинитная рекурсия; биекция; взаимно однозначное соответствие; нестандартный анализ; сюрреальные числа; гросс-объекты; условная непротиворечивость; невыразимые объекты; реификация; основания математики.}

\end{abstract}

\selectlanguage{english}

\begin{abstract}

The paper continues the critical analysis of the limits of applicability of Cantor's theorem on the uncountability of the continuum and the associated diagonal argument, begun in the earlier articles of the series. It is shown that the classical proof (including Hessenberg's generalization) is valid only within well-ordered structures under sufficiently strong assumptions regarding the universality and totality of infinity. Three fundamentally distinct algebraic structures for describing infinite sets are identified: the internal (cardinal), the hybrid (ordinal), and the external (indexical, associated with numerosities, surreal numbers, and grossone objects). The modern axiomatic interpretation (ZFC, NBG) conflates these structures, leading to a breakdown of isomorphism with Cantor's original constructions and to a persistent conditional consistency of the theory. Architectural shortcomings of modern axiomatic systems—related to points of discontinuity in predicative description at limit ordinals, the fallacy of reification, and the ambiguity of the membership relation—are demonstrated as the cause of this breakdown of isomorphism. It is shown that each algebraic structure requires its own balance of predicativity and impredicativity, and that conflating them under a single notion of bijection leads to the loss of ontological support for operations (in particular, Cantor's distinction between left and right addition and multiplication). A close connection between Russell's paradox and Cantor's theorem is revealed, pointing to a single structural conflict underlying both phenomena. The phenomenon of inexpressible objects, the role of impredicativity in the closure of algebraic structures, and connections to non-standard analysis, surreal numbers, and grossone objects are discussed. A fundamental incompatibility between the global closure of an algebra and the total ordering of its objects is established. Restoring Cantor's original ontology makes it possible to reconcile cardinal, ordinal, and indexical algebras, clarify the status of the continuum (as an object that is not a set in the strict sense), and outline a path toward resolving a number of paradoxes and methodological collisions in the foundations of mathematics.

{\it Keywords: set theory; Cantor's theorem; diagonal argument; potential and actual infinity; predicativity and impredicativity; Russell's paradox; cardinal algebra; ordinal algebra; indexal algebra; axiom of regularity; transfinite recursion; bijection; one-to-one correspondence; non-standard analysis; surreal numbers; grossone objects; conditional consistency; inexpressible objects; reification; foundations of mathematics.}

\end{abstract}

\selectlanguage{russian}

Одним из самых мощных инструментов математики открывшимся человеку в 20 веке стала теория множеств, разработанная Георгом Кантором, способная обуздать и удерживать одну из самых противоречивых и парадоксальных сущностей - бесконечность, которая в момент из избегаемой абстракции превратилась в объект пристального изучения. Безусловно, наивно было предполагать, что понятие над которым великие мудрецы ломали головы тысячелетиями вдруг покорится без сопротивления. Оказалось, что бесконечность можно упаковать в некоторый контейнер - множество, изолировав и связанную с ней процедуру, не имеющую завершения. Однако, формализация таких контейнеров тут же столкнулась с большими трудностями, связанными с тем, что множества наследовали противоречивость и парадоксальность заключённой в них сущности. Даже сегодня взаимосвязь бесконечности, соответствующей ей никогда не завершающейся процедуры, с заключающим их контейнером остаётся недоисследованной, что при неаккуратном обращении приводит к притягательным иллюзиям, которые не просто развеять.

В статьях ''Несчётность множества континуум, шах'' и ''Несчётность множества континуум, ещё шах'' были подробно рассмотрены проблемы вызванные сочетанием бесконечных множеств и процедур, без аккуратного обозначения границ применимости выводов. Оказалось, что не всякую процедуру можно продолжать неограниченно долго, а тем более бесконечно, не делая при этом ряда допущений. Достоверность диагональных доказательств зависит от довольно вольных предположений относительно сочетания индуктивных эффектов с предельным поведением во время сравнения объектов. Рассматривая геометрические доказательства, мы обнаружили ещё больше утверждений, требующих аккуратного обозначения границ. Как и в случае с диагональным методом, нам необходимо отождествлять количество элементов двух бесконечных множеств на основании итеративного сопоставления конечных элементов. Данное наблюдение мы связали с двумя интерпретациями взаимно однозначного соответствия: сильной, рассматривающей такое соответствие как биекциею, преобразующую все элементы одного множества во все элементы другого, и слабой, замкнутой внутри множества, позволяющей определить лишь его класс эквивалентности, но не количественную тождественность. Кроме того, мы обнаружили, что все доказательства зависят от предположений об исходной структуре бесконечности, которая должна обладать достаточно сильными предикативными свойствами универсальности (возможности одновременного сопоставления элементов произвольного числа множеств) и тотальности (возможности теории говорить обо всех элементах совокупности без исключения), в то время как, предположение о гомогенности (неограниченной импредикативности) бесконечности делает доказательства внутренне неконсистентными. В случае геометрических доказательств неконсистентность условий в общем виде удаётся явно доказать, таким образом, они требуют предположений способных обосновать частичную консистентность построения, опирающихся на некоторую форму ограничения бесконечности. Более того, геометрические доказательства основаны на предположении представимости континуума множеством точек, которое также может оказаться не верным, особенно в современном контексте. Мы обнаружили, что все оригинальные доказательства непосредственно связаны с вполне упорядоченными структурами, в то время как современная интерпретация теоремы Кантора обобщает результат на произвольные неупорядоченные множества. Стоит отметить, что доказательства принадлежащие Кантору отличаются от современных онтологической и эпистемологической согласованностью, позволяющей легче анализировать области их применимости, и по всей видимости связаны, с противопоставлением потенциальной и актуальной бесконечности и синонимичными с ними импредикативностью и предикативностью, ключевую роль которых во всей этой истории мы отмечаем в своём повествовании. В оригинальных доказательствах прототипы альтернативных и мейнстрим теории соседствуют друг с другом и определяют две связанных картины, которые сильно разделились в современном контексте. Теорема Кантора зависит от исходных предположений о структуре бесконечности, а не определяет их, что существенно меняет современные акценты. Более того, как будет показано в настоящей работе, построения Кантора могут быть легко скорректированы, или скорее даже дополнены, как только мы формализуем их область применимости, что приводит к их полной согласованности с работами, которые в современном контексте считаются маргинальными. Мы проанализируем слабые места современных аксиоматических систем, которые окажутся некритическими, но когда мы разберём их и попытаемся заново собрать вокруг построений Кантора, то увидим где возникает расхождение и какие дополнительные предположения требуются чтобы сохранить современную интерпретацию, а какие приводят к альтернативным теориям.

Обратимся теперь к ключевому доказательству несчётности множества вещественных чисел, сформулированному в терминах современной теории множеств, утверждающему, что мощность множества всех подмножеств некоторого множества превосходит мощность этого множества. По всей видимости, его можно отнести к авторству Герхарда Хессенберга \cite{hessenberg}, поскольку именно у него мы впервые встречаем рассуждения о диагонализации биективного отображения определённого на множестве произвольной мощности, включая счётные множества. В то время как Кантор явно указывает, что данное доказательство можно распространить исключительно на {\it вполне определённые} многообразия (почему-то этот аккуратный термин Кантора не замечает мейнстрим) произвольной мощности, вероятно, прекрасно осознавая необходимость предикативного порядка на множестве, с которым неразрывно должны быть связаны функциональные отношения. Иными словами, вполне определённость также как и вполне упорядоченность указывают на сильные свойства универсальности и тотальности позволяющие биекции предикативно связывать множественные сущности. Формально он утверждает, что если аннотировать линейный континуум совокупностью вещественных чисел, то внутри алгебры определяющей данную совокупность как универсальный объект с тотальным порядком, не может существовать функция сопоставляющая данный объект с совокупностью двузначных функций доменом которых является данный объект. Вряд ли можно усомнится в данном утверждении, тем не менее, на некоторые неоднозначные детали можно указать. В частности, когда мы рассматриваем множество функций $f(x)$ из области вещественных чисел в значения $0$ и $1$, мы вынуждены принимать в расчёт аксиому непрерывности (требующую предикативности сечений), которая накладывает ограничения на построенную универсальную функцию $\varphi(x, y)$. Например, на её диагонали соотношение между группами точек со значение 1 и значением 0 определяются сечениями только двух типов, то есть к точке со значением 1, может прилегать лишь интервал со значение 0 (бесконечное множество точек) и наоборот, что с формальной точки зрения не даёт нам возможности утверждать, что построение диагонального элемента является следствием свойств множества определённой мощности, а не дополнительных ограничений принятых для него \cite{diagonal}. Анализ непрерывности имеет смысл, в том числе, потому, что эта аксиома играла центральную роль в исследовании геометрических доказательств в предыдущей статье, именно её конфликт с импредикативностью всюду плотности приводил к парадоксу, который Кантор использовал в своем доказательстве, а её комбинация с архимедовостью обеспечивает несчётность и уникальность поля вещественных чисел. Вместе с тем, искажение становится существенной частью формализма теории множеств, точнее формализм был создан в предположении верности искажённой интерпретации этого доказательства. Так или иначе, данный формализм в определённой степени оказался верен, по крайней мере, так полагают авторы, но наша интерпретация этого формализма распространила его за пределы области объектов которую он описывает, что связанно с неопределённостью в такой трансцендентной сущности как бесконечность, к пониманию которой, по всей видимости, приблизился только Кантор, к сожалению, уже не принимавший участие создании этого формализма. Данное обстоятельство в значительной мере усугубило кризис основ математики, появившись несколько позже периода творческой активности Кантора, подменив тщательно изолированную им область, требующую дополнительных предположений, на более обширную область, сохранив при этом неопределённость, которой дало разрастись до невероятных масштабов. В настоящее время это доказательство можно встретить в книгах по теории множеств \cite{cantor-common}, рассмотрим его более подробно.

Теорема: {\sl Никакое множество $X$ не равномощно множеству всех своих подмножеств.}

Доказательство: 

Предположим, что множество $X$ и множество всех его подмножеств $\mathcal{P}(X)$ равномощны. Пусть $\varphi$ - взаимно однозначное соответствие между ними. Рассмотрим те элементы $x \in X$, которые не принадлежат соответствующему им подмножеству. Образованное ими множество обозначим как $Z$:
\begin{equation}
Z = \{ x \in X \mid x \notin \varphi(x) \}
\end{equation}

Докажем теперь, что подмножество $Z$ не соответствует никакому элементу множества $X$. Пусть это не так и $Z$ соответствует некоторому элементу $z \in X$, то есть $Z = \varphi(z)$. Тогда
\begin{equation}
z \in Z \Leftrightarrow z \notin \varphi(z) \Leftrightarrow z \notin Z
\end{equation}

Первая равносильность вытекает из определения множества $Z$, как множества элементов, которые не принадлежат соответствующему им подмножеству, а вторая равносильность следует из предположения, что $Z = \varphi(z)$. Полученное противоречие показывает, что $Z$ действительно ничему не соответствует, так что соответствие $\varphi$ не взаимно однозначно. Теорема доказана.

Во-первых, заметим, что в рамках данного доказательства не осталось места классической концепции потенциальной бесконечности, которая долгое время главенствовала в математике надёжно предохраняя её от парадоксальных заключений. Удостоверимся, что приведённые выше рассуждения не совместимы с понятием потенциальной бесконечности на примере натуральных чисел. 

Теорема: {\sl Никакое конечное и актуально бесконечное множество $X$ не равномощно множеству всех своих подмножеств.}

Доказательство: 

Предположим, что мы можем распространить предыдущее доказательство на потенциально бесконечное множество натуральных чисел $\mathbb{N}$. Как и в оригинальной теореме, предполагаем, взаимно однозначное отображение его на $\mathcal{P}(\mathbb{N})$, с помощью соответствия $\varphi$, то есть:
\begin{align}
\begin{split}
1 &\leftrightarrow (x^1_1, x^1_2, \dots, x^1_{k^1})\\
2 &\leftrightarrow  (x^2_1, x^2_2, \dots, x^2_{k^2})\\
\dots &\leftrightarrow \dots\\
n &\leftrightarrow (x^n_1, x^n_2, \dots, x^n_{k^n})\\
\dots &\leftrightarrow \dots
\label{constr:n-function}
\end{split}
\end{align}

Рассмотрим последовательность конечных множеств $Z_i$, состоящих из элементов не превосходящих $i$, которые не принадлежат соответствующим им подмножествам. То есть, $Z_1$ пусто, если $1 \in \varphi(1)$, и $Z_1 = \{1\}$ в противном случае, $Z_2 = Z_1 \bigcup \emptyset$ если $2 \in \varphi(2)$, и $Z_2 = Z_1 \bigcup \{2\}$. Точно также, дойдя до пары $n - 1$, имеем $Z_n = Z_{n - 1} \bigcup \emptyset$ если $n \in \varphi(n)$, и $Z_n = Z_{n - 1} \bigcup \{n\}$. Получилась потенциально бесконечная последовательность множеств, по построению очень похожих на множество $Z$ из доказательства. Попробуем теперь разобраться, что же такое множество $Z$ в случае натуральных чисел и как оно связано с последовательностью множеств $\{Z_i\}$. Некоторое множество может быть связано с последовательностью множеств, в частности, двумя следующими способами: либо оно содержится в последовательности в качестве элемента, либо не содержится (закон исключённого третьего).

\begin{enumerate}
    \item Последовательность $\{Z_i\}$ содержит множество $Z$, то есть существует некоторый индекс $m$, такой, что $Z = Z_m$. По построению множество $Z_m$ полностью определено первыми $m$ натуральными числами. Поскольку для любого $m \in \mathbb{N}$ всегда есть $m + 1 > m$, есть и неограниченное множество натуральных чисел превосходящих $m$, значит может существовать некоторый элемент с индексом $w > m$, такой что $Z_m = \varphi(w)$ и мы не в состоянии прийти к противоречию. В частности если $\varphi$ таково, что содержит последовательно все комбинации длины $1$ из элемента $1$, затем все комбинации длины $2$ из элементов $\{1, 2\}$, все комбинации длины $3$ из элементов $\{1, 2, 3\}$ и так далее, то $w$ непременно входит в первые $2^m$ элемента.
    \item Множество $Z$ не совпадает ни с одним из множеств в последовательности $\{Z_i\}$. Однако легко убедиться, что такой случай представиться не может. Действительно, поскольку последовательность $\{Z_i\}$ монотонна, то есть может только пополняться новыми элементами, и состоит исключительно из конечных множеств, с которыми мы легко можем поэлементно сравнивать любое другое множество, то могут реализоваться лишь 2 возможности приводящие к тому, что множество $Z$ не содержится в последовательности: если на каком-то шаге $k$ множество $Z_k$ пополнилось элементом, которого нет в множестве $Z$, либо наоборот множество $Z$ содержит некоторый элемент, который не был добавлен ни на одном из шагов, то есть отсутствует в последовательности $Z_i$. Первая возможность не может реализоваться, так как $Z_k$ может пополниться только элементом $k$ тогда и только тогда когда $k \notin \varphi(k)$, из чего следует $k \in Z$, что противоречит предположению. Остаётся единственный вариант, что $k \in Z$ и ни одно из множеств $Z_i$ не содержит $k$, но из того, что $k \in \mathbb{N}$ и $k \in Z$ однозначно следует, что $k \notin \varphi(k)$, но тогда $k$ по построению входит в $Z_k$, снова приходим к противоречию с предположением. Заметим, что остаётся ещё случай $k \notin \mathbb{N}$, но он запрещён потенциальной бесконечностью. 
\end{enumerate}

Таким образом, требуемое множество $Z$ не может существовать. Теорема доказана.

Заметим что, чтобы была возможность прийти к противоречию в доказательстве множество $Z$ должно содержать некоторый элемент или элементы, которые не является натуральными числами. Судя по всему, среди таких элементов и содержится то, что мы называем актуальной бесконечностью, которая позволяет множеству $Z$ стать бесконечным и делает дальнейшие рассуждения, по крайней мере, обоснованными. Иными словами, если мы упорядочим все подмножества множества натуральных чисел, то для обозначения некоторых из них нам неизбежно потребуются трансфинитные порядковые числа, не являющиеся натуральными. Но для этого нам требуется отказаться от импредикативности (потенциальной бесконечности) натуральных чисел, которая обеспечивала их замкнутость, и принять существование некоторой трансфинитной предикативной иерархии, которую можно замкнуть только введя новый тип совокупностей - классы. На данном этапе может возникнуть вопрос, насколько данные рассуждения можно считать корректными, в частности, множество $\mathbb{N}$ не определено в контексте потенциальной бесконечности, хоть и ясно, что оно здесь не играет роли бесконечного множества натуральных чисел определяемого в теории множеств, а скорее является классом объектов полученных из одного построения, в связи с чем, использование отношения принадлежности является не совсем законным. В то же время, оно выражает наши априорные знания о числах и потенциальной бесконечности, а значит, с позиций интуиционизма заслуживает к себе такого же серьёзного отношения, как и любое другое математическое доказательство. Впрочем, данный тезис интуиционистов вряд ли был услышан и в достаточной мере осознан, поскольку зачастую атаки против альтернативных интерпретаций множественных построений направлены на математическую форму, а не на математическое содержание, игнорируя простую истину сформулированную Ричардом Фейнманом: нужно слушать не то, что человек говорит, а то что он имеет ввиду. С формальной же точки зрения, вопрос корректности данного доказательства упирается в систему предположений и определений в их естественных границах, исследование которых и станет нашей задачей. 

Как бы то ни было, обратим внимание на любопытный феномен: по какой-то причине, мы постоянно сталкиваемся с фактом того, что существует математическая картина (в настоящее время альтернативная), где утверждения справедливые для потенциальной бесконечности (даже в интуитивной форме) при тщательном рассмотрении справедливы и для актуальной бесконечности, где последняя становится неподвижной точкой, некоторым абсолютно рефлексивным объектом. Данное наблюдение не согласуется с тем, что в рамках доказательства бесконечные объекты появляются в множестве $Z$, которое по построению является подмножеством множества натуральных чисел, определяемого как множество всех конечных натуральных чисел. Что указывает на некоторую неопределённость и необходимость в тщательной проработке концепции бесконечных объектов, ведь в рамках современной теории множеств бесконечные объекты приобретают свойства, которые мы не могли наблюдать в потенциально бесконечной математической реальности. Например, доказательство Евклида об отсутствии максимального натурального числа верно в случае потенциальной бесконечности, но когда мы декларируем существование бесконечного множества натуральных чисел в современных аксиоматиках, мы расширяем границы применимости данного доказательства, при этом недостаточно чётко осознавая до каких пределов допустимо такое расширение. Границы применимости имеют все без исключения математические концепции, и именно их аккуратная формализация делает математику строгой. Ключевой задачей настоящего исследования является демонстрация пределов области применимости некоторых утверждений теории множеств, простирающихся значительно дальше достоверных границ, что приводит к существенным трудностям. В частности, в истории математики полно примеров определения максимального натурального числа, которые в современном контексте становятся неотъемлемой частью алгебраических структур, таких как нестандартные натуральные числе \cite{infinitesimal-nonstandard-numbers} ($\dot{\omega} - 1$) (в том числе нумеросити \cite{numerosities}), сюрреальные числа \cite{conway-surreal-numbers} ($\dot{\omega} - 1$) или гросс-объекты \cite{grossone} ($\raisebox{.5pt}{\textcircled{\scalebox{0.8}{$\omega$}}} - 1$). Здесь и далее, мы используем обозначения $\dot{\omega}$, $\dot{\omega_1}$ для нестандартных натуральных чисел и омнифических целых чисел и обозначения $\raisebox{.5pt}{\textcircled{\scalebox{0.8}{$\omega$}}}$, $\raisebox{.5pt}{\textcircled{\scalebox{0.8}{$\omega_1$}}}$ для гросс-объектов, чтобы подчеркнуть их различную с ординалами природу и избежать неоднозначности. Авторы полагают, что данные объекты совместимы с современными аксиоматиками гораздо естественней и проще, чем представляется мейнстриму в нестандартном анализе, и являются ещё одним классом числоподобных объектов (наряду с кардиналами и ординалами), ускользнувшим от внимания Кантора (или подразумевавшийся в следующих работах), вследствие чего неучитываемый при аксиоматизации теории множеств, что сейчас усложняет их возвращение в математику.

Напомним, что мы проводим чёткое различие между понятиями биекции и взаимно однозначного соответствия. Биекция между двумя множествами гарантирует, что образом одного множества является в точности другое множество. Взаимно однозначное соответствие, в свою очередь, лишь обеспечивает возможность продолжения процедуры сопоставления полностью как для одного, так и для другого множества. Сильная интерпретация полностью предикативна, в то время как слабая - нет, но допускает некоторые ограничения импредикативности. Такими ограничениями импредикативности выступают предполагаемые свойства бесконечных объектов: универсальность или воспроизводимость, обеспечивающая одновременное сопоставление неограниченных совокупностей объектов, что можно рассматривать как существования эталона, и тотальность, заключающаяся в способности делать предикативные утверждения о всех элементах, ассоциированная с порядковыми структурами. Для окончательного прояснения ситуации авторы полагают, что Кантор фактически использовал две различные слабые интерпретации взаимно однозначного соответствия. Первая из них связана с процедурами перечисления и смещения элементов и применяется в кардинальной арифметике, где есть универсальность, но нет тотальности. Вторая связана с процедурой конкатенации (сцепления) упорядоченных отрезков множеств и используется в арифметике порядковых чисел, где импредикативность ограничена также тотальностью. Обе интерпретации требуют классов эквивалентности, обладающих свойством подобия своим частям, и различаются явно, проявляя себя прежде всего через алгебраическую структуру соответствующих чисел, с которой они находятся в неразрывной связи. Судя по всему, сам Кантор осознавал это различие, что подтверждается введением им существенно разных систем понятий для кардинальных и порядковых чисел \cite{cantor-definitions}. Современным языком это можно выразить следующим образом: Кантор построил не одну, а две полноценные алгебры - одну для кардинальных чисел и другую для порядковых, каждая из которых включает свои функциональные отношения (в том числе взаимно однозначные соответствия). 

Алгебраическую структуру кардинальных чисел вместе с арифметическими операциями, основанными на процедурах перечисления и смещения элементов, мы называем {\it внутренней}. В данном контексте структура множеств рассматривается как их внешнее свойство расширяющее понятие множества, которое не влияет на результат операций (можно добавить элемент в любое множество не заботясь о его структуре). Соответствие здесь называется эквивалентностью, а сами кардинальные числа интерпретируются как мощности множеств. Алгебраическую структуру порядковых чисел мы называем {\it гибридной}, поскольку арифметические операции, основанные на конкатенации упорядоченных отрезков, частично отражаются на общей структуре результата. Соответствие таких структур называется подобием, а сами они порядковыми типами. Наконец, алгебраическую структуру, которую Кантор не выделил явно, связанную с натуральными суммой и произведением \cite{hessenberg-sum}, сюрреальными числами, нестандартными натуральными числами и гросс-объектами, которая, как выяснится, также обладает строгим процедурным воплощением в форме операций над множествами, мы называем {\it внешней}. В этой структуре нетривиальные операции могут приводить к изменению организации объектов (результат добавления элемента зависит от того, как он встраивается в структуру). Отметим, что такое деление не имеет никакого отношения к делению на внутренние и внешние объекты в нестандартном анализе, где оно является синтаксическим и связано с классификацией формул, в то время как, в нашем случае речь идёт о семантике связи между операциями алгебры и структуры множеств. В настоящее время алгебраическая структура кардинальных чисел не имеет единого общепринятого названия. В области конечных кардиналов она представляет собой аддитивно-мультипликативное кольцо, а в трансфинитной области - алгебраическую систему, в которой максимальный операнд поглощает минимальный. Гибридная структура, тщательно выстроенная Кантором, состоит из аддитивного и мультипликативного моноидов, связанных гомоморфизмом. Примечательно, что благодаря явному акценту на этот гомоморфизм Кантору удалось избежать неопределённостей с определением пределов, характерных для классического анализа (что будет рассмотрено далее). Внешняя структура представляет собой кольцо - комбинацию аддитивной и мультипликативной групп, также связанных гомоморфизмом. По мнению авторов, именно с этой структурой логично связывать понятие биекции в строгом смысле, поскольку она определяет наиболее обширный класс объектов, классами эквивалентности которых являются порядковые и кардинальные числа. 

Во внутренней структуре понятие предела в классическом смысле отсутствует; существует лишь супремум, который, как правило, недостижим с помощью арифметических операций. Точнее говоря, предел здесь можно определить только как супремум. Напротив, арифметические операции гибридной и особенно внешней структур позволяют естественным образом определить понятие предела как элемента некоторой универсальной структуры. Мы также не рассматриваем в данной статье связь индексных чисел с аксиомой Евклида - целое больше своих частей, в данном контексте нас интересует исключительно морфология этих объектов, связь которой с ограничениями Евклида будет рассмотрена в следующей статье. Структурного определения внутренних и внешних аспектов, предложенного в данной статье, на текущий момент достаточно. Благодаря выявленным закономерностям в дальнейшем мы планируем дать этим понятиям более фундаментальное и строгое определение. Интуитивно различие между структурами проявляется следующим образом. В кардинальной арифметике объединение двух множеств рассматривает их как единое целое, независимо от внутренней структуры. Допустимо даже изменение структуры (например, ''выталкивание'' элемента ''вверх'' до бесконечности), не влияющее на кардинальный тип, то есть понятие всех объектов несёт универсальный характер по включению элементов, но тотальность отсутствует, так как система лишена структурной информации позволяющей однозначно распространить некоторое утверждения с конечного набора на все элементы. В ординальной же арифметике подобное ''всплытие'' элемента изменяет порядковый тип (например, превращает $\omega$ в $\omega + 1$), вследствие тотального вполне упорядочения. Важно отметить, что аддитивно-мультипликативный гомоморфизм играет ключевую роль во всех трёх алгебраических структурах, обеспечивая их внутреннюю согласованность и регулярность. К настоящему времени накоплено достаточно материала, чтобы чётко идентифицировать и разделить указанные структуры. Во времена Кантора это было существенно сложнее из-за философского неприятия актуальной бесконечности и её противопоставления потенциальной. Все эти концепции оказались тесно переплетены: потенциальная бесконечность связана преимущественно с внутренним строением, актуальная - с внешним. Порядковые числа, допускающие трансфинитную индукцию и рекурсию, указывают на связь с потенциальной бесконечностью, в то время как строгое определение пределов свидетельствует об их связи с актуальной бесконечностью. Каждый тип числоподобных объектов определяется через собственный баланс импредикативности потенциальной бесконечности и предикативных свойств актуальной. Однако, судя по всему, из-за отсутствия языка для описания данных тонкостей, в процессе формализации теория Кантора была существенно упрощена и три различные алгебраические структуры стали описываться одним и тем же понятием биекции. Которого оказывается явно недостаточно для адекватной характеристики всего многообразия отношений между множествами, что вынуждает делать дополнительные предположения, существенно сужающие класс рассматриваемых объектов.

Чтобы избежать необоснованной критики, необходимо сразу уточнить, что конкретно мы подвергаем сомнению в приведённом выше доказательстве и для чего собственно мы тщательно исследовали исходные формулировки доказательств Кантора, найдя в них ряд неоднозначностей. Действительно, невозможно себе вообразить, чтобы в настолько ключевом доказательстве определяющем целую область математики, являющуюся её фундаментом, перепроверенному множество раз, содержалась ошибка. Усилиями многих выдающихся математиков фундамент теории множеств стал невероятно прочным, а само по себе доказательство логически безупречно, что делает поиск уязвимостей в нём задачей на грани возможного. Тем не менее, именно это доказательство является краеугольным камнем и ключом, отмыкающим двери к пониманию источников противоречий и парадоксов связанных с теорией множеств, и именно поэтому оно привлекает внимание стольких исследователей, пытающихся понять, что же мы упускаем. Мы уже показали, что доказательство теоремы Кантора не имеет отношения к идее потенциальной бесконечности, и требует полноценного существования бесконечных объектов, что по смыслу является выражением того, что обычно подразумевается под термином актуальная бесконечность. Но и в случае актуально бесконечных построений, общепринятая интерпретация нуждается в ряде дополнительных предположений, формализация которых открывает двери для интересных наблюдение и альтернативных концепций. В рамках современных аксиоматик, если мы считаем равномощными множества, между которыми можно установить биективное отображение, то теорема абсолютно верна и никаких вопросов к доказательству не возникает. В то же время, из-за смешения понятий, числоподобные объекты Кантора и числоподобные объекты современных теорий представляют собой разные сущности и, судя по всему, последние из них нуждаются в переосмыслении, так как существенно расходятся в интерпретации со своими прототипами. В итоге, доказательство прекрасно работает в области внешних структур, где у нас есть строгое определение понятия предела, однако, оно неприменимо в области внутренних структур, где сопоставление - это процесс, не имеющий предела, но ограниченный супремумом, которого он никогда не достигнет. Острая математическая интуиция не может обойти неопределённостей связанных с определением границ применимости данного доказательства, что и обуславливает огромное количество работ посвящённых его опровержению. Помимо всего прочего, с данным доказательством тесно связан неимоверно любопытный парадокс Рассела, указывающий на некоторую фундаментальную проблему в описании актуально бесконечных объектов, которая судя по всему и лежит в самом сердце кризиса оснований математики, не позволяя математическим теориям обрести полную мощь.

Давайте попытаемся исследовать общие рассуждения о множестве всех подмножеств в контексте обнаруженных нами уязвимостей в геометрических доказательствах, где противоречие возникает вследствие неконсистентности построения, а попытка обосновать его, как следствие предположения, ссылаясь на индукцию приводит к незаурядным результатам, наталкивающимся на строгую алгебраическую структуру континуума. Это вынуждает нас вводить ограничение на последовательность интервалов, например, постулируя архимедовость, в том числе в обобщённом смысле, где натуральное число становится ординалом непревосходящим определённой мощности, что позволяет сделать вывод только о том, что континуум больше данной мощности, но не зафиксировать эту мощность. Рассмотрим произвольное множество $X \neq \emptyset$, в соответствии с аксиомой степени существуют множества $\mathcal{P}(X)$ и $\mathcal{P}(\mathcal{P}(X))$. Из множества $\mathcal{P}(X)$ существует сюръекция в $X$, $f : \mathcal{P}(X) \to X$, достаточно сопоставить синглтоны своим элементам, а всем остальным множествам $\mathcal{P}(X)$ сопоставить произвольный элемент, но можно воспользоваться и аксиомой выбора. Пусть $F$ - множество пар сюръекции $f$. Определим отображение $g : F \to \mathcal{P}(\mathcal{P}(X))$ следующим образом: $\forall p \in \mathcal{P}(X) \; g((p, f(p))) := \{ q \in \mathcal{P}(X) \mid f(p) \in q \oplus f(p) \in p \}$, то есть каждой паре $(p, f(p))$ мы сопоставляем правильное подмножество $g(p) \subset \mathcal{P}(\mathcal{P}(X))$, так как $p \notin g(p)$, все элементы которого отличаются от $p$ принадлежностью элемента $f(p)$, это можно сделать используя схему аксиом спецификации для каждой пары. При этом $g(p) \neq \emptyset$ для всех элементов $p$, в самом деле, если $f(p) \in p$, то $\emptyset \in g(p)$, в противном случае $\{ f(p) \} \in g(p)$. Множества пар отображения $g$ обозначим как $G$. По теореме 1 из статьи ''Несчётность множества континуум, ещё шах'' такого отображения не существует, если мы накладываем на него условие, что пересечение всех $\bigcap\limits_{p \in \mathcal{P}(X)} g(p) \neq \emptyset$. В рамках настоящего построения, это даёт нам следующую теорему:

Теорема 2: {\sl Никакое множество $X$ не позволяет определить сюрьекцию $f : \mathcal{P}(X) \to X$, удовлетворяющую условию: $f$: $\forall x \in X \; (\forall p \in \mathcal{P}(X) : f(p) \neq x \oplus x \in p) \oplus (\forall p \in \mathcal{P}(X) : f(p) \neq x \oplus x \notin p)$.}

Иными словами, каким бы ни было множество $X$ сюрьекция $f$ из множества степени $\mathcal{P}(X)$ такая что, либо все прообразы произвольного элемента $x$ содержат этот элемент, либо все не содержат, не существует.

Доказательство: 

Предположим, что требуемая сюрьекция существует. Воспользуемся схемой аксиом спецификации и рассмотрим множество $Z = \{ f(p) \mid f(p) \in \bigcap g(p) \} = \{ x \in X \mid \forall p \in \mathcal{P}(X) : f(p) \neq x \oplus x \in \bigcap g(p) \}$. Поскольку $Z \in \mathcal{P}(X)$, то существует некоторый элемент $z$ множества $X$, в который оно отображается, то есть $f(Z) = z$. Для данной пары также определено значение функции $g$: $g(Z) = \{ q \in \mathcal{P}(X) \mid z \in q \oplus z \in Z \}$. Тогда, если $z \in Z$, то $z \in \bigcap g(p)$ для всех $p$, которые отображаются в $z$, в том числе и для $p = Z$, значит $z \in \bigcap g(Z)$, но если $z$ принадлежит всем элементам $g(Z)$, то по определению $g$ имеем $z \notin Z$, таким образом, $z \in Z \Rightarrow z \notin Z$. В обратную сторону, если $z \notin Z$, тогда существуют такие $p$, для которых $z \not \bigcap g(p)$, если $z \notin \bigcap g(Z)$, то поскольку $g(Z) \neq \emptyset$, есть хотя бы одно множество $q \in g(Z)$, для которого $z \notin q$, что в силу определения $g$ влечёт $z \in Z$. Случай $z \in \bigcap g(Z)$ был бы корректным выходом из ситуации, но условие на сюръекцию $f$ запрещает данный случай, следовательно, $z \in Z \Leftrightarrow z \notin Z$, что опровергает наше предположение. Теорема доказана.

Таким образом, сюръекция $f$ с подобным условием невозможна, более того, она невозможна независимо от нашего предположения о биекции, то есть, о том что $f(p) = x$ единственно. Тем не менее, в контексте аксиоматики теории множеств, наше предположение в силу схемы аксиом выделения влечёт существование сюръекции $f$, что приводит к противоречию. Усомниться в схеме аксиом выделения можно \cite{diagonal-antinomy}, но проблема в том, что аксиоматическая теория множеств создавалась с той целью, чтобы ограничить противоречивость данной схемы аксиом, именно для этого универсум теории множеств был кардинально сужен до регулярной кумулятивной иерархии порождаемой классом порядковых чисел, существующим вследствие аксиомы бесконечности. Более того, при построении аксиоматик теории множеств, акцент делался именно на непротиворечивости диагонального построения и нейтрализации связанного с ним парадокса Рассела \cite{russell-paradox}, что оказалось возможным, благодаря тому, что формула Рассела ($x \notin x$) может быть непротиворечивой в некоторых универсумах, одним из которых и выступает кумулятивная иерархия. Однако, фокусируясь на столь сложном объекте, обнаруженном ещё в античности в форме парадокса лжеца, мы обошли вниманием ключевой вопрос: является ли формула Рассела уникальной в своём роде, или подобными свойствами могут обладать другие формулы? В контексте теоремы Кантора, схема выделения в универсумах современных теорий множеств консистентна с другими предположениями и не приводит к противоречиям, которые могут возникать вне жёстких рамок аксиоматики. Поэтому единственный вывод который мы можем сделать в данном контексте, заключается в том, что теорема корректна, она связана с невозможностью более общего построения, но предположение в доказательстве однозначно влечёт существование невозможного построения. Более того, поскольку по аксиоме степени данный объект существует как множество, у нас, в отличии от геометрических доказательств, появляется возможность зафиксировать эту мощность. Таким образом, проблема которая возникала в первоначальных доказательствах несчётности множества вещественных чисел была полностью устранена в рамках современных аксиоматик, в то же время, её артефакты до сих пор проявляются в виде обнаруженных нами уязвимостей в рамках строгих структур анализа, указывая на то, что если в теории множеств и есть уязвимости, то они связаны с аксиомами бесконечности и выбора обуславливающими индукцию требуемую в геометрических построениях.

Чтобы рассмотреть ещё одно затруднение, необходимо вернутся к самому началу, а именно к наивной теории множеств. Представим некоторую вселенную состоящую из двух типов сущностей: объектов, которыми могут быть любые предметы, и множеств, содержащих как объекты так и сами множества. Если мы хотим узнать каким отношением связаны объекты и множества, мы можем попробовать сопоставить каждому множеству уникальный объект, однако, у нас возникнут трудности с объектом, который соответствовал бы множеству объектов, которые не содержатся во множестве, которому они сопоставлены. Очевидно, что сущностей не меньше чем множеств, поскольку любое множество является сущностью, а кроме того существуют атомарные сущности, такие как молекулы, яблоки, камни, планеты и прочие, которые называют объектами. Попробуем сопоставить каждому множеству некоторую сущность, но у нас снова возникает загвоздка, с множеством сущностей, которые не содержатся в соответствующих им множествах, ведь по определению это множество не может содержать сущность, которой оно соответствует, но в таком случае сопоставленная ему сущность должна удовлетворять критерию сущностей в него входящих, а следовательно должна в нем присутствовать. Данный парадокс можно формулировать по-разному, можно попытаться дать определение всем множествам и рассмотреть множество определений, которые не содержатся в определяемых ими множествах, либо попытаться назвать или обозначить все множества и так далее. Он возникает в том случае, если областью определения некоторого отображения $\psi$ являются множества, а областью значений - объекты, содержащиеся в этих множествах, тогда множество объектов $O = \{ \psi(s) \notin s \mid \forall s_1, s_2 \; \psi(s_1) \neq \psi(s_2) \}$ должно содержать и не содержать соответствующий ему объект одновременно. Требование уникальности значений совместно с использованием абсолютно всех объектов в качестве значений (если вдруг нам не требуются все объекты, можно рассмотреть образ множества множеств во множестве объектов), равносильно тому, что отображение взаимно однозначно, тогда заменив переменную $x = \psi(s)$, можно записать $O = \{ x \notin \psi^{-1}(x) \}$, или, переобозначив $\psi^{-1} = \varphi$, имеем $O = \{ x \notin \varphi(x) \}$. Его частный случай, когда используется тождественное отображение $\psi(x) = \varphi(x) = x$ (множеству сопоставляется оно само) широко известен как парадокс Рассела. 

По какой-то причине, в начале прошлого века было решено, что только отображение множества множеств самого в себя вызывает затруднения, вероятно, ещё отображение множества множеств в множество сущностей не имеет смысла, поскольку множество множеств является его подмножеством, а вот с множеством объектов совсем другая история, посчитали, что мощность множества объектов меньше чем мощность множеств из них составленных. Хорошо, значит парадокс для одних отображений мы считаем противоречием и всячески избегаем, а для других используем в доказательствах. Но присмотримся в таком случае к тому, что в действительности происходит в нашем доказательстве. 

Теорема 3: {\sl Биекция между множеством и множеством его подмножеств эквивалентна антиномии Рассела.}

Доказательство: 

Мы предполагаем взаимно однозначное отображение $\varphi$ между множеством $X$ и множеством всех его подмножеств $\mathcal{P}(X)$, тогда множество подмножеств $X$ можно представить как $\mathcal{P}(X) = \{ p \mid p \subset X \} = \{ \{ \varphi^{-1}(p) \mid p \in u \} \mid u \subset \mathcal{P}(X) \}$. Рассмотрим отображение $\psi$ множества $\mathcal{P}(X)$ на другое множество $S(X) = \{ \psi(p) \mid p \in \mathcal{P}(X) \}$, где $\psi(p) = \{ \vartheta(x) \mid x \in p \} = \{ \vartheta(\varphi^{-1}(v)) \mid v \in u \; \wedge \; u \subset \mathcal{P}(X) \}$ некоторое преобразование множества, такое что $\vartheta(\varphi^{-1}(p)) = \psi(p)$, то есть $S(X) = \{ \{ \psi(p) \mid p \in u \} \mid u \subset \mathcal{P}(X) \}$. Благодаря нашему предположению, такое отображение полностью определено на множестве $\mathcal{P}(X)$, кроме того нет ничего препятствующего уникальности его значений, ведь любой образ $u$ является множеством образов $p$, и если уникальны исходные элементы то их образы могут быть уникальными, а множество уникальных образов уникально (идентифицируемо с помощью аксиомы экстенсиональности), значит ничто не препятствует тому, что некоторое подмножество таких отображений является взаимно однозначным. Будем считать, что $\psi$ - одно из таких взаимно однозначных отображений, тогда $S(X) = \{ \{ t \mid \psi^{-1}(t) \in \{ \psi^{-1}(r) \mid r \in s \} \} \mid s \subset S(X) \}$ и заметив, что из определения $\psi(p) = \{ \psi(v) \mid v \in u \}$, в силу $v \in u = \psi^{-1}(t) \in \{ \psi^{-1}(r) \mid r \in s \}$ следует $t \in s$, получаем $S(X) = \{ \{ t \mid t \in s \} \mid s \subset S(X) \}$. Иными словами полученное множество является множеством образов $\mathcal{P}(X)$ содержащих исключительно образы $\mathcal{P}(X)$. Теорема доказана.

Стоит уточнить, что наши действия вполне корректны, если мы исключаем аксиому регулярности из теории (вследсвии которой не любое подмножество $S(X)$ может являться множеством). Так с точки зрения аксиоматики Цермело - Френкеля, множество $S(X)$ является следствием аксиомы схемы преобразования \cite{fraenkel-axioms}, они также корректны в аксиоматике фон Неймана - Бернайса - Гёделя, где это следствие аксиомы ограничения мощности \cite{von-neumann-axioms} (если существует взаимно однозначное отображение $S(X)$ в $V$, допустим $\rho$, то $\mathcal{P}(X)$ - также не может быть множеством, ведь $\psi \circ \rho$ взаимно однозначно отображает его в $V$), и вероятно, во всякой аксиоматике допускающей определённую степень свободы в преобразовании множеств такое построение допустимо. Значит, предполагая взаимно однозначное отображение $\varphi$ между множеством $X$ и множеством всех его подмножеств $\mathcal{P}(X)$ мы одновременно предполагаем существование множества $S(X)$ связанного с обоими множествами взаимно однозначным отображением, но образ множества $Z$ в множестве $S(X)$ воспроизводит парадокс Рассела в чистом виде, противоречие на уровне прообразов возникает одновременно с противоречием на уровне образов, что, вообще говоря, логично. Обратим внимание, что нам даже не потребовалось множество всех множеств, факт свидетельствующий о том, что ключевую роль в парадоксе Рассела играют вовсе не размеры множества, а алгебраическая структура, которой мы его описываем, чему мы найдём и другие подтверждения. Связь противоречия возникающего в доказательстве с парадоксом Рассела, оказывается весьма глубокой и указывает нам на фундаментальный конфликт между предикативностью объекта как части домена и импредикативностью самого домена без фиксации его структуры, конфликт потенциальной и актуальной бесконечности сформулированный в апориях Зенона. Иными словами, мы допускаем ошибку когда разбиваем парадокс на части, одну из которых продолжаем считать антиномией, а вторую используем в доказательстве теоремы Кантора. В самом деле, если противоречие является неотъемлемой частью вселенной или нашего восприятия, то изолировав его источник, сделав его правомерным математическим объектом, мы открыли бы доступ к исследованию данного феномена и его природы, в свою очередь в рамках доказательства, нам было бы достаточно уточнить область его применимости. Если же противоречие связано с какой-то фундаментальной ошибкой описания реальности, то её идентификация привела бы к устранению коллизий как в случае парадокса Рассела, так и в случае доказательства. Однако, вместо этого, одну часть мы выводим за пределы описываемого универсума, а вторую кладём в самое сердце теории, ставшей фундаментом математики, что, судя по всему, и приводит к началу кризиса её оснований, поскольку источник противоречий оказался не полностью изолированным. 

Обратим внимание, что есть некоторая фундаментальная странность во множествах, которым позволено содержать самих себя. Воспроизводство парадокса Рассела на образе биекции произвольного множества с множеством всех своих подмножеств, продемонстрированное нами выше, позволяет построить любопытные объекты, используя аппарат теории множеств:
\begin{center}
\begin{tabular}{c c c r r r r r r r c}
$P$ & $=$ & $\{$ & $p_1$, & $p_2$, & $p_3$, & $p_4$, & $\dots$, & $p_n$, & $\dots$ & $\}$\\
$p_1 = \{\}$ & $=$  &  & 0, & 0, & 0, & 0, & $\dots$, & 0, & $\dots$ & \\
$p_2 = \{ p_1 \}$ & $=$  &  & 1, & 0, & 0, & 0, & $\dots$, & 0, & $\dots$ & \\
$p_3 = \{ p_2 \}$ & $=$  &  & 0, & 1, & 0, & 0, & $\dots$, & 0, & $\dots$ & \\
$p_4 = \{ p_1, p_2 \}$ & $=$  &  & 1, & 1, & 0, & 0, & $\dots$, & 0, & $\dots$ & \\
$p_5 = \{ p_3 \}$ & $=$  &  & 0, & 0, & 1, & 0, & $\dots$, & 0, & $\dots$ & \\
$p_6 = \{ p_1, p_3 \}$ & $=$  &  & 1, & 0, & 1, & 0, & $\dots$, & 0, & $\dots$ & \\
$p_7 = \{ p_2, p_3 \}$ & $=$  &  & 0, & 1, & 1, & 0, & $\dots$, & 0, & $\dots$ & \\
$p_8 = \{ p_1, p_2, p_3 \}$ & $=$  &  & 1, & 1, & 1, & 0, & $\dots$, & 0, & $\dots$ & \\
$\dots$ & $=$  &  &  &  &  & \dots &  &  &  & 
\end{tabular}
\end{center}

В данном примере мы просто используем множества $U[0, 1]$ и $C[0, 1]$ из статьи ''Несчётность множества континуум, шах'' в качестве матрицы, определяющей какие предыдущие элементы входят в данный элемент, $1$ указывает на то, что элемент присутствует во множестве, $0$ - не присутствует. Ясно, что соответствие множествам $U[0, 1]$ и $C[0, 1]$ необходимо исключительно для простоты восприятия и правила построения множества $P$ могут быть заданы без подобных трафаретов. Так формально данную конструкцию в рамках наивной теории множеств мы можем определить следующим образом:
\begin{equation}
\emptyset \in P \wedge \forall x, y \; (x \in P \wedge y \in P \rightarrow y \cup \{ x \} \in P)
\end{equation}

Подобная конструкция является неудобной для современных аксиоматик, поскольку обнажает их архитектурные несовершенства. Если мы попытаемся определить класс объектов удовлетворяющей ей в NBG, то обнаружим, что данный класс совпадает со всем универсумом. Обратим внимание на то, что мы умышленно используем рекурсивный приём, аналогичный аксиоме бесконечности, чтобы проиллюстрировать некоторые тонкости современной интерпретации аксиоматик. В самом деле, давайте индуктивно развернём множество $P$: $P_0 = \emptyset$, $P_1 = \{ \emptyset \}$, $P_2 = \{ \{ \emptyset \} \}$,  $P_{n + 1} = P_n \cup \{ y \cup \{ x \} \mid x, y \in P_n \}$ и так далее. Не сложно видеть, что эта индуктивная последовательность соотносится с кумулятивной иерархией фон Неймана $V$ следующим образом: $P_n = V_{n}$. При этом мощность слоёв растёт одинаково $|P_{n + 1}| = 2^{|P_n|}$ и $|V_{n + 1}| = 2^{|V_n|}$. В рамках стандартной интерпретации трансфинитная индукция позволяет получить предельный слой: $P_{\omega} = ( \bigcup\limits_{i < \omega} P_i) \cup \{ y \cup \{x\} \mid x, y \in (\bigcup\limits_{i < \omega} P_i) \}$, где $\bigcup\limits_{i < \omega} P_i$ есть ничто иное как $V_{\omega}$ - множество всех наследственно конечных множеств. Но добавление элемента из $V_{\omega}$ в некоторое наследственно конечное множество оставляет его наследственно конечным, в силу чего $P_{\omega} = V_{\omega}$. Более того, $P_{\omega}$ оказывается неподвижной точкой для нашей формулы: $P_{\omega + 1} = P_{\omega} \cup \{ y \cup \{x\} \mid x, y \in P_{\omega}\} = P_{\omega}$. Иными словами, формула которая представляла весь универсум нашей теории замыкается на счётном уровне и не позволяет продвигаться дальше по трансфинитной лестнице. Но эта же формула включает в себя конструктивный принцип аксиомы бесконечности, задающий объект, допускающий трансфинитное продолжение. Данный факт демонстрирует, что внутри стандартной теории отсутствует сквозная индукция, а предельные элементы, связанные с разрывами непрерывности свойств, являются {\it точками разрыва предикативного описания} (точками прерывности индукции). Сквозная индукция была бы здесь крайне продуктивна, так как обеспечила бы продолжение построения $P_{\alpha}$, которое, в отличие от кумулятивной иерархии фон Неймана, уважительно относилось бы к порядковой структуре универсума. Проблема классической иерархии $V$ в том, что по достижении $V_{\omega}$ она начинает совершать экспоненциальные скачки мощности, размеры которых (в силу независимости континуум гипотезы) мы не в состоянии определить внутри теории. Индуктивное продолжение $P_{\alpha}$ могло бы стать гладкой иерархией, схожей с конструктивной иерархией Гёделя $L$, но доступной непосредственно внутри теории и не требующей метатеоретического аппарата выразимости. Однако предикативная прерывности предельных точек означает, что стандартное поведение математических операций не переносится на них из индуктивной области автоматически. Любое такое поведение, не выводимое строго из аксиом, представляет собой лишь соглашение и интерпретацию теории, допускающую альтернативные способы расширения операций на предельные объекты. Данное наблюдение имеет ключевое значение для понимания метаматематических развилок, возникших на этапе становления теории множеств. Забегая вперёд, обратим внимание на связь описанных конструкций с объектом вида $2^{\omega}$. Здесь и далее мы строго различаем {\it ординалы} - множества возникающие в аксиоматических теориях и {\it порядковые типы} с соответствующими {\it порядковыми типами}, возникающие как часть алгебраической структуры связанной с операцией конкатенации упорядоченных множеств. Ординал $2^{\omega}$ представляет собой множество всех функций $f : \omega \to \{ 0, 1 \}$ с конечным носителем (то есть $\{ i \in \omega \mid f(i) = 1 \}$ конечно, иногда также говорят о конечной поддержке), при этом множество носителей всех функций соответствующих ординалу $2^{\omega}$ является подмножеством $V_{\omega}$. В оригинальной конструкции Кантора вопрос об интерпретации счётного порядкового типа $2^{\omega}$, по мнению авторов, намеренно оставался открытым - в частности, он мог сопоставляться со всем множеством функций $f : \omega \to \{0, 1\}$, что связано с естественным гомоморфизмом между аддитивным и мультипликативным монойдами. Лишь в рамках аксиоматических теорий множеств зафиксировалось конкретное доопределение операций для предельных объектов, обеспечивающее тождество $2^{\omega} = \omega$ и жестко разделяющее кардинальное и ординальное возведение в степень - разделение, которое, как мы намерены показать, вовсе не является единственно возможным логическим следствием из базовых представлений о порядке. 

Давайте немного отвлечёмся от стандартной интерпретации, и рассмотрим данное множество вне аксиоматик, а как некоторый рекурсивный объект. Любопытно, каким именно образом парадокс Рассела воспроизводится на данном множестве: множества не содержащие себя - это множества соответствующие строкам имеющим $0$ в диагональной позиции, то есть абсолютно все множества, которые только могут быть определены данной конструкцией, но множество таких множеств по построению также является элементом $P$ (должно соответствовать строке $0.1(1)$), иными словами, определение множества приводит к тому, что оно должно быть своим собственным элементом, но в то же самое время оно не может быть своим собственным элементом, так как оно отличается от всех множеств принадлежащих ему хотя бы одним элементом. Данное наблюдение весьма любопытно, в свете того, что в теориях, где теорема Кантора недоказуема действительно возникают множества эквивалентные множеству своих подмножеств. При анализе данного множества проблемы начинаются уже с начальными бесконечными сегментами, то есть уже предельный переход к первому трансфинитному порядковому числу связан с каким-то затруднением, заставляющем наше сознание колебаться. Отметим здесь, что такое колебание имеет непосредственное отношение к противоречию возникающему в антиномии Рассела и понятиям потенциальной и актуальной бесконечности, но в рамках существующих теоретических конструкций мы не можем к нему подступиться. Несложно видеть, что данное множество противоречит аксиоме регулярности и теореме Кантора, и потому, корректней называть его классом или даже совокупностью, в том случае если мы принимаем аксиому регулярности на уровне классов, как в NBG. Но дело в том, что такого рода аксиоматические системы, просто ограничивает, то что мы можем считать множеством или классом, а что нет, то есть является своего рода {\it цензурой} фильтрующей объекты в нашем описании. Аксиома регулярности не влечёт математическую невозможность таких объектов, а лишь выводит их за рамки теории. Интересно, что присмотревшись к данному построению можно обнаружить связь с сюрреальными числами, в самом деле, на каждой итерации мы получаем набор множеств, характеризующих отношения объектов в предыдущем наборе, то есть даже такие замечательные структуры были воспроизводимы в терминах теории множеств, но мы поспешили исключить их из описания теории, которая легла в основу всей математики. Данная конструкция не представляет собой какой-то вырожденный случай, как множество всех множеств, подобных множеств можно построить бесконечно много. 

Такого рода наблюдения выявляют тесную связь аксиомы регулярности и противоречий, возникающих в контексте диагональных аргументов. Вообще говоря, подобные множества признаются невозможными и без аксиомы регулярности, что связывается с теоремой Кантора, выводящей данный факт из схемы аксиом выделения, принятой абсолютно для всех формул. Если мы обратимся к различным попыткам показать несостоятельность утверждения Кантора \cite{count-random}, то во многих случаях они сводятся к необходимости сопоставить некоторому вещественному числу натуральное число, имеющее бесконечную позиционную запись. О подобных попытках, в частности, пишет Уилфрид Ходжес \cite{hodges}, указывая, что термин натуральные числа имеет общепринятый смысл, принципиально отличающийся от предлагаемого в таких работах. Ошибочность такого рода построений проще всего продемонстрировать именно через аксиому регулярности: конструируя последовательность порядкового типа $\omega$, авторы фактически вводят элементы, индекс которых оказывается не меньше этого трансфинитного числа, что означало бы $\omega \in \omega$, напрямую противоречащее аксиоме регулярности. Конечно, в классической теории считается, что невозможность $\omega \in \omega$ вытекает уже из самого определения ординалов как фундированных транзитивных множеств и не требует отдельных аксиом. Однако авторы связывают подобные доказательства с определенными методологическими установками относительно допустимого использования такого инструмента как импредикативность. При ближайшем рассмотрении эти установки не находят удовлетворительного обоснования, а реальная сфера применения указанного инструмента оказывается значительно шире. Импредикативность является абсолютно необходимым понятием в математике, без которого невозможно мыслить континуум - среду отношений между объектами, где зафиксированное отношение само становится объектом, способным порождать новые отношения. Кантор определил порядковые типы, которые по своей природе сохраняют предикативность для начальных сегментов, но в предельной области проявляют свойства импредикативности, необходимые для однозначного разделения конечного и бесконечного. Поэтому, строя аксиоматику таких объектов, мы уже принимаем импредикативность как неотъемлемую часть нашего универсума. Для доказательства $\omega \notin \omega$ нам необходима минимальность $\omega$ как объекта, которую пытаются получить из формулы, выражающей такое стремление: $\omega = \{ x \in I \mid \forall Y \; (Ind(Y) \to x \in Y) \}$, где $I$ - некоторое бесконечное множество, существующее вследствие аксиомы бесконечности, а $Ind(Y)$ - предикат, фиксирующий индуктивную замкнутость множества $Y$ (аналогично аксиоме бесконечности). Однако сама по себе синтаксическая фиксация выражения не гарантирует, что в теории реализуется требуемая семантика. Например, формула $z = \{ x \in \omega \mid \exists y \; (x \in y \wedge y \in \omega) \}$ семантически описывает максимальный внутренний элемент $\omega$, но фактически представляет $\omega$. Аналогичным образом, пересечение индуктивно замкнутых множеств может не реализовывать искомую минимальность вследствие способности универсума порождать новые внутренние объекты. Минимальность может оказаться недостижимой просто потому, что для любого бесконечного объекта удаётся сконструировать новый объект, строго меньший данного. Здесь мы не можем сослаться на статичность универсума в платоновском смысле: импредикативность существует в мире платоновских идей точно так же, как и натуральные числа, являясь необходимой основой континуума. Более того, мы уже привносим импредикативность в наш универсум в направлении роста (при движении вверх по трансфинитной лестнице). Поэтому нам требуется отдельное утверждение, фиксирующее нашу веру в то, что универсум предикативен хотя бы снизу - в отношении минимальных элементов. Именно эту функцию и выполняет аксиома регулярности, обеспечивая предикативную фундированность отношения принадлежности $\in$. Давайте теперь зафиксируем наше наблюдение об отсутствии сквозной трансфинитной индукции, которое является следствием импредикативности связанной с бесконечностью, и не позволяет произвольно распространять свойства на предельные ординалы:

Теорема 4: {\sl предельные ординалы представляют собой точки разрыва предикативного описания.}

Доказательство: 

Мы воспользуемся тем фактом, что $\bigcup\limits_{ x \in \omega } x = \omega$. Обозначим множество $\bigcup\limits_{ x \in \omega } x$ как $X$, тогда указанный факт следует из $X \subseteq \omega$ и $\omega \subseteq X$ по аксиоме экстенсиональности. Чтобы доказать $X \subseteq \omega$ обратим внимание на то, что для всякого $y \in X$ по аксиоме объединения существует $x \in \omega$, такое что $y \in x$. Рассмотрев подмножество $T = \{ n \in \omega \mid n \subseteq \omega \}$ множества $\omega$, легко получить, что $T = \omega$, в силу минимальности $\omega$ и транзитивности $T$, которая непосредственно проверяема: $0 \in T$ из $0 \subset \omega$ и $i \in T$ влечёт $i \cup \{ i \} \in T$, так как $y = i \in \omega$ (в силу $T \subseteq \omega$) и $i \in T \to \forall j \in i \; (j \in \omega)$. То есть, все элементы $\omega$ являются его подмножествами, а значит из $y \in x$ следует $y \in \omega$ и требуемое $X \subseteq \omega$. Наоборот, для всякого $y \in \omega$ по аксиоме бесконечности имеем $y \cup \{ y \} \in \omega$, тогда замечая, что данный факт в точности отражает определение объединения заключаем $y \in X$ и $\omega \subseteq X$. Однако факт того, что $\bigcup\limits_{ x \in \omega } x = \omega$ вступает в противоречие с индуктивной цепочкой $\bigcup\limits_{ x \in \{ \emptyset \} } x \in \{ \emptyset \}$ и $\bigcup\limits_{ x \in n } x \in n \to \bigcup\limits_{ x \in (n \cup \{ n \}) } x \in (n \cup \{ n \})$, в силу того что $n \cup \{ n \}$ состоит из всех элементов $n$ и самого множества $n$. Таким образом, некоторые свойства не могут быть непрерывно перенесены на предельные ординалы, а значит такие ординалы представляют собой точки прерывности индукции. Теорема доказана.

Аксиома регулярности признаётся большинством исследователей и является весьма показательной: она отражает стремление математиков начала XX века в корне избавиться от парадокса Рассела и от объектов, противоречащих канонической интерпретации. В ней заключена попытка выразить интуицию универсальности индуктивной структуры в домене теории множеств, как одного из выделенных нами свойств необходимых для доказательства теоремы Кантора. Вывод о необходимости универсальности индуктивной структуры для доказательства теоремы Кантора, справедлив и в случае аксиоматических теорий, в чём мы убедимся в следующей статье настоящей серии, поэтому ссылка на теорему Кантора без поддержки универсальности, например, через соответствующую интерпретацию аксиомы регулярности, не обладает достаточной достоверностью. Однако, как мы выясним, и аксиома регулярности сама по себе имеет куда более ограниченный эффект, чем принято считать в рамках стандартной интерпретации. Она вовсе не блокирует альтернативные морфологии (структурные типы) бесконечности внутри теории, а лишь изолирует области функциональных отношений с фиксированной структурой - в том числе со структурой вполне упорядочения. В то же время аксиома регулярности вносит ряд существенных ограничений, сужающих область описания теории. Она жестко закрепляет наши установки относительно допустимого использования импредикативности, что затрудняет освоение этого инструмента, которым сам Кантор пользовался значительно более свободно. Рассмотрим наглядные примеры трансфинитной рекурсии, демонстрирующие разрыв между первоначальной интуицией создателя теории множеств и сформировавшимся впоследствии консенсусом. Применим трансфинитную рекурсию, чтобы переупорядочить область значений тождественной биекции $f : \mathbb{N} \to \mathbb{N}$, где $f(n) = n$. Положим $f_0(n) = n$ для всех $n \in \mathbb{N}$. Для каждого $j = i + 1$ определим перестановку $f_j$: $f_j(j) = f_i(2j), \quad f_j(2j) = f_i(j)$ и $f_j(n) = f_i(n)$ для всех $n \notin \{ j, 2j \}$. На предельном шаге $\omega$ зафиксируем поточечный предел: $f_{\omega}(x) = \lim\limits_{i < \omega} f_i(x)$. Легко видеть, что последовательность функций стабилизируется в каждой конечной точке, гарантируя существование данного поточечного предела. Обратим внимание: на каждом конечном шаге $j$ функция $f_j$ остается строгой биекцией на $\mathbb{N}$, а перестановки затрагивают только конечные позиции. Однако после $\omega$ шагов рекурсии мы получаем функцию $f_{\omega} : \omega \to \omega$, заданную формулой: $f_{\omega}(x) = 2x$. Ее область значений $Im(f_{\omega})$ содержит только чётные числа. Нечётные числа полностью выпадают из образа предельной функции, хотя на любом конечном шаге функции оставались биективными. Аналогичная ситуация возникает, если мы строим рекурсию, ориентируясь на область определения. Положим $f_0(n) = n$. На шаге $j = i + 1$ для четного индекса $j = 2k$ меняем местами значения $f_i(k)$ и $f_i(j)$: $f_j(k) = f_i(j)$, $f_j(j) = f_i(k)$ и $f_j(n) = f_i(n)$ при $n \notin \{ k, j \}$. Переходя к пределу $f_{\omega}(x) = \lim_{i < \omega} f_i(x)$, мы вновь получаем инъективное отображение, не являющееся сюръекцией. На каждом конечном этапе нечётное число успевает стать значением аргумента, не превосходящего номер шага; мы явно пробегаем по всем элементам области определения, но в пределе нечётные числа полностью исчезают из области значений. В стандартной интерпретации область определения функции не может измениться, поскольку $\mathbb{N} = \omega$ строго зафиксирована (а $2\omega$ - это уже совершенно иной ординал экстенсионально различный с $\omega$). Построение предельного объекта $f_\omega$ гарантируется правилами трансфинитной рекурсии, привязанными к ординалам. В результате возникающий феномен оставляет единственную возможность: предельный переход разрушает сюръективность, принудительно исключая элементы из области значений. Несложно задать и такую трансфинитную рекурсию, которая упорядочит нечётные числа по возрастанию так, что число $1$ окажется первым элементом после всех чётных чисел. В самом деле, можно итерироваться по чётным числам $2i$, перемещая их в половинную позицию ($2i \to i$), и параллельно сдвигать все промежуточные элементы, меньшие $2i$, в следующую позицию ($\forall j: i \le j < 2i \implies j \to j + 1$), подобно тому как пузырьки всплывают на поверхность. Правда, в таком случае число перестановок растёт квадратично, что требует дополнительных соглашений о предельном переходе для трансфинитной рекурсии, которых хотелось бы избежать на данном этапе.

Если мы обратимся к исходным работам Кантора, где $\omega$ понимается как порядковый тип (класс эквивалентности вполне упорядоченных множеств) и где Кантор явно указывает на допустимость перестановок \cite{cantor-relation}, то можно обнаружить совсем иные результаты. Действительно, область определения функции сама по себе не является порядковым типом - она представляет собой множество (в нашем случае $\mathbb{N}$), которое лишь аннотируется (описывается) порядковым типом $\omega$. Перестановка значений функции разделяет область определения на два сегмента (отображающийся в чётные числа и отображающийся в нечётные), тем самым изменяя порядковый тип, посредством которого мы описываем эту область. Иными словами, область определения как базовое множество остаётся неизменной - меняется лишь её структурная аннотация. Без дополнительной информации об области значений она описывается порядковым типом $\omega$, однако с использованием переупорядочивающей функции мы можем описать ту же самую область определения порядковым типом $\omega + \omega$ (или $2\omega$). Кантор явно акцентировал внимание на абстрагировании от природы элементов множеств, то есть от их внутренней идентичности. Поскольку Кантор работал с классами эквивалентности алгебраических структур, понятие порядка на множестве и взаимно однозначного отображения на некоторое множество определённого порядкового типа - это по сути одно и то же. Чтобы задать порядок, нам достаточно указать такое отображение; при этом, переупорядочивая область значений этого отображения, мы меняем порядок и на области определения, поскольку переупорядочиваем множество пар. Формально порядок задаётся как биекция $f : X \to O$, где $O$ - некоторая упорядоченная структура, представляющая порядковый тип и индуцирующая его на множестве $X$. Если сам носитель $X$ обладает естественным (интуитивным) порядком, то его можно связать с тождественным отображением $id : X \to X$, что уже сильно проясняет ситуацию: мы способны идентифицировать не сами объекты, а их позиции в определённой структуре. Таким образом, в случае алгебраических структур множество может быть изоморфно (подобно) своему собственному подмножеству, в том числе и собственному начальному сегменту. Однако невозможно аннотировать одним и тем же порядковым типом начальный сегмент и всё множество одновременно. В алгебраической структуре порядковых типов Кантор фактически занимается вопросом непротиворечивой аннотации множеств. Когда же в аксиоматической теории множеств возникает бесконечный ординал $\omega$, ассоциируемый с соответствующим порядковым типом Кантора, некоммутативная арифметика порядковых чисел, которую мы пытаемся спроецировать на множество $\omega$, указывает на то, что это множество должно стать неподвижной точкой, поскольку $1 + \omega = \omega$. У нас нет изящного способа перенести это свойство непосредственно на отношение принадлежности $\in$, так как в разных контекстах выражение $\omega \in \omega$ имело бы совершенно разный смысл. Всё дело в том, что Кантор требовал предикативности исключительно от структуры, абстрагируясь от конкретного содержания элементов множества. В то время как стандартный формализм пытается одновременно потребовать и предикативности структуры, и жесткой идентифицируемости элементов по включению (экстенсиональности).

Мы могли бы обойти возникающие противоречия, отказавшись от концепции {\it рефлексивности} - возможности биективного соответствия бесконечных множеств со своей собственной частью. Тогда арифметика ординалов стала бы коммутативной, а порядковые типы возникли бы как классические классы эквивалентности. Однако мы выбираем иной путь: сохраняя и однозначность отношения принадлежности $\in$ (определяющего жесткую вполне упорядоченную структуру ординалов), и рефлексивность (импредикативность), для чего вводим аксиому регулярности. Это позволяет объявить выражение $\omega \in \omega$ бессмысленным или неопределённым (аналогично операции $\omega - 1$), что концептуально напоминает механизм стратификации формул в системе новых оснований Куайна \cite{quine-new-foundations}. На этом шаге и возникает понятийная путаница: мы вводим отношение порядка на множестве, но не даём строгого ответа на вопрос о том, чем является порядок внутри самой теории. Само по себе множество - это лишь набор элементов; отношение порядка - это множество упорядоченных пар, удовлетворяющих определенным условиям. Таким образом, упорядоченное множество - это всегда пара $(X, R)$, которую можно эквивалентно переписать через биекцию с эталоном $(X, f)$. Порядок - это отдельный математический объект, самостоятельное множество (в контексте ординалов это класс пар определяющий само отношение принадлежности $\in$). Следовательно, даже в случае ординалов нам остается доступной исходная интерпретация Кантора: предельная функция $f_{\omega}$ на множестве $\mathbb{N}$ индуцирует новый порядок $(\mathbb{N}, f_{\omega})$, изоморфный ординалу $\omega + \omega$, в то время как естественный порядок $(\mathbb{N}, \in)$ остаётся неизменным и изоморфным $\omega$. В таком рассмотрении мы отчетливо видим проявление той самой импредикативности, которую Кантор заложил в определение порядковых типов. Мы не смогли бы прямо выразить элементы $f^{-1}(2k + 1)$ в теории не потому, что утратили информацию о самих элементах, а потому что биекция $f$ не связана с канонически отношением $\in$ для ординалов, через которую мы идентифицируем эти элементы внутри упорядоченной структуры в рамках теории. Элементы сами по себе не теряют своей идентичности, но импредикативность взаимоотношения объектов, определяющих порядок препятствует использованию структуры $\in$ для их синтаксической выразимости, в силу того, что теория не содержит требования о возможности одновременной интерпретации двух разных упорядочивающих преобразований, аннотирующих одно и то же множество. Это явление служит прямым намеком на идеи нестандартного анализа, распространяющего предикативность существенно дальше классических границ. 

В частности, нам уже хорошо знаком данный эффект на примере всё того же упорядочивающего отношения $\in$ и морфологии континуума, импредикативность взаимосвязи данных объектов не позволяет явно выразить биекцию между структурами вполне упорядочения и замыкания всюду плотности. Кантор выстраивает гениальную конструкцию, демонстрирующую причины данного явления, которые заключаются в том, что вполне упорядочение связано с импредикативностью аддитивной структуры, в то время как замыкание всюду плотного множества требует предикативности минимум экспоненциального уровня. Иными словами, предикативность континуума связана с мультипликативными уровнями воспроизводства аддитивной структуры, элементы которой попадают в импредикативные участки аддитивной структуры вполне упорядочения. Так например объект $2^{\aleph_0}$ не имеет никакого следа в богатой порядковой структуре, перепрыгивая из конечной области сразу в область неопределённой мощности, что объясняется именно тем, что данный след просто не выразим в аддитивной структуре. Однако мы выбираем интерпретацию, насильственно отождествляющую порядок и внутреннюю принадлежность ординалов. Отсутствие объективной математической необходимости за этим выбором явно указывает на его идеологическую и методологическую природу. Не вполне ясно также, почему аксиоматики с таким упорством стремились избежать неподвижных точек, связанных с синтаксисом теории множеств. К примеру, они широко и плодотворно используются в математическом анализе, где неподвижные точки (пределы) представляют собой классы эквивалентности фундаментальных последовательностей. Даже их соотношение друг с другом может варьироваться (как в случае последовательностей $1/n$ и $1/3n$), а их рекурсивная интерпретация точно так же может приводить к циклам ($x_0 = 1, \; x_{n+1} = x_n / 2, \; 0 = 0 / 2$). Любопытно, что когда Джон фон Нейман предлагает ординалы в качестве модели порядковых типов, он прямо указывает: эта конструкция имеет строгий смысл только тогда, когда трансфинитная индукция уже определена, что позволяет выбрать в классе эквивалентности ординалов, соответствующем порядковому типу $\omega$, одного канонического представителя \cite{von-neumann-ordinals}. Данный факт как раз и должен был способствовать аккуратному формированию интерпретации при адаптации теории Кантора - если мы стремимся сохранить изоморфизм между структурами ординалов и порядковых типов, который фон Нейман доказывает именно в предположении первичности последних. Так или иначе, аксиома регулярности существенно затрудняет возврат к исходной интерпретации Кантора, в которой соотношение количественных характеристик множеств в рамках одного класса эквивалентности порядкового типа $\omega$ могло варьироваться.

Чтобы понять как всё это связано с диагональным построением Кантора, приведём ещё один пример с бесконечным множеством натуральных чисел. Рассмотрим подмножество $\mathcal{P}^*(\mathbb{N})$ множества $\mathcal{P}(\mathbb{N})$ состоящее из начальных сегментов натурального ряда (включая $\emptyset$ и исключая множество всех чисел $\mathbb{N}$). Легко видеть, что такое множество взаимно однозначно соответствует бесконечному множеству натуральных чисел $\mathbb{N}$, а в контексте ординальной алгебры это и есть само множество натуральных чисел: $\varphi(n) = \{ 0, 1, 2, \dots, n - 1 \}$. Запишем данную биекцию сохраняющую порядок в позиционной записи, где $0$ означает, что элемент, соответствующий данному столбцу, не принадлежит множеству, а $S$ - принадлежит:
\begin{center}
\begin{tabular}{c c c r r r r r c} 
 & & & 0 & 1 & 2 & 3 & $\dots$ & \\
$0$ & $\Leftrightarrow$  & $|$ & 0 & 0 & 0 & 0 & $\dots$ & $|$\\
$1$ & $\Leftrightarrow$  & $|$ & S & 0 & 0 & 0 & $\dots$ & $|$\\
$2$ & $\Leftrightarrow$  & $|$ & S & S & 0 & 0 & $\dots$ & $|$\\
$3$ & $\Leftrightarrow$  & $|$ & S & S & S & 0 & $\dots$ & $|$\\
$\dots$ & $\Leftrightarrow$  & $|$ & S & S & S & S & $\dots$ & $|$
\label{constr:n-bijection}
\end{tabular}
\end{center}

Такую матрицу можно также рассматривать как унарное кодирование, где $S$ - отношение следования. Мы используем подмножества натуральных чисел из множества степени, чтобы иметь универсальную конструкцию, которая не зависит от их представления в теории, что позволяет нам сравнить интерпретацию в система Кантора и современный вариант, а кроме того демонстрирует некоторую связь между множеством степени и кодированием в метатеории, осознание которой, к сожалению, было утеряно при аксиоматизации. В современной интерпретации мы имеем диагональный элемент $Z = \{ x \in \mathbb{N} \mid x \notin \varphi(x) \}$, который равен множеству $\mathbb{N}$. В то время как в исходной конструкции порядковых типов Кантора и в метатеории происходит настоящая магия: диагональный элемент представляет из себя объект отличный от всех натуральных чисел благодаря изменению $0$ на $S$, но такое отличие в контексте нашей матрицы - это фактически операция добавления единицы $+1$ или перехода к следующему элементу $S(x)$. Иными словами, диагональное построение позволяет для такого импредикативного объекта как порядковый тип - класс эквивалентности подобных множеств, предикативно сконструировать элемент непосредственно следующий за всеми элементами данной аморфной совокупности, для чего требуется лишь операция локального отрицания тождества элементов через повышение ранга одного из них. Не сложно видеть, что данный трюк работает, в том числе, для любого конечного множества, что даёт нам строгое обоснование правому (внешнему) сложению, которое как бесконечная сумма ряда становится поддержкой интуиции пределов Кантора: $\omega = \sum\limits_{i \in \mathbb{N}} 1$. Если теперь имея порядковый тип $\omega$ добавить в конец полученный объект $\omega$ мы получаем новый порядковый тип $\omega + 1$, последний элемент которого предикативен и равен $\omega$. Таким образом, всплытие элемента $\varphi(e_{n}) = e_{n + 1}$ (здесь критична именно правая операция), которое в кардинальной алгебре обеспечивало эквивалентности $\aleph_0 + 1 = \aleph_0$, в порядковой алгебре приводит к возникновению элемента строго отличного от всех элементов последовательности. Невозможно в очередной раз не обратить внимание насколько в этом отношении точны формулировки Кантора, он не использует данный кардинальный метод в порядковой алгебре, вместо этого он демонстрирует подобие $E = (e_1, e_2, \dots, e_n, \dots)$ и $(f, E) = (f, e_1, e_2, \dots, e_n, \dots)$ чисто синтаксически, то есть через обращение к метатеории, что можно транслировать в формальную конструкцию $\varphi(e_1) = f, \varphi(e_n) = e_{n - 1}$ \cite{cantor-inner-addition} (которая в силу тривиальных соображений является биекцией). Даже в современном контексте это имеет смысл, поскольку функция представляет собой множество упорядоченных пар, то есть имеет выделенное направление, и вообще говоря $\varphi(e_1) = f, \varphi(e_n) = e_{n - 1}$ и $\phi(f) = e_1, \phi(e_n) = e_{n + 1}$ не являются эквивалентными друг другу объектами. Описанная диагональная конструкция позволяет также видеть, что преобразование всех элементов порядкового типа $\omega$ правым (внешним) сложением даёт некоторый объект, где все строки смещены отношением следования, при этом диагональный элемент становится неотъемлемой частью данного объекта (в матрице \ref{constr:n-bijection} диагональ заполняется символом $S$). В теории порядковых типов такой объект имеет начальный сегмент подобный $\omega$ и последний элемент отличный от всех, фактически $(\cdot + 1) : \omega \to \omega + 1$. При этом левое (внутреннее) сложение остаётся рефлексивным и замкнутым внутри множества $\omega$: $(1 + \cdot) : \omega \to \omega$. Таким образом, операции левого (внутреннего) и правого (внешнего) сложения в конструкции Кантора имеют онтологическую поддержку, позволяющую нам однозначно разрешить конфликт предикативности и импредикативности, в том числе используя диагональный метод. Но давайте обратимся к современной интерпретации теории множеств. У нас есть единственная операция следования $x \cup \{ x \}$, которую мы интерпретируем исключительно во внутреннем смысле $x \cup \{ x \} = (1 + x)$, мы выводим её рефлексивность из аксиомы бесконечности и схемы аксиом преобразования. Внешние операции появляются благодаря тем же аксиомам теории, но выражаются идентичной формулой $(x + 1) = x \cup \{ x \}$. Для того, чтобы описать разницу между левым и правым сложением, мы вынуждены возвращаться к концепции порядковых типов, поскольку наша теория не способна их различить. Вследствие этого мы вынуждены отвергнуть диагональную конструкцию \ref{constr:n-bijection}, поскольку у нас в теории просто нет функции, которая способна соответствовать правому сложению Кантора $(\cdot + 1)$ и выскочить за пределы индуктивного замыкания, такое отрицание оказывается неконтролируемым и распространяется на фрагменты интерпретации теоремы Кантора, чего мы даже не замечаем. Даже несмотря на то, что наш собственный синтаксис сохраняет рудименты внешней семантики $\bigcup\limits_{ x \in \omega } x = \omega$, гарантирующие существование внешней интерпретации отношения следования, содержащей $\omega$ в своей области значений. Иными словами, мы полностью утрачиваем онтологическую поддержку наших операций, вместе с инструментами разрешения конфликта предикативности и импредикативности, что судя по всему и становится одной из основных причин потери изоморфизма между построениями Кантора и современной интерпретацией. Как ни парадоксально, сам по себе используемый нами формализм допускает возвращение к онтологической глубине исходных построений, то есть наш синтаксический аппарат оказывается адекватнее собственной семантической интерпретации. Дело в том, что современные теории множеств не гарантируют универсальную корректность теорем во всех моделях, что представляет собой особую форму некатегоричности, которая является одним из существенных архитектурных недостатков этих теорий, в настоящее время уже в значительной мере осознанный.

Вернёмся к нашему примеру с трансфинитной рекурсией и изменим биекцию $\psi = \mathbb{N} \to \mathbb{N}$, так чтобы пустому множеству не соответствовало ни одно конечное натуральное число. Задаём $\psi_0 = \psi$ и для $j = i + 1$ определяем $\psi_j(i) = \psi_i(j)$, $\psi_j(j) = \psi_i(i)$ и $\psi_j(x) = \psi_i(x)$ для $x \notin \{ i, j \}$, $\psi_{\omega} = \lim\limits_{i < \omega} \psi_i(x)$. Мы снова видим как пустое множество в результате исключается из области значений функции $\psi_{\omega}$. В то же время, если мы возьмём готовое переупорядоченное множество $\mathbb{N}^{*}$ изоморфное $\omega + 1$, которое начинается с $1$ и содержит $0$ на позиции $\omega$, то изоморфизм $\varphi : \mathbb{N}^{*} \to \omega + 1$ будет иметь вид: $\varphi(0) = \omega, \varphi(y) = x - 1$. Указывая на то, что в множестве $\mathbb{N}$ достаточно элементов, сопоставляя которые с метотеоритическими отношениями следования мы способны выйти за пределы множества $\mathbb{N}$, получив элемент $\omega$ отличный от всех конечных. У нас появляется абсолютно законная конструкция, которая восстанавливает правую операцию сложения Кантора. Для этого требуется выход в метатеорию, но это лишь говорит о том, что метатеория обладает каким-то свойством позволяющим различать замкнутое левое сложение от сквозного правого. Но вместо того, чтобы задаться вопросом, что это за свойство, и может ли оно быть определено внутри теории, мы продолжаем упорствовать в том, что множество $\omega$ предикативно задано через экстенсиональность и не допускает импредикативных элементов. Несмотря на то, что Джон фон Нейман явно формулирует, что ординальное представление имеет смысл для уже заданной порядковой арифметики, то есть де-факто является метатеоритическим обозначением порядковых типов. Таким образом, когда мы трансфинитной рекурсией позволяем элементу $0$ всплыть в области значений, мы можем рассматривать дилемму с областью определения как необходимость изменения её семантической аннотации в метатеории. То есть, нам вовсе не обязательно предполагать изменение области значений, мы вполне законно можем описывать явление, как рассинхронизацию синтаксиса и семантики вызванную импредикативностью порядковых типов. Такого рода сложности весьма показательны и рассредоточены по всей современной теории множеств, они указывают нам на то, что нашего инструментария недостаточно для описания и раскрытия конфликта импредикативности и предикативности, в том числе той роли которую в нём играет метатеория. Об этом факте почему-то умалчивает Ходжес, высмеивая состояние ума авторов критикущих аргумент, предлагая другой аргумент, который не приводит к тому же выводу, в остальном дающий, на взгляд авторов, практически безупречный анализ возможных уязвимостей теоремы Кантора. Впрочем, когда профессиональные математики комментируют работы непрофессионалов, стоит акцентировать внимание, на том, что в математике не существует запретов (''нельзя'' не имеет никого другого смысла кроме идеологического), и во многих случаях размышления последних оказываются математически корректными, но выходящими за пределы общепринятого консенсуса, например, построенные объекты могут не являться множествами в рамках конкретной аксиоматики, хотя вполне могут существовать. Состояние ума о котором говорит Ходжес называется рациональностью, поскольку наличие двух противоположных друг другу аргументов в теории может либо говорить о её противоречивости, либо свидетельствовать об отсутствии общезначности её утверждений в моделях. Если мы не разделяем веру в то, что пустое множество исчезает из области значений функции $\psi$, то наблюдаем идентификацию импредикативного элемента внутри $\mathbb{N}$ отличного от всех предыдущих, что семантически означает его соответствие $\omega$ и изменение ординальной аннотации области определения.

Данный пример чрезвычайно показателен, поскольку открывает сразу несколько интерпретационных тонкостей современных аксиоматических теорий множеств. Во-первых, по недостаточно ясным причинам классы эквивалентности бесконечных множеств, которые исследовали Кантор и Дедекинд, подчёркивая подобие объектов в контексте алгебраических структур, в рамках аксиоматических теорий превратились в фиксированные множества, которые приобрели свойство рефлексивности. Важно различать свойство бесконечности по Дедекинду \cite{dedekind-infinite-set} и свойство рефлексивности. Первое определяет классы эквивалентности в рамках алгебраических структур систем элементов, исследованных Дедекиндом, и означает лишь, что бесконечные системы подобны своим частям (то есть часть бесконечного может описываться так же, как целое). Второе возникает, по всей видимости, в результате случайного упрощения, сливающего две независимые арифметики - кардинальную и ординальную - в единую алгебраическую систему, в которой приобретает смысл отображение объекта на свою часть. Иными словами, эквивалентность и подобие, о которых говорили Кантор и Дедекинд, реализуются посредством отображения структуры, аннотирующей целое, на структуру, аннотирующую часть, тогда как рефлексивность связана с буквальным отображением целого объекта на свою часть. Сама по себе рефлексивность не приводит к проблемам, поскольку допустимо считать, что объектами теории являются исключительно элементы структур, не имеющие иной онтологической природы. Однако ключевую роль играет место рефлексивности в существующей интерпретации теории множеств, где наблюдается некритическое расширение домена её семантического значения, выходящее далеко за рамки реально доказуемого смысла этого понятия. Именно для прояснения этого смысла мы и планируем опираться на исходные интуиции Кантора и Дедекинда. Во-вторых, элементы множества при переупорядочении способны структурироваться в различные порядковые типы, что связано с некоммутативностью операции конкатенации (различием левого и правого сложения). В отсутствие какого-либо порядка операция объединения множеств коммутативна, вследствие чего описание с помощью единой биекции без учёта алгебраической структуры приводит к неоднозначности. Совершенно не случайно Кантор подчёркивал возможность любых перестановок в любом количестве при определении неупорядоченного взаимно однозначного соответствия, тогда как в случае взаимно однозначного соответствия с сохранением порядка перестановки могут изменить порядковый тип. Данный факт свидетельствует о том, что биекция в кардинальной алгебре и биекция в порядковой алгебре могут представлять собой совершенно разные объекты, выразимость которых друг через друга требует специального доказательства. В частности, равномощность в смысле биекции в произвольном универсуме (случай Кантора) и равномощность в смысле биекции исключительно в упорядоченном универсуме (современная интерпретация) - это два принципиально разных понятия, не образующих изоморфных друг другу структур. В-третьих, замена импредикативных классов эквивалентности фиксированными экстенсионально предикативными множествами затрудняет понимание наблюдаемых в теории явлений. Феномен, при котором элементы аннотации структуры становятся полноценными объектами (карта воспринимается как территория), принято называть {\it реификацией}. В большинстве случаев контролируемая реификация - например, та, которую осуществляет фон Нейман, сопоставляя порядковым типам ординалы как канонические представители, - наоборот, делает описание более простым и наглядным. Однако реификация всегда происходит в метатеории, которая сама оказывается связанной с конфликтами импредикативности и предикативности, которые мы и стремимся прояснить. Это можно наблюдать, в частности, в относительности понятия истины, необходимой в фундаменте, но допускающей лишь интуитивное определение на философском уровне вне формальных систем. Таким образом, современная интерпретация теории множеств обусловлена синергетическим эффектом множества нюансов, каждый из которых сам по себе не играет решающей роли. Среди этих нюансов: архитектурные недостатки, связанные с требованиями к аксиоматическим системам и затрудняющие корректную атрибуцию понятий теории; некритическое смешение алгебраических структур, нарушающее изоморфизм аксиоматизируемым построениям; и реификация, вплетающая знание о метатеории в данную историю. Математические формулы обретают смысл в контексте тех построений, в которых они используются: объекты неотделимы от структуры, а структуры - от объектов. В современных же теориях множеств фокус смещён на формулы, которыми пытаются описывать две (а иногда и три) различных алгебраических структуры. Сама по себе реификация - замена класса эквивалентности на объект - не столь опасна: любое определение и любая аксиома представляют собой ограничение области описания, которому удовлетворяет некоторая совокупность объектов; не имеет принципиального значения, считать ли эту совокупность единой сущностью или коллекцией сущностей. Однако когда в при такой реификации одной формулой начинают описываться объекты различной природы могут возникать существенные сложности. В свою очередь, при использовании формул для аксиоматизации необходимо чётко понимать, какой смысл они обретают в контексте определяемых теорией объектов, и строго ограничивать семантику данного смысла.

Пожалуй, самым безобидным феноменом современных аксиоматических теорий является существование объектов, невыразимых на языке теории, - объектов, которые мы не можем определить или идентифицировать. Уже само по себе это свидетельствует о некотором несоответствии синтаксиса и семантики. В языке формальной теории, состоящем из счётного набора символов, можно построить лишь счётное количество формул, тогда как объекты, существующие в модели (например, во вселенной фон Неймана $V$), могут иметь несчётную мощность. В результате многие объекты не имеют соответствующих формул. Формально множество $x$ называется {\it неопределимым}, если не существует формулы $\phi(x)$ такой, что $\exists! x : \phi(x)$. Например, невозможно идентифицировать большую часть вещественных чисел и подмножеств натуральных чисел. Данный феномен активно используется исследователями: идеи Коэна легли в основу метода форсинга - расширения моделей теории путём добавления новых обобщённых объектов, не определяемых формулами теории в исходной модели \cite{forcing}. Однако в действительности этот феномен связан вовсе не с мощностью множеств, а с импредикативностью объектов, необходимой для рефлексивности бесконечных множеств. Тот факт, что чаще всего об отсутствии выразимости говорят именно в контексте вещественных чисел, указывает на то, что для описания континуума требуется импредикативность, которая и обеспечивает его замкнутость и подобие целого части. В качестве иллюстрации можно привести и более экзотические примеры. Если отказаться от аксиомы регулярности, становятся возможными бесконечные убывающие цепочки вложенности и множества вида $x = \{x\}$. Это позволяет определить класс ординальных степеней синглтонов $Sin$, удовлетворяющий следующим условиям:
\begin{itemize}
    \item $\emptyset \in Sin$;
    \item если $s \in Sin$, то $\{ s \}^1 = \{ s \} \in Sin$;
    \item отношение порядка $<$ на $Sin$ вводится рекурсивно и удовлетворяет условиям: $\forall s \in Sin \; (s < \{ s \}^1)$, $n < m \to \{ s \}^n < \{ s \}^m$;
    \item для предельного $\lambda$ и возрастающей последовательности $(\{ s \}^{\alpha})_{\alpha < \lambda}$ элементов $Sin$, полагаем $\{ s \}^{\lambda} = \sup \{ \{ s \}^{\alpha} \mid \alpha < \lambda \}$, являющийся пределом этой последовательности, причём $\{ s \}_{\lambda} \in Sin$.
\end{itemize}

Сказанное следует из приводимой ниже теоремы, в доказательстве которой используется подход, аналогичный введению степеней порядковых чисел Кантором \cite{cantor-power}. 

Теорема 5: {\sl Существует единственная класс-функция $f : Sin \times Ord \to Sin$ удовлетворяющая следующим условиям: 
\begin{enumerate}
    \item $f(x, 0) = x$;
    \item $\forall \alpha, \beta$ если $\alpha < \beta$, то $f(x, \alpha) < f(x, \beta)$;
    \item $\forall \alpha \; f(x, \alpha + 1) = \{ f(x, \alpha) \}$.
    \item если $\lambda = \lim\limits_{i \in I} \alpha_i$, то $f(x, \lambda) = \lim\limits_{i \in I} f(x, \alpha_i)$.
\end{enumerate}}

Доказательство: 

Поскольку порядок на $Sin$ изоморфен порядку ординалов, предельные синглтоны являются топологическими пределами: в любой открытой окрестности такого синглтона содержатся все синглтоны, начиная с некоторого номера. Тогда по теореме о трансфинитной рекурсии для каждого $x \in Sin$ существует единственная функция $F_x$, удовлетворяющая рекурсивным соотношениям 1, 2 и 4. Определяем $f(x, \alpha) = F_x(\alpha)$. Свойство 2 доказывается трансфинитной индукцией. Покажем единственность. Пусть $g$ — любая другая функция с теми же свойствами. Тогда множество $\{\alpha \mid f(x, \alpha) = g(x, \alpha)\}$ содержит 0 и индуктивно замкнуто, следовательно совпадает с классом ординалов. Теорема доказана.

Таким образом, у нас возникают весьма интересные объекты: $\{ s \}^0 = s$, $\{ s \}^1 = \{ s \}$, $\{ s \}^2 = \{\{ s \}\}$, $\{ s \}^3 = \{\{\{ s \}\}\}$, $\dots$, $\{ s \}^{\omega} = \underbrace{\{\dots\{\{}_{\omega} s \underbrace{\}\}\dots\}}_{\omega}$, $\dots$. Они помогают лучше понять природу ограничений теории на невыразимость некоторых объектов. Интересно, что Абрахам Френкель показал эквивалентность подходов, основанных на конечных синглтонах и ординалах; в частности, аксиому бесконечности можно формулировать как для множества конечных синглтонов, так и для ординалов \cite{fraenkel-singleton-ordinal-equivalence}. Арифметика степеней синглтонов полностью определяется арифметикой ординалов, а сами эти объекты не противоречат никаким аксиомам, кроме, возможно, аксиомы регулярности. Точнее, мы полагаем, что данные объекты противоречат аксиоме регулярности потому, что естественными кандидатами на роль предельных синглтонов видятся либо неподвижные точки $\{ s \}^{\omega} = \{ \{ s \}^{\omega} \}$, либо объекты допускающие бесконечные убывающие цепочки $\dots \{ s_2 \}^{\omega} \in \{ s_1 \}^{\omega} \in \{ s_0 \}^{\omega} = \{ s \}^{\omega}$. В частности, к равенству $\{ s \}^{\omega} = \{\{ s \}^{\omega}\}$ можно прийти, если использовать порядковые типы Кантора для описания синглтонов, где доступна формула разности $\omega - 1 = \omega$ \cite{cantor-limit-minus}. В случае ординалов такая формула нуждается в дополнительных соглашениях, поскольку в силу аксиомы экстенсиональности множество без одного элемента не может быть равно самому себе. Таким образом, оригинальные построения Кантора вполне допускали конструкции вида $x = \{ x \}$. В сущности, множество $\{ s \}^{\omega}$ мало чем отличается от множества $\omega$: точно так же можно убедиться в отсутствии непосредственного предшественника $\{ x \in \{ s \}^{\omega} \mid \{ x \} = \{ s \}^{\omega} \} = \{ x \in \omega \mid x \cup \{ x \} = \omega \} = \emptyset$. Единственное различие в том, что неопределимый объект в нём не соседствует с другими элементами. В то же время в рамках нестандартного анализа во внутренней модели синглтон $\{ s \}^{\omega}$ никак не конфликтует с аксиомой регулярности. Это указывает на возможность реализации некоторых альтернативных конструкций, позволяющих таким объектам существовать в рамках полноценных теорий ZFC и NBG. Более того, аксиома регулярности не конфликтует с аксиомой бесконечности, которая определяет топологически идентичную конструкцию с изоморфным уровнем вложенности множеств. Это также свидетельствует в пользу законности предельных синглтонов. Дело в том, что импредикативность представляет собой неполноту информации - синтаксическую неопределённость, которую мы инстинктивно пытаемся заполнить некоторым семантическим смыслом, например в форме неподвижных точек. Когда мы пишем предикат ''максимального натурального числа'', а теория выдаёт нам само множество натуральных чисел, это не следствие импредикативности как таковой, а следствие дополнительных соглашений о том, как именно заполнять информационную пустоту. Кантор вкладывает в структуру вполне упорядочения именно такой смысл: порядковые типы - это объекты, предикативные в области минимальных элементов и импредикативные в области максимальных. Мы же добавляем аксиому экстенсиональности и стремимся добиться полной предикативности внутреннего содержания этих объектов. Проблема, однако, заключается в том, что для этого используется импредикативная форма аксиомы бесконечности, которая не содержит информации о максимальных элементах. Иными словами, онтологически контролируемое взаимодействие импредикативности и предикативности, реализованное в теории Кантора, подменяется двумя взаимоисключающими намерениями. Порядковые типы используют импредикативность, чтобы перестраиваться, когда мы пытаемся идентифицировать объекты в области максимальных элементов; именно это обеспечивает полное отсутствие таких элементов. Экстенсиональность же ограничивает способность ординалов к перестройке и накладывает довольно жёсткие требования к теории, призванные гарантировать полное отсутствие альтернативной семантики, к которой могли бы редуцироваться синтаксические выражения. Когда речь идёт о предельных синглтонах, имеется в виду исключительно семантика решения предельного уравнения $\forall N \; \exists \{ s \}^{N} : \forall n > N \; (\{ s \}^{N} < \{ s \}^{n} < \{ s \}^{\omega})$. Мы ищем такое решение среди стандартных семантических конструкций, связанных с синтаксическими соглашениями (неподвижные точки и убывающие цепочки вложенности). Однако перебор стандартных вариантов не означает доказательства отсутствия соответствующих объектов в теории. Такое наблюдение позволяет говорить о некотором решении уравнения, семантический смысл которого пока неизвестен. $\{ s \}^{\omega}$ - это не аналитическое решение, а константа, обозначающая решение, которое должно существовать вследствие стандартной топологии. Мы по какой-то причине упорно пытаемся маскировать явное указание на использование такого инструмента как импредикативность. Между тем такое указание нисколько не разрушает теорию, но делает её более честной. Поскольку $\{ s \}^{\omega}$ является синглтоном, мы знаем, что он содержит некоторое множество, но формула этого множества нам недоступна: она может быть определена только через степень, которая сама является неопределимым ординалом. Иными словами, ограничения на идентифицируемость объектов связаны не только и не столько с мощностью, сколько с сознательным ограничением информации внутри теории для реализации импредикативности, необходимой для замыкания алгебры. В сущности, именно на этом построен метод форсинга: если для теории $\{s\}^{\omega}$ является внешним обобщённым объектом, то теория внутри себя не может вывести противоречие с аксиомой регулярности, сформулированной на внутреннем языке. Для теории этот объект фундирован; лишь в метатеории известно, что он был сконструирован через предельную цепочку. Это обеспечивает независимость таких объектов от теории. В действительности далее мы увидим, как придать подобным конструкциям более явную форму в рамках современных аксиоматик.

В контексте теоремы Кантора данное наблюдение представляет интерес со следующей стороны. Мы не можем утверждать, что описанный феномен является следствием формальной системы аксиом, поскольку на языке теории невозможно сформулировать соответствующее утверждение, которое могло бы быть доказано. По той же причине его нельзя ввести в качестве аксиомы теории или сделать частью логики первого порядка. Тем не менее существование объектов, невыразимых на языке теории, является её неотъемлемой частью. В частности, такая теория совместима лишь с метатеориями, также допускающими существование неопределимых объектов; в противном случае пришлось бы заключить, что теория противоречива. Мы не можем анализировать такую теорию, например, средствами теории категорий, если в рамках последней область описания по каким-либо причинам ограничена исключительно идентифицируемыми объектами. В рамках интуиционистского подхода в широком смысле мы вольны анализировать любые мыслеобразы, подвластные нашему определению, включая бесконечность, - при условии, что мы способны сформулировать то, что понимаем под этим понятием. Единственное ограничение состоит в том, что мы не вправе делать утверждения об объектах, которые сознание не умеет воображать: все объекты теории должны быть определимы; утверждения о неидентифицируемых объектах не проверяемы и, с точки зрения интуиционизма, лишены смысла. Сама по себе аксиоматика теории множеств вполне может рассматриваться в контексте интуиционизма: каждая аксиома представляет собой некоторый мыслеобраз, а аксиома бесконечности - идею множества, замкнутого относительно отношения следования. Однако при анализе теоремы Кантора можно возразить следующее. Используя Гёделеву нумерацию, мы способны пронумеровать все предельные формулы вещественных чисел и тем самым получить счётное множество всех иррациональных чисел (именно это, по сути, делает Анджей Буркет \cite{diagonal-antinomy}). Следовательно, в теории можно доказать как утверждение, так и его отрицание, то есть теория оказывается противоречивой. При этом нет необходимости подниматься на метауровень: метауровень как необходимая сущность возникает лишь в контексте формализма, истолковывающего современные аксиоматические теории множеств. По существу данное рассуждение связано с неполнотой теории множеств: утверждение о существовании невыразимых объектов недоказуемо в теории, недоказуемо также его отрицание, однако оно может быть только истинным, поскольку противное делает теорию противоречивой. Можно было бы сказать, что такие объекты требуют бесконечных формул в метатеории. Однако бесконечные формулы у нас уже имеются: в большинстве случаев в качестве метатеории используется сама теория множеств, которая прекрасно оперирует подобными объектами. Проблема заключается в том, что как только допускаются бесконечные формулы, бесконечные ряды $\sum_{i \in \omega} a_i 2^i$, где $a_i \in \{ 0,1 \}$, становятся полноправными объектами теории. С точки зрения современной интерпретации все такие ряды с бесконечным числом ненулевых коэффициентов тождественны $\omega$, поскольку внутреннее и внешнее сложение, различавшиеся Кантором, не разграничиваются, а замена через универсальную операцию объединения строго замкнута внутри $\omega$. Чтобы выйти за пределы $\omega$, требуется именно внешняя операция, которая не замкнута внутри, позволяя конечным слагаемым формировать различные бесконечные суммы. Точно также, в случае отрицательных степеней разные бесконечные последовательности, оказываются эквивалентными одному вещественному числу заданному аддитивным идеалом нуля. В результате континуум остаётся эффективно счётным: сквозного трансфинитного аппарата нет. Даже введя бесконечные формулы, мы не располагаем средствами для записи невыразимых объектов континуума. Данный факт связан вовсе не с мощностью, которая относительна по своей природе (вследствие теоремы Лёвенгейма–Сколема \cite{skolem-set-theory}), а с тем, что континуум обладает предикативным порядком, принципиально отличным от вполне упорядочения. Вследствие этого одну предикативную структуру невозможно выразить через другую.

Как бы то ни было, факт существования неопределимых объектов, являющийся частью интерпретации формальной системы - набора методологических правил и допущений, принимаемых для обеспечения непротиворечивости и согласованности теории, - накладывает ограничения на возможность некоторых доказательств. Действительно, если в рамках теории предполагается существование объекта, удовлетворяющего синтаксической системе уравнений, то по существу делаются два предположения: 1) требуемый объект существует и 2) он идентифицируем в теории с помощью формулы. Фактически записав синтаксическую систему уравнений, мы ищем семантику объекта удовлетворяющего им, но если в теории есть невыразимые объекты, мы можем иметь ситуацию, когда такие объекты структурированы как некоторые семантические инварианты, введение которых в теорию позволяет выразить искомый объект через них. Если из предположения о существовании такого объекта выводится противоречие, то нельзя в полной мере воспользоваться законом исключённого третьего: в строгом смысле неверно лишь одно из двух предположений - либо объект действительно не существует, либо не существует объекта, выразимого в известном синтаксическом аппарате теории (как мы увидим далее, формулировку можно усилить: либо не существует объекта, выразимого в синтаксическом аппарате теории с известным семантическим значением). В этой связи невозможно не вспомнить Лёйтзена Брауэра, с исключительной проницательностью указавшего, что бесконечность вовсе не обязательно должна описываться законом исключённого третьего \cite{brouwer}. Феномен невыразимых объектов даёт строгий критерий того, когда данный закон может не работать. В частности, нельзя быть уверенным, существует ли максимальное натуральное число, которое просто не может быть определено в теории, или же оно определимо, но не существует, чему мы действительно находим подтверждение в нестандартном анализе, где такое числа существует во внутреннем представлении. Доказанные нами теоремы несуществования теории множеств имеют куда более скромный смысл, чем им обычно приписывают: они утверждают лишь, что в рамках определённых нами семантических конструкций некоторые синтаксические уравнения не имеют решения, что существенно отличается от отсутствия решений в самой теории. Невыразимость объектов в теории - простое следствие использования такого инструмента, как импредикативность, необходимого для замыкания алгебраических структур. В то же время этот инструмент таит опасность: некоторые противоречия теории могут сместиться в зону импредикативности, вследствие чего теория окажется неспособной вывести их явным образом. Что будет препятствовать развитию и способствовать идеологизации определённой интерпретации, что в конечном итоге приведёт к кризису. Описанный феномен имеет непосредственное отношение к недопониманию между Куртом Гёделем и Людвигом Витгенштейном, вызванному аналогичным дрейфом интерпретаций за пределы формальной системы, - когда теория позволяет установить истинность высказываний, лежащих за границами её доказательной силы (если они не истинны, теория становится противоречивой) \cite{wittgenstein-godel}.

Потребность в теореме Кантора возникает вследствие импредикативности бесконечности, которая необходима для представления континуума в анализе, а также в качестве свойства строго разделяющего конечное и бесконечное на начальных этапах развития теории множеств. Теория Кантора была примечательно именно высоким уровнем контроля за балансом предикативности и импредикативности. Без рефлексивности множество степени естественным образом превосходит исходное множеств, поскольку помимо синглтонов всех элементов содержит огромное количество других множества. Данное наблюдение вскрывает ещё одну линию атаки на доказательство, о которой Ходжес не счёл нужным написать. В частности, теорема Кантора имеет смысл только в том случае, если мы предполагаем свойство рефлексивности бесконечных множеств, в противном случае один единственный элемент способен привести к невозможности биекции. Перед Кантором стояла сложная задача понять какой именно структурный паттерн предикативности даёт необходимые свойства анализа, которую он блестяще решил построив конструктивные объекты - порядковые типы, которые позволяют воспроизвести подобный экспоненциальный паттерн. В результате, первый трансфинитный порядковый тип $\omega$ обладает весьма интересными свойствами: с одной стороны, он не может аннотировать никакой свой начальный сегмент, с другой стороны способен биективно отображаться на собственную часть. Такая комбинация свойств достигается благодаря некоммутативности конкатенации упорядоченных отрезков, которая легла в основу разделения левых и правых операций. Давайте попробуем разобраться, получилось ли в рамках аксиоматических систем воспроизвести эту красивую находку Кантора. В ZFC или NBG, действительно, не сложно показать, что благодаря аксиоме выбора и рефлексивности $\omega$ всякое бесконечное множество рефлексивно \cite{fraenkel-reflection-choice}. В самом деле, рассмотрим произвольное бесконечное множество $X$ по теореме Цермело, его можно вполне упорядочить и взять начальный сегмент изоморфный $\omega$: $\{ x_0, x_1, x_2, \dots \}$. После чего определим функцию $\varphi : X \to X \setminus \{ x_0 \}$: $\varphi(x_i) = x_{i + 1} \; \forall x_i \in S$ и $\varphi(x) = x$ в остальных случаях, легко видеть что, данная функция является биекцией сопоставляющей множество $X$ с его частью $X \setminus \{ x_0 \}$. То есть, если множество $\omega$ рефлексивно, то таким же свойством обладают все бесконечные множества. Перейдём к доказательству рефлексивности множества $\omega$ \cite{fraenkel-reflection-omega}.

Теорема: {\sl Множество $\omega$ рефлексивно.}

Доказательство: 

На множестве ординалов операция $n + 1$ определена через отношение следования $s(n) = n \cup \{ n \}$. В соответствии со схемой аксиом преобразования (см. \ref{axiom:replacement}) для формулы $\phi(x, y): y = x \cup \{ x \}$ совокупность значений $y$ является множеством $Y$, а поскольку по аксиоме бесконечности (см. \ref{axiom:infinity}) $\omega$ замкнуто по отношению следования, как наименьшее индуктивное множество, то $Y \subset \omega$. Таким образом, функция $s : \omega \to Y$ задаёт биекцию $\omega$ со своей частью. Теорема доказана.
\begin{equation}
\forall A \; ( [ \forall x \in A \; \exists! y \; \phi(x, y) ] \Rightarrow \exists B \; \forall y \; [ y \in B \Leftrightarrow \exists x \in A \; \phi(x, y) ])
\label{axiom:replacement}
\end{equation}
\begin{equation}
\exists I (\emptyset \in I \wedge \forall x (x \in I \Rightarrow x \cup \{ x \} \in I))
\label{axiom:infinity}
\end{equation}

Существует множество доказательств данного факта, но во всех случаях, как и здесь, ключевую роль играют схема аксиом преобразования и аксиома бесконечности. Обратим внимание, что даже в формульной записи уже возникают тонкости, действительно, мы можем утверждать, что $\forall x \in \omega \; x \cup \{ x \} \neq \omega$, но в то же время $\bigcup\limits_{ x \in \omega } x = \omega$. Первое утверждение принято трактовать как отсутствие максимального конечного ординала, но в действительности, гораздо важнее то, что множество $\omega$ отсутствует в области значений определённого нами отношения следования $s(n) = n \cup \{ n \}$. В самом деле, предположим, что $\omega$ является значением отношения следования, то есть для формулы $\phi(x, \omega): \omega = x \cup \{ x \}$ найдётся $x = \upsilon$, делающий её истинной, что означает $\upsilon \in \omega$. Но тогда в силу аксиомы бесконечности $\omega = \upsilon \cup \{ \upsilon \} \in \omega$, пришли к противоречию с аксиомой регулярности, следовательно область значений отношения следования не содержит $\omega$. Иными словами, объединением конечных ординалов множество $\omega$ получить невозможно, так как такое объединение просто не может иметь своим значением $\omega$. Но тогда наше второе соотношение $\bigcup\limits_{ x \in \omega } x = \omega$ не имеет смысла, то есть предел индуктивной последовательности ординалов $\lim n = \bigcup x$ может существовать и не быть равным $\omega$ или вообще не существовать как объект теории. В современной интерпретации, по какой-то причине мы избегаем рассмотрения данного противоречия. С одной стороны, отношения следование не может иметь $\omega$ в области значений, с другой - оно явно имеет его. В действительности, никакого противоречия нет, поскольку как справедливо указывает Кантор речь идёт о двух принципиально различных типах операций: левых (внутренних) обеспечивающих замыкание и отсутствие $\omega$ в образе, и правых (внешних) обеспечивающих сквозную индукцию. Однако, вместо предложенного Кантором решения, мы строим некоторый неявный каркас методологических ограничений, за пределы которого не позволяем себе выходить. Чтобы импредикативность унаследованная из аксиомы бесконечности коллапсировала в рамки требуемой нами семантики, нам необходимо соглашение, которое позволило бы доказать минимальность множества $\omega$. Такому доказательству препятствует всё та же импредикативность бесконечности, вследствие которой мы не можем гарантировать стабилизацию бесконечного пересечения множеств даже в контексте аксиомы выделения. Мы принимаем предположение о предикативности минимальных объектов как базовое свойство нашего универсума, поддерживая его аксиомой регулярности. Однако, мы не доказываем, что в рамках принятой системы аксиом невозможна альтернативная интерпретация, которая которая обеспечивает коллапс импредикативности таким образом, что синтаксическая формула выражающая намерение выделить наименьшее бесконечное подмножество $\omega$ даёт объект не обладающий универсальным свойством минимальности. Поэтому мы называем такого рода предположения {\it методологическими допущениями} или {\it методологическими аксиомами}. У Кантора данное методологическое затруднение разрешается элегантно. Он вводит две независимые алгебраические структуры связанные с кардинальными и порядковыми числами. Так, арифметика кардинальных чисел замкнута, сколько бы мы ни складывали конечные кардинальные числа мы всегда получаем конечное кардинальное число, а бесконечное кардинальное число $\aleph_0$, являющееся совокупностью всех конечных, не может быть представлено их суммой \cite{cantor-cardinal}. В этом смысле, первое трансфинитное кардинальное число минимально, как класс эквивалентности множеств обладающих данным свойством рефлексивности. В то же время, в арифметике порядковых чисел появляется правое (внешнее) сложение, которое непосредственно связано с понятием предела $\lim n_i = n_1 + \sum (n_{i + 1} - n_i)$ \cite{cantor-ordinal}, в частности, $\omega = \lim n = \sum 1$, то есть наименьшее трансфинитное порядковое число является бесконечной суммой единиц отличной от любого конечного числа. При этом такое трансфинитное число является классом эквивалентности множеств по левой (внутренней) операции сложения, что также обеспечивает минимальность порядкового типа как рефлексивного класса эквивалентности, обладающего заданными свойствами, а не как фиксированного множества. В рамках аксиоматической теории множеств мы могли бы сделать тоже самое и $\forall x \in \omega \; x \cup \{ x \} \neq \omega$ выполнялось бы в кардинальном смысле, в то время как $\lim x = \bigcup\limits_{ x \in \omega } x = \omega$ в ординальном смысле. В таком случае, рефлексивность была бы свойственна кардинальным операциям, и левым ординальным, а правые ординальные операции были бы нерефлексивными, позволяя получать предельные ординалы и выходить за их границы. При этом вовсе не обязательно, чтобы одни и те же формулы теории имели два разных смысла в зависимости от контекста, Кантор уже предложил синтаксическое разделение на левое и правое, которое прекрасно работает. Методологические аксиомы довольно распространённое явление, например, определение предела в анализе связано с такими предположениями. В целом, их допустимо считать частью интерпретации аксиом и в большинстве случаев нет смысла акцентировать на них внимания, но наша попытка определить предельное множество через интерпретацию универсальной фундированности аксиоматики является исключением. К сожалению, последствия данной методологической редукции оказываются чересчур серьезными, ввиду того, что аксиоматики теории множеств прочно заняли позиции оснований математики. Более того, точная формулировка этих допущений вызывает затруднения, поскольку утверждается, что ординалы изоморфны порядковым типам, импредикативность которых ограничена левыми операциями, что довольно сложно согласовать с ординалами рефлексивность которых универсальна.

Отвлечёмся, однако, от минимальности. Она действительно присутствует в индуктивном фрагменте аксиоматических теорий множеств и позволяет говорить о корректности всех утверждений, выводимых из неё в данном фрагменте. Обобщение минимальности на весь домен в строгом формальном смысле неправомерно (является следствием из формальной ошибкой преждевременного обобщения): оно существенно искажает смысл аксиом и усложняет анализ достоверности теоретических выводов, что приводит к идеологизации и необоснованной маргинализации альтернативных идей. В рамках настоящей статьи мы лишь указываем на данный архитектурный недостаток теории как на один из элементов, препятствующих обнаружению скрытых предположений, ответственных за нарушение изоморфизма с построениями Кантора. Усомнимся в рефлексивности бесконечных множеств и рассмотрим альтернативные размышления, не позволяющие считать соответствующее доказательство, приведённое выше, безупречным. Следующие рассуждения будут довольно тонкими, и может показаться, что методологическая проблема, которой они касаются, не столь существенна. В определённой степени это верно: рефлексивность является неотъемлемой частью импредикативного фрагмента теории множеств, а её обобщение на весь домен также неправомерно и представляет собой проявление того же архитектурного дефекта. Наша цель, однако, состоит в том, чтобы подчеркнуть невероятную сложность понятийного аппарата, возникающего из, казалось бы, простых аксиом, - сложность, которая в конечном итоге успешно маскирует неявные допущения, лежащие в основе расхождения современной интерпретации аксиоматических теорий с оригинальной системой Кантора. Из всех возможных формулировок аксиомы бесконечности нам, приключилось выбрать одну из самых неудачных. Это связано с тем, что формула, используемая в данной аксиоме, в определённом смысле родственна формуле Рассела: она обладает с ней общими свойствами и приводит к противоречиям в некоторых контекстах. На этом мы подробно остановимся далее; станет понятно, что для исправления ряда проблем современных аксиоматических теорий достаточно заменить формулировку аксиомы бесконечности. Неопределённость, свойственная таким выражениям, связана не с формулами как таковыми, а скорее с нашим непониманием их смысла. В действительности рефлексивность не имеет никакого отношения к формуле аксиомы бесконечности, и было бы весьма странно пытаться сохранить данную интерпретацию даже в рамках ортодоксальной теории. Дело в том, что замкнутость операций над элементами внутри математической структуры и замкнутость операций над самой математической структурой - два совершенно различных свойства. Это обусловлено тем, что замкнутость внутренней составляющей не позволяет выйти за пределы математической структуры для описания внешних свойств. В свою очередь, данное обстоятельство связано с онтологическими ограничениями нашего мышления, в силу которых математические теории представляют собой идеализированные аннотации или описания явлений, а не сами эти явления. Например, одна и та же формальная теория может описывать (аннотировать) объекты совершенно разных моделей. Отождествление алгебраических абстракций с сущностями, наблюдаемое в аксиоматических теориях множеств, является примером довольно распространённой в математике ошибки реификации. Впрочем, эпистемологические аспекты отождествления абстракций и сущностей, приводящие к возникновению устойчивых иллюзий, формирующих догмы и, в конечном итоге, идеологические системы, заслуживают отдельного обсуждения. Когда мы говорим о замкнутости по отношению следования, мы подразумеваем, что образы элементов замкнуты в некоторой алгебраической структуре, описывающей устройство некоторого объекта - совокупности или класса; такую аннотацию мы называем множеством. Говоря же о замкнутости множества относительно преобразования, мы имеем в виду, что трансформации подвергается сам объект, устройство которого аннотирует наша алгебраическая структура. В силу своей внутренней замкнутости эта структура не способна описывать внешние свойства и подобное преобразование. Так, отношение следования $x \cup \{ x \}$ формирует индуктивную последовательность, которая вовсе не обязательно имеет предел в рамках нашей алгебраической структуры: она может расходиться, продолжаться неограниченно и не образовывать законченного объекта. Это важное наблюдение в контексте последующего обсуждения: аналогичным образом арифметика Пеано PA ''живёт'' внутри $\omega$, но не определяет его как объект. Аксиома бесконечности утверждает, что существуют объекты, к которым применима наша множественная аннотация, замкнутая по отношению следования, однако при этом не отождествляет их с пределом последовательности, образованной данным отношением; этот предел, вообще говоря, может быть элементом внешнего, а не внутреннего описания. Аксиома довольно явно указывает на необходимость дополнительных предположений относительно природы таких объектов - предположений, которые могли бы согласовать внутренние и внешние аспекты их существования. По какой-то причине в рамках современных теорий, делая такие предположения, мы не придаём им формального характера.

Для того, чтобы сделать наши рассуждения более наглядными введём некоторые вспомогательные понятия. Совокупность объектов индуктивного построения будем называть {\it областью регулярности}. Иными словами, всё, что может быть сформировано многократным повторением элементарных операций и их комбинаций, составляет область регулярности. В частности, все конечные ординалы являются элементами такой области, и именно здесь происходит замыкание отношения следования. Область регулярности не является множеством, поскольку не существует в законченной форме: она воплощает идею потенциальной бесконечности, импредикативность которой и обеспечивает замкнутость её алгебраической структуры Нетрудно видеть, что счётные бесконечные ординалы $\omega, \omega + 1, \omega + 2, \dots, 2\omega, \dots, \omega^2, \dots, \omega^\omega, \dots$ также образуют область регулярности, замкнутую относительно алгебраических операций. Переход между областями регулярности будем называть  {\it областью сингулярности}, или просто {\it сингулярностью}. В частности, конечная и бесконечная области регулярности разделены такой областью сингулярности. Согласно аксиоме бесконечности ни один объект конечной области регулярности не может непосредственно предшествовать бесконечному множеству; отсюда возникает потребность в термине, позволяющем атрибутировать свойства, связывающие области регулярности. Отметим, что пока мы не делаем никаких дополнительных предположений, а лишь формируем язык, позволяющий описывать связь областей регулярности как сущность, которая разделяет конечное и бесконечное и служит мостом между ними, а вместе с тем определяет смысл понятия ''все элементы'', отсутствующий у области регулярности. Это позволяет связать преобразование всех элементов именно с сингулярностью. Наша задача состоит в том, чтобы показать: рефлексивность, сопряжённая с преобразованием всех элементов множества, является свойством, отличным от замкнутости алгебраической структуры и независимым от аксиом современных теорий множеств. Для этого необходимо описать сингулярность таким образом, чтобы бесконечное множество оставалось нерефлексивным, но при этом удовлетворяло всем аксиомам теории множеств. Иными словами, требуется привести контрпример, иллюстрирующий независимость замкнутости и рефлексивности. Нужно смоделировать свойства сингулярности так, чтобы в результате применения схемы аксиом преобразования к множеству $\omega$ действительно получалось множество значений $\omega \setminus \emptyset$, однако сингулярность области определения не формировалась бы. В этом случае функция $f : \omega \to \omega \setminus \emptyset$ оказывается невозможной. Если же преобразовать элементы $\omega$ непосредственно в пары $(i, i+1)$, то сингулярность формируется и функция как множество пар существует вследствие схемы аксиом преобразования, но областью её значений оказывается множество $(\omega + 1) \setminus \emptyset$, что соответствует $f : \omega \to (\omega + 1) \setminus \emptyset$. 

Сконструируем теперь модель, удовлетворяющую указанным требованиям. Возьмём за основу стандартную кумулятивную иерархию и сделаем в модели дополнительное предположение: всякое бесконечное множество $X$ содержит сингулярности $\mathcal{S}(X)$ в областях разрыва непрерывности предикативного описания. Формально достаточно привязаться к аксиоме бесконечности и потребовать, чтобы всякому ординалу, не имеющему непосредственного предшественника, соответствовала такая область сингулярности: $\forall x \in Ord \; (\nexists y \; (y \cup \{ y\} = x) \to \exists \mathcal{S}(x))$. Сингулярность $\mathcal{S}(X)$ характеризуется совокупностью атрибутов, описываемых омнифическими целыми числами, нумеросити или гросс-объектами, попадающими в данную область при сопоставлении с ординалами. В действительности для описания таких атрибутов достаточно множества конечных ординалов, что позволяет избежать введения дополнительного типа объектов и выхода в многосортную логику. Однако наша цель в данном контексте - сделать построение как можно более наглядным, а не сформулировать законченный формализм теории; поэтому мы воспользуемся готовым хорошо изученным классом омнифических целых чисел. Полноценные аксиоматики, в которых реализуются подобные построения, станут предметом рассмотрения следующей статьи серии. Кроме того, совокупность атрибутов можно сопоставить непосредственно предельным ординалам, не вводя вспомогательной сущности. Тем не менее понятие сингулярности лучше раскрывает смысл импредикативности, разрывающей непрерывность свойств. Класс ординалов ${\backslash}Ord$, образующий полукольцо, изоморфно вкладывается в подклассы положительных элементов колец омнифических целых чисел ${\backslash}Oz$, сохраняя нормальную форму Кантора; при этом кольцо содержит полукольцо ординалов ${\backslash}Ord$ как подструктуру \cite{conway-omnific-integers}. Нетрудно видеть, что обобщение гросс-объектов, для которых постулируются арифметические операции, совпадающие с таковыми для натуральных чисел, также приводит к алгебраической структуре кольца, для которой класс ординалов становится подструктурой, когда операция возведения в степень определяется в ординальном смысле. Будем обозначать сюрреальные числа, соответствующие ординалам, точками над символом ординала; например, ординалу $\alpha$ соответствует сюрреальное число $\dot{\alpha}$. Что позволяет нам доказать теорему:

Теорема 6: {\sl Рефлексивность $\omega$ независима от ZFC (NBG).}

Доказательство: 

Воспользуемся тем фактом, что класс ординалов является подструктурой класса омнифических целых чисел и сформулирем дополнительную к ZFC или NBG аксиому модели:
\begin{equation}
\begin{gathered}
\forall \alpha \bigl(\emptyset \in \alpha \wedge \forall x (x \in \alpha \Rightarrow x \cup \{ x \} \in \alpha) \wedge \nexists y \; (y \cup \{ y\} = \alpha)\bigr) \Rightarrow \\
 \forall \dot{x} (\dot{x} > 0 \wedge \dot{x} < \dot{\alpha} \Rightarrow x \in \alpha \vee \dot{x} \triangleleft \mathcal{S}(\alpha))
\label{axiom:singularity}
\end{gathered}
\end{equation}

Здесь $x \triangleleft \mathcal{S}(\alpha)$ означает, что $x$ является атрибутом сингулярности $\mathcal{S}(\alpha)$. Вместо сюрреального числа $\dot{\alpha}$ можно использовать гросс-объект $\raisebox{.5pt}{\textcircled{\scalebox{0.8}{$\alpha$}}}$ или любой другой объект индексной алгебры. Удобнее всего рассматривать атрибуты сингулярности $\mathcal{S}(\alpha)$ как мнимые множества - объекты, к которым применимы некоторые множественные операции, но которые не являются полноценными множествами исходной теории. Все операции для сюрреальных чисел уже заданы, поэтому нам не требуются дополнительные аксиомы. Такой подход позволяет легко перенести существующую нотацию множеств на интерпретацию с сингулярностью. В частности, $\omega = \{0, 1, 2, \dots, \dot{\omega} - 2, \dot{\omega - 1}\}$, где реальные множества и атрибуты сингулярности разделены троеточием. Это даёт интуитивное понимание операций над мнимыми множествами; например, $(\dot{\omega} - 3) \cup \{ \dot{\omega} - 3 \} = (\dot{\omega} - 3) + 1 = \dot{\omega} - 2 = \{0, 1, 2, \dots, \dot{\omega} - 3 \}$. Рассмотрим взаимодействие аксиом теории множеств, используемых в доказательстве рефлексивности, с объектами, обладающими сингулярностью, заданной бесконечным множеством атрибутов. Применяя схему аксиом преобразования с формулой $\phi(x, y) : y = x \cup \{ x \}$, мы действительно получаем множество образ $Y = \omega \setminus \emptyset$: для всякого элемента $x$ схема требует существования $y$, удовлетворяющего формуле, и, наоборот, для всякого элемента $Y$ существует элемент $\omega$, связанный с ним через $\phi$. Однако поведение атрибутов сингулярности обоих множеств при этом остаётся неопределённым. В схеме аксиом преобразования невозможно потребовать, чтобы формула выполнялась для всякого атрибута сингулярности множества $A$, поскольку $\dot{\omega} - 1 \cup \{\dot{\omega} - 1\} = \omega$, а если $\omega \in Y$, то должно существовать множество, связанное с ним формулой; между тем $\dot{\omega} - 1$ является мнимым множеством. Тем не менее согласованности области сингулярности множества $B$ с формулой можно ожидать:
\begin{equation}
\begin{gathered}
\forall A \; ( [ \forall x \in A \; \exists! y \; \phi(x, y) ] \Rightarrow \exists B \; \forall y \; [ y \in B \Leftrightarrow \exists x \in A \; \phi(x, y) ]) \Rightarrow \\
 \forall \dot{y} (\dot{y} \triangleleft \mathcal{S}(B) \Leftrightarrow \exists \dot{x} \triangleleft \mathcal{S}(A) \; \phi(\dot{x}, \dot{y}))
\label{axiom:singularity-replacement}
\end{gathered}
\end{equation}

Следовательно, несмотря на то что множество $Y = \omega \setminus \emptyset$, действительно существует как того требует схема аксиом преобразования, его невозможно связать с множеством пар, соответствующим функции $f : \omega \to \omega \setminus \emptyset$. Сингулярность множества пар $\mathcal{S}(f)$ не может иметь атрибут вида $(\dot{\omega} - 1, x)$: если бы такой $x$ существовал, то в силу формулы $\phi$ он был бы равен $\omega$, однако атрибуты сингулярности $\mathcal{S}(\omega \setminus \emptyset)$ могут быть только мнимыми множествами, а не реальными. Таким образом, в данном случае речь может идти лишь о функции $f: \dot{\omega} - 1 \to \omega \setminus \emptyset$ отображающей мнимое множество в реальное и представленной множеством пар $\{ (0, 1), (1, 2), \dots, (\dot{\omega} - 3, \dot{\omega} - 2), (\dot{\omega} - 2, \dot{\omega} - 1) \}$. Если же воспользоваться формулой $\phi(x, y) : y = (x, x \cup \{ x \})$ для образования множества пар, то в рамках наших предположений может существовать также множество пар $\{ (0, 1), (1, 2), \dots, (\dot{\omega} - 3, \dot{\omega} - 2), (\dot{\omega} - 2, \dot{\omega} - 1), (\dot{\omega} - 1, \omega) \}$, где сингулярность имеет атрибут $(\dot{\omega} - 1, \omega)$, являющийся мнимым множеством. Это соответствует функции $f : \omega \to (\omega + 1) \setminus \emptyset$. Иными словами, в рамках полученной системы аксиом единственной функцией, связанной с отношением следования и отображающей реальное множество на реальное, является $f : \omega \to (\omega + 1) \setminus \emptyset$. Данный факт связан с тем, что стандартная форма схемы аксиом преобразования не связывает множества с множеством пар, которое можно интерпретировать как функцию. Нетрудно убедиться, что, используя схемы аксиом выделения и преобразования, невозможно прийти к противоречию, поскольку мы вводим лишь дополнительные символы не интерферирующие с базовым языком теории. Это свидетельствует о том, что наше расширение полностью удовлетворяет стандартной аксиоматике $\mathrm{ZFC}$. С помощью схемы аксиом выделения из множества пар $f : \omega \to (\omega + 1) \setminus \emptyset$ в принципе можно получить $f : \dot{\omega} - 1 \to \omega \setminus \emptyset$, поскольку требуется лишь существование произвольного множества, удовлетворяющего условию. Однако в данной схеме аксиом, как и во всех остальных, мы можем оперировать только реальными множествами, что не позволяет вывести противоречие для мнимых множеств. Интересно, что схема аксиом преобразования в текущем виде не способна дать область значений $(\omega + 1) \setminus \emptyset$: элемент $\omega$ требовал бы реального множества в качестве прообраза, которого не может существовать; в то же время мнимые множества не могут иметь реальные множества в качестве образа. Это указывает на то, что для полноценной работы с данным расширением потребуются дополнительные аксиомы. Мы также использовали интерпретацию функций, учитывающую сингулярность. Поскольку функции не связаны с аксиомами, а вводятся через определения, не составляет труда адаптировать определения функций к объектам с сингулярностью. Таким образом, построенная модель интерпретирует ZFC и NBG, то есть свидетельствует о существовании соответствующего непротиворечивого расширения этих теорий, и при этом допускает нерефлексивность $\omega$. Теорема доказана. 

Не вызывает трудностей построить подобную модель и без привлечения индексных множеств: достаточно лишь предположить, что сингулярность (или предельные ординалы как точки разрыва непрерывности предикативности) характеризуется некоторым множеством $S(X)$. Для ординала $\omega$ характеристическим множеством может служить само $\omega$, в котором множество $n$ интерпретируется как свойство $\omega - n$. Такая интерпретация потребует лишь более тщательной проработки вспомогательных аксиом, что не имеет большого смысла, поскольку мы не планируем использовать данную теорию. Кроме того, выход за пределы логики первого порядка можно обосновать тем, что ZFC и NBG (как и любая достаточно сложная система, интерпретирующая арифметику) не могут быть последовательно истолкованы в такой логике вследствие отсутствия базовых понятий, в частности понятия истины. Нам требуется допущение о существовании невыразимых объектов, которое является аксиомой высшего порядка. Смысл приведённого построения состоит в том, что мы используем уязвимость схемы аксиом преобразования, не связывающей существование образа с существованием множества пар. Это позволяет гарантировать существование объекта, удовлетворяющего аксиоме, без существования соответствующей функции. Данная теорема представляет собой не более чем простую констатацию факта: когда мы аксиоматизируем некоторые структурные построения, упрощая при этом онтологию, всегда остаётся возможность восстановить исходную полную онтологию - при условии, что в аксиомах не было допущено ошибки. Отметим ещё раз, что сама по себе рефлексивность не является источником противоречий. Если предположить, что помимо уже постулированных объектов никаких других сущностей не возникает, то рефлексивность выводима и оказывается универсальной (точнее, в таком случае отсутствует инструментарий для описания какой-либо внешней структуры, и логичным представляется расширение внутреннего описания на весь домен). Однако неправомерное позиционирование данного свойства в теории существенно затрудняет анализ и служит точкой расхождения между общепринятыми и некоторыми альтернативными взглядами. Это обстоятельство потребуется нам далее для того, чтобы собрать всю картину воедино. К вопросу об онтологии Кантора, которая легла в основу аксиоматик теории множеств, мы ещё вернёмся в рамках настоящей работы. Кроме того, в серии статей будет приведено ещё одно, существенно более прозрачное доказательство данного факта. Стоит обратить внимание на то, что мнимые множества в доказательстве позволяют избежать квантификации в контексте современных аксиоматик, особенно в случае аксиомы бесконечности, о которой речь пойдёт далее. Это выявляет любопытную особенность логики первого порядка, допускающей импредикативность. В результате в рамках некоторой аксиомы можно ввести класс новых объектов, определяемых через квантификацию, которая происходит в том числе по объектам нового типа: допускается определённая зацикленность, обеспечивающая существование объектов, связанных с квантификацией по ним самим. Данное наблюдение имеет прямое отношение к проблемам выбранной нами формы аксиомы бесконечности, о чём мы поговорим далее.

Мы воспользовались сюрреальными числами исключительно потому, что эти объекты хорошо известны и в достаточной мере исследованы. В то же время данный контрпример, демонстрирующий уязвимость вывода, как уже отмечалось, может быть сформулирован и без использования таких объектов. Например, даже простое переупорядочение множества $\tilde{\omega} = \{ 0, 2, 4, \dots, 5, 3, 1 \}$ приводит к тому, что необходимо определить образ следования для $1$. Это указывает на то, что рефлексивность - свойство порядка, то есть алгебраической структуры нашей аннотации (где оно называется подобием, как у Кантора и Дедекинда, чьи построения онтологически более корректны), а не непосредственно самого объекта. Подчеркнём, что мы доказали гораздо более сильную теорему, чем требуется в действительности, - исключительно для того, чтобы указать на уязвимость стандартной схемы аксиом преобразования. Чтобы закрыть данную уязвимость, достаточно изменить схему аксиом преобразования так, чтобы для всякого $y$ в полученном множестве $B$ существовала пара $(x, y)$, где $x \in A$; множество таких пар и стало бы требуемой функцией. Такое исправление безусловно улучшит ситуацию. Однако оно никак не повлияет на вторую и существенно более важную проблему: схема аксиом преобразования делает некоторое утверждение о связи множеств через объекты теории, но при этом не требует отсутствия альтернативных связей через другие объекты. Данная аксиома не закрепляет нашу интерпретацию наделяющую один класс объектов эксклюзивным статусом функциональных преобразований множеств. Иными словами, исправление действительно сделает невозможной сильную формулировку о независимости рефлексивности, поскольку рефлексивность в строго зафиксированном смысле станет неотъемлемой частью системы. Но оно не повлияет на возможность существования альтернативных смыслов функциональных связей между множествами, реализуемых через другие объекты, уже присутствующие в теории. Ключевая проблема заключается в том, что мы допускаем логическую ошибка преждевременного обобщения: одно свидетельство не может распространятся на всю теорию. Для такого обобщения необходимо доказать, что никакой другой класс объектов теории не может описывать функциональные отношения между множествами. Таким образом, принимая намерение воплощённое в схеме аксиом преобразования, рефлексивность действительно можно считать частью системы, но лишь по отношению к фиксированному классу объектов теории, интерпретируемых как функции. Следовательно, рефлексивность присуща некоторому фрагменту теории, а не всему универсуму. Такой класс объектов - множеств пар - мы называем {\it внутренними функциями}. Альтернативный класс объектов, которые не проявляют рефлексивности, мы называем {\it внешними функциями}. Обратим внимание, что внутренние и внешние функции непосредственно связаны с различными интерпретациями взаимно однозначного соответствия, которые в рамках современных аксиоматик объединены в одну сущность - биекцию, причём такое объединение аксиоматически не узаконено. Существенно то, что оба класса объектов уже присутствуют в теории, не требуя разделения множеств на стандартные и нестандартные: достаточно лишь восстановления утраченной аннотации внешних отношений из теории Кантора. Возможность доказать независимость появляется благодаря дополнительным сущностям - мнимым множествам, которые не описываются аксиоматикой, но и не запрещаются ею. Это отсылает нас к проблемам, возникавшим на самых ранних этапах формирования формализма теории множеств. В частности, Френкель предлагал ввести аксиому ограничения, постулирующую, что помимо множеств, существование которых абсолютно необходимо по аксиомам, других множеств не существует \cite{fraenkel-restriction-axiom}. Однако такого рода предположения довольно трудно обосновать, поэтому математическое сообщество ограничилось аксиомой регулярности, которая по замыслу должна обеспечивать универсальность индуктивной структуры. Аксиомы цензуры подобного рода допускают неверное истолкование, перекликающееся с духом современной интерпретации, в которой объекты и конструкции, плохо согласующиеся с принятыми убеждениями, безжалостно исключаются из рассмотрения. Стоит отметить, что исходная формулировка Френкеля аксиомы преобразования - если $M$ - множество, то заменяя каждый элемент $M$ ''объектом из области $B$'', мы снова получаем множество \cite{fraenkel-replacement-axiom}, мы снова получаем множество» \cite{fraenkel-replacement-axiom} - совершенно не годится для доказательства рефлексивности, поскольку не накладывает никаких структурных ограничений на полученное множество. Кроме того, необходимо сказать и о том, что в NBG, где схема аксиом преобразования является теоремой, вытекающей из аксиомы ограничения мощности и схемы аксиом неограниченного усложнения классов, рефлексивность множеств оказывается ещё более спорным утверждением. Это связано с тем, что структура классов значительно более разнообразна, чем структура множеств, в частности вследствие ограничений, накладываемых на структуру множеств теоремой Кантора.

В действительности, нам не требуются даже мнимые множества, для того чтобы разглядеть данный архитектурный недостаток теории. Аксиома бесконечности не декларирует никаких свойств сингулярности, связывающей конечную и бесконечную области: известно лишь, что эти области должны быть каким-то образом непротиворечиво связаны между собой. Чтобы доказать рефлексивность, необходимо показать, что множество $\omega$ не может каким-либо образом входить в сингулярность множества пар, образующих функцию $f : \omega \to \omega \setminus \emptyset$, или что тоже самое, что в теории не существует объектов содержащих $\omega$, которые можно интерпретировать как функциональные отношение такого рода. Для этого требуются дополнительные предположения относительно природы бесконечности. Свойство сингулярности переносить преобразование всех элементов из области регулярности в следующую область регулярности назовём {\it транзитивностью}. Мнимые множества нужны лишь для того, чтобы лучше рассмотреть детали механизма, ответственного за транзитивность. Когда мы преобразуем множество $\omega$, его реальная часть преобразуется стандартным образом: $0 \to 1, 1 \to 2, \dots, n \to n + 1, \dots$, в то время как транзитивность сингулярности обеспечивает одновременное преобразование его мнимой части: $\omega - (n + 1) \to \omega - n, \dots, \omega - 2 \to \omega - 1, \omega - 1 \to \omega$. В рамках современных аксиоматик мнимая часть остаётся невидимой; тем не менее сингулярность, возникающая вследствие импредикативности, связанной с аксиомой бесконечности, может обладать данным свойством. Благодаря этому образом множества $\omega$, преобразуемого внешним отношением следования, оказывается множество $\omega + 1$, в котором элемент $\omega$ является частью сингулярности. Вопрос сводится лишь к существованию объекта теории, выражающего данное отношение между множествами. Такой объект вполне может существовать, например в виде $f = (\{ (x, y) \mid x, y \in \omega \wedge (y = x + 1 \vee y = x = \{ \omega \}) \})$. Так или иначе, свойство транзитивности позволяет понять различие между замкнутостью алгебраической структуры относительно внутренних операций и рефлексивностью объекта самого по себе, аннотированного с помощью алгебраической структуры отношения принадлежности. Замкнутость преобразований внутри множества - это всего лишь свойство алгебраической структуры, которой мы аннотируем данное множество; сами преобразования остаются внутренними и не позволяют делать произвольные выводы о множествах как объектах. Для подобных выводов требуется обобщение преобразований, позволяющее им стать внешними. От того, как именно осуществляется такое обобщение - транзитивно или нет, - и зависит рефлексивность. Если исследуемая совокупность не обладает свойством транзитивности (вследствие чего этим свойством не обладает и сингулярность, описывающая её свойства в контексте теории множеств), то такая совокупность называется {\it рефлексивной}. Совокупность, обладающая свойством транзитивности, называется  {\it транзитивной}. Таким образом, в строгом смысле рефлексивность независима от современных аксиоматик и должна предполагаться явно. Формы такого предположения могут быть различны, включая отсутствие транзитивности области сингулярности при преобразовании области регулярности и существование биекции бесконечных множеств со своими частями. В случае принятия методологического допущения, что схема аксиом преобразования говорит именно о функциональной связи, рефлексивность становится неотъемлемой частью импредикативного фрагмента теории. Данное наблюдение отчётливо демонстрирует последствия отказа от онтологической обоснованности Кантора и Дедекинда: вне философского базиса формальные системы становятся существенно зависимыми от интерпретаций, возникающих при попытке заполнить импредикативные пустоты и добиться осмысленности всех синтаксических выражений. Интересно также, что существующие аксиоматики вполне могут оказаться совместимыми с 8-ой аксиомой Евклида: ''целое больше части''. Например, если мы постулируем транзитивность сингулярности, предположение которое, вообще говоря, логично, в силу того, что транзитивность является ключевым свойством ординалов, которое мы искусственно ограничиваем для случае рефлексивности предельных ординалов. 

Обращает на себя внимание весьма тесное переплетение обсуждаемых нюансов. Из аксиомы бесконечности мы заключаем, что в бесконечном множестве не существует максимального элемента. Это позволяет определить конечные множества как те, любое непустое подмножество которых имеет максимальный элемент. В свою очередь, данное определение позволяет найти одно из решений (но не единственное) предельного уравнения, заданного аксиомой бесконечности, - минимальное бесконечное множество $\omega$, которое оказывается рефлексивным, поскольку пара $(\omega - 1, \omega)$ не может существовать. Вместе с тем в теории появляются невыразимые объекты - объекты, которые невозможно ни определить, ни идентифицировать, лежащие за пределами нашего описания. Технически внутри теории мы не в состоянии дать ответ на вопрос: противоречивость конструкции $\omega - 1$ означает невозможность описываемой ею сущности или же неопределённость данного понятия в рамках существующей системы предположений (отсутствие семантического объекта среди уже зафиксированных понятий). Отметим, что в исходной теории Кантора $\omega - 1$ - строго определённый объект: множество без первого элемента, порядковый тип которого равен $\omega$, что уже свидетельствует о нарушении изоморфизма ординалов и порядковых типов. Такого рода неопределённости можно исследовать и попытаться закрыть новыми предположениями, делающими теорию однозначной. Вместо этого мы принимаем предположение, что противоречивость означает невозможность существования такого объекта. Это, в сущности, понятно: мы опираемся на классическую логику, а не на интуиционистскую. Однако последняя создавалась как раз с целью стать инструментом анализа проблемных мест математических теорий. Загвоздка состоит в том, что для полноценной опоры на классическую логику требуется однозначность семантического содержания отношения принадлежности: необходимо доказать, что теория не допускает альтернативных смыслов данного отношения в зависимости от контекста. Между тем Лёвенгейм и Сколем доказали ровно обратное, что ещё сильнее ослабило позицию мейнстрима. Если рассматривать выражение $\omega - 1$ как неопределённое, доказательства рефлексивности бесконечных множеств и минимальности множества $\omega$ становятся проблематичными. Дело в том, что если неопределённость $\omega - 1$ трактуется как отсутствие последнего элемента, то невозможность и неопределённость для нас существенно не различаются. В то же время в случае пары $(\omega - 1, \omega)$ мы не можем использовать её синтаксическую конструкцию, однако ничто не мешает ей реализоваться несмотря на нашу неспособность её описывать. Данные наблюдения имеют непосредственное отношение к форме аксиомы бесконечности. Если заменить аксиому бесконечности на более абстрактную (''существует бесконечное множество''), то окажется невозможным сделать вывод об отсутствии максимальных элементов, о минимальности $\omega$ и о его рефлексивности. Дело здесь вовсе не в необходимости определить, что такое бесконечное множество (хотя это, разумеется, играет немаловажную роль при аксиоматизации теории, в которой бесконечность возникает, но по какой-то причине её статус игнорируется). Дело в том, что аксиома бесконечности обладает побочными эффектами - дополнительными утверждениями, которые кажутся нам логическими следствиями из неё. Как мы увидим далее, существует множество формул, определяющих бесконечные множества, но при этом не имеющих побочных эффектов. В связи с этим обращает на себя внимание следующее. Аксиомы формальной системы должны быть минимальными или атомарными (постулировать ровно одно свойство или феномен), чистыми (не накладывать условий, кроме абсолютно необходимых), однозначными (не допускать различных интерпретаций) и независимыми (не пересекаться друг с другом). Эталоном такой системы могут служить постулаты Евклида: когда через тысячелетия мы столкнулись с проявлениями геометрии за их пределами, мы просто исключили пятую аксиому из списка предположений. Но присмотримся к аксиоме бесконечности. Её смысл заключается в том, чтобы задекларировать свойство замкнутости алгебраической структуры. Однако выбранная нами форма не является независимой: в неочевидном переплетении со схемами аксиом выделения и преобразования она порождает побочные эффекты, такие как отождествление замкнутости и рефлексивности. Она содержит неоднозначность, проявляющуюся как в теореме Лёвенгейма–Сколема, так и в невозможности сделать вывод о том, является ли бесконечное множество результатом объединения конечных или же никак не связано с индуктивным построением его элементов. Она не является чистой, поскольку постулирует не только алгебраическую замкнутость, но и существование объекта, в рамках которого эта замкнутость реализуется, а кроме того определяет конкретную структуру реализации. Наконец, она не является атомарной: из неё выводится отсутствие максимального элемента, а в комбинациях с другими аксиомами - и другие побочные эффекты. Иными словами, данное предположение в стандартной интерпретации не удовлетворяет почти всем требованиям, предъявляемым к аксиомам формальных систем. Его побочные эффекты вступает в противоречие с фундаментальным принципом Евклида о том, что целое больше своих частей. При этом оно объявляет рефлексивными сами объекты теории, а не их аннотации - в отличие от работ Кантора и Дедекинда, исследовавших подобие алгебраических структур (причём Дедекинд даже явно фиксирует предположение подобия в определении). Но и это ещё не всё: теорию, основанную на такой интерпретации данного предположения, мы нарекаем основанием математики и помещаем в фундамент, откуда позволяем ей расползтись по всем областям. Сложно вообще дать какой-либо комментарий относительно этого немотивированного упрямства.

Проблемность выбранной формы аксиомы бесконечности можно отчётливо и математически строго продемонстрировать. Намёки на это проявляются даже в давно существующих теориях, например в ''Новых основаниях'' (NF) Уилларда Куайна \cite{quine-new-foundations}. В этой системе принимаются всего две аксиомы: аксиома экстенсиональности в стандартной форме и схема аксиом неограниченного свёртывания для стратифицированных формул. Формула считается стратифицированной, если каждой переменной можно присвоить натуральный номер (тип) так, чтобы отношение принадлежности увеличивало данный номер, а равенство сохраняло его. Легко видеть, что формула стандартной аксиомы бесконечности $\emptyset \in I \wedge \forall x \; (x \in I \Rightarrow x \cup \{ x \} \in I)$ не является стратифицированной. Действительно, $x \cup \{ x \} = \{ y \mid y \in x \vee y = x \}$, и тип $y$ должен одновременно равняться типу $x$ и быть меньше его. Аналогично обстоит дело с формулой Рассела $x \notin x$. Отметим, что исходная формула аксиомы бесконечности циклична; в NF допустимо преобразовать её к виду $\{ x \mid \exists y \; (\emptyset \in y \wedge \forall z \; (z \in y \Rightarrow z \cup \{ z \} \in y)) \wedge x \in y \}$, который также не проходит стратификацию. Стоит также отметить, что на сегодняшний день непротиворечивость NF установлена относительно теории запутанных типов \cite{holmes-nf-consistent}. В то же время непротиворечивость относительно арифметики Пеано показана Вильгельмом Аккерманом лишь для ZFC без аксиомы бесконечности \cite{ackermann-zfc-consistent}. Конечно, вследствие теоремы Курта Гёделя о неполноте доказательство непротиворечивости ZFC невозможно средствами самой ZFC \cite{godel-incompletness}. Однако даже относительная непротиворечивость ZFC доказана лишь в рамках расширений ZFC с недостижимыми кардиналами, что не приводит к каким-либо противоречиям в случае, если ZFC всё-таки противоречива \cite{kinamori-zfc-consistent}. Иными словами, аксиома бесконечности по какой-то причине препятствует непротиворечивости теории; если быть более точными - препятствует ей наша интерпретация данной аксиомы. На ошибку реификации, ответственную за проблемную интерпретацию, впервые указал Туральф Сколем в статье ''Некоторые заметки о аксиоматической теории множеств'' \cite{skolem-set-theory}, значительно опередившей своё время. В ней он демонстрирует относительную природу мощности и бесконечности, указывая, что аксиоматика теории множеств представляет собой не более чем набор правил формирования определённого функционального отношения - множества пар, которое мы называем отношением принадлежности $\in$. Это отношение, в свою очередь, задаёт алгебраическую структуру, аннотирующую объекты интуитивного представления - множества. Иными словами, мы пытаемся определить явный вид некоторого функционального отношения через правила ограничивающие его поведение. Но если, такая система правил однозначно (с точностью до изоморфизма) не определяет данное отношение, то есть допускает альтернативные объекты, которые мы можем интерпретировать как принадлежность элементов множеству, то классическая логика, требующая жесткого детерминизма, просто не применима в этой системе правил, без явного указания контекста фиксирующего конкретную альтернативу. Поэтому следующую аксиому, навеянную работами Сколема, мы назовём аксиомой бесконечности Сколема:
\begin{equation}
\begin{gathered}
\exists I \; [\emptyset \in I \wedge \forall x \; \exists y, z (\forall u \; (u \in y \Leftrightarrow \exists v  \; (u \in v \wedge v \in x) \wedge x \in z) \Rightarrow x \in I]
\label{axiom:skolem-infinity}
\end{gathered}
\end{equation}

Обозначая объединение всех элементов через $\bigcup x = y \Leftrightarrow \forall u \; (u \in y \Leftrightarrow \exists v  \; (u \in v \wedge v \in x))$ получаем более компактную форму::
\begin{equation}
\begin{gathered}
\exists I \; [\emptyset \in I \wedge \forall x \; (\bigcup x \in x \wedge \exists z \; (x \in z)) \Rightarrow x \in I]
\label{axiom:skolem-infinity-short}
\end{gathered}
\end{equation}

Легко видеть, что в отличие от стандартной аксиомы бесконечности данная форма стратифицируема, и в NF существует соответствующее множество $I = \{ x \mid x = \emptyset \; \vee \; (\bigcup x \in x \; \wedge \; \exists z \; (x \in z))\}$. Кроме того, можно показать, что теория $\mathrm{ZFS} = \mathrm{ZFC}^- + \mathrm{Sk}$ (ZFC без аксиомы бесконечности плюс аксиома бесконечности Сколема) непротиворечива относительно арифметики Пеано.

Теорема 7: {\sl ZFS непротиворечива относительно PA.}

Доказательство: 

$\mathrm{ZFC}^-$ имеет стандартную модель - множество наследственно конечных множеств $V_{\omega}$, которая интерпретирует PA $V_{\omega} \models PA$. Для любого конечного $n$ множество $V_n$ является моделью конечных фрагментов $\mathrm{ZFC}^-$. Однако во всяком $V_n$ аксиома бесконечности Сколема истинна, для чего мы позаимствовали у Джона фон Неймана требование существования множества содержащего данное множество. Пусть $\Phi$ - конечное подмножество аксиом ZFS. Тогда $\Phi$ содержит конечные фрагменты $\mathrm{ZFC}^-$ и Sk. Существует $n$ такое, что все формулы из $\Phi$ интерпретируются $V_n$, то есть $V_n \models \Phi$. Таким образом, по индукции заключаем, что любая конечная подсистема ZFS имеет модель. Отсюда по теореме о компактности \cite{godel-compactness-logic} следует, что ZFS непротиворечива. Теорема доказана.

Полученная теория весьма любопытна, поскольку предлагает альтернативную архитектуру ординалов, которые определяются стандартным образом: $\alpha - \text{ординал} \Leftrightarrow \forall \beta \in \alpha \; (\beta \subseteq \alpha \wedge \forall \gamma \in \beta \; (\gamma \subseteq \beta))$. Ординал $\alpha$ называется преемственным если $\exists \beta \; (\alpha = \beta \cup \{ \beta \})$. Для преемственных ординалов справедливо $\bigcup \alpha \subset \alpha$, а значит, все они являются элементами множества $I$, поскольку по аксиоме пары для каждого из них существует содержащий его синглтон. Следовательно, по схеме аксиом выделения существует множество всех преемственных ординалов $O_{regular}$, принадлежащих множеству аксиомы бесконечности. Это обеспечивает базовую индукцию и рекурсию и позволяет говорить о последовательностях, индексированных преемственными ординалами. Рекурсивно определим многократное взятие объединения: $\bigcup^{0} x = x$, $\bigcup^{n + 1} x = \bigcup (\bigcup^{n} x)$ и для всякого преемственного ординала, если $\exists \beta \in O_{regular} \; (\alpha = \beta + 1)$, то $\bigcup^{\alpha} = \bigcup (\bigcup^{\beta} x)$, и $\bigcup^{\alpha} = \bigcup x$ в противном случае. По схеме аксиом преобразования существует множество $trace(x) = \{ x, \bigcup x, \bigcup^2 x, \dots \}$, которое мы называем следом множества $x$. Это даёт элегантный способ разделить конечное и бесконечное внутри универсума ZFS в соответствии с концепцией уровней в работе Сколема. Корнем множества $x$ назовём пересечение всех элементов следа: $root(x) = \bigcap trace(x)$. Множество $x$ называется {\it конечным}, если $root(x) = \emptyset$ (множества первого уровня по Сколему), и  {\it бесконечным}, если $root(x) \neq \emptyset$. Существование $\dot{\omega}$ следует непосредственно из аксиомы выделения: $\dot{\omega} = \{ x \mid x \in I \wedge x - \text{ординал} \wedge root(x) = \emptyset \}$. Интересно, что в ZFS не возникает аналога теоремы ZFC $\bigcup \omega = \omega$. В ZFC доказательство следует из аксиомы экстенсиональности в силу того, что ординалы конечны тогда и только тогда, когда принадлежат $\bigcup \omega$, что в равной степени справедливо и для $\omega$. В ZFS индукция замкнута в множестве $I$: для всякого преемственного ординала следующий ординал принадлежит $I$. Однако если $\bigcup \omega = \omega$, то $\omega \notin I$ и индукция здесь не конструктивна; аксиома выделения выступает в роли внешнего преобразования. Тем не менее равенство $\bigcup \omega = \omega$ можно принять в качестве аксиомы, делающей $\omega$ неподвижной точкой объединения элементов и сохраняющей смысл обозначения $\omega$ для класса эквивалентности. То есть, мы вынуждены явным образом фиксировать, что ищем решение рефлексивное по отношению к внутренней индукции. При этом, поскольку $\bigcup \omega = \omega$, само множество конечных ординалов не принадлежит множеству, заданному аксиомой бесконечности Сколема: $\omega \notin I$. Тем не менее по аксиомам пары и объединения существует множество $\omega + 1= \omega \cup \{\omega\}$, для которого $\omega + 1\in I$, так как $\bigcup (\omega + 1) = \omega \in (\omega + 1)$. Это открывает возможность доопределения индукции и рекурсии до трансфинитных, позволяя задать инварианты внешних операций. Таким образом, мы получаем альтернативную стратификацию - разделение на предельные и преемственные ординалы. Преемственные ординалы содержатся в бесконечном множестве $I$, которое по сути является объединением областей регулярности разных уровней, в то время как предельные ординалы ему не принадлежат. Это напоминает стратификацию формул в ''Новых основаниях'' Куайна и разделение на множества и классы в ZFC и NBG. Данная аксиома вновь поднимает вопрос о непротиворечивости, который упирается непосредственно в интерпретации рефлексивности в теории. Таким образом, ZFS позволяет подняться на одну ступеньку выше, включив замкнутость отделённою от рефлексивности в непротиворечивое ядро. Поскольку класс ординалов формируется из двух частей - предельных и преемственных, такого объекта как множество всех ординалов не возникает: из множества $I$ с помощью аксиомы выделения можно получить только преемственные ординалы, а неограниченного выделение (усложнения) у нас нет. Благодаря этому парадокс Бурали-Форти \cite{burali-forti} в контексте аксиомы $\bigcup \omega = \omega$ избегается, и необходимости в дополнительной сущности - классах, как в $\mathrm{ZFC}$, - не возникает. 

В такой теории невозможность вывести рефлексивность бесконечных множеств (существование биекции со своей правильной частью) оказывается существенно более явной: предельные ординалы становятся точками разрыва предикативного описания или внешними замыканиями, в то время как алгебраическая структура относительно следования замкнута в $I$. Благодаря этому предельные ординалы возвращают себе смысл классов эквивалентности, аннотирующих части $I$, или неподвижных точек, где $\bigcup \omega = \omega$ равносильно $\omega - 1 = \omega$ в построениях Кантора. Если же рефлексивность не предполагается, можно зафиксировать это альтернативным предположением, то $\bigcup \dot{\omega} \in \dot{\omega}$. $\dot{\omega}$ существует как множество - объект теории, в силу схемы аксиом выделения, при этом этот объект принадлежит $I$, а значит единственная возможность состоит в том, что $\dot{\omega}$ - это первый объект с корнем отличным от пустого множества. В отсутствие рефлексивности необходимость в теореме Кантора не возникает: множество всех подмножеств содержит по синглтону для каждого элемента плюс пустое множество, делая биекцию невозможной. Если же принять аксиому рефлексивности для бесконечных множеств, теорема Кантора становится доказуемой в полученной системе. Любопытно также, что в ZFS воспроизводится парадокс Сколема \cite{skolem-set-theory}. Действительно, если смотреть изнутри на счётную модель универсума, можно наблюдать всю ординальную иерархию, тогда как с внешней точки зрения предельные ординалы являются просто конечными множествами. Данный подход оказывается весьма иллюстративным с онтологической точки зрения. Не важно внутри модели какой мощности мы находимся, сначала мы приходим к идее неограниченного продолжения, позволяющей придать смысл выражению: "и так далее до бесконечности" (индукция и рекурсия, где ''и так далее...'' может означать и $1000$ итераций, и $10000$, и $2^{1000}$), затем к идее актуальной бесконечности - неподвижной точки ($\omega$), которая оказывается изолированной от регулярных структур ($\notin I$), и в конце концов, выходим за пределы бесконечности в трансфинитную область. Концепция бесконечности как класса эквивалентности оказывается проще для понимания и кажется естественней, в то время, как концепция объективированной актуальной бесконечности $\dot{\omega}$ требует представлений о сингулярности, которую непросто предположить без знания о классах эквивалентности. Множество $\omega$ становится точкой прерывности индукции, отдельной от индуктивного построения, индукция заперта внутри множества $\omega$, но не конструктивна для самого $\omega$, где требуются дополнительные предположения. Это именно то, что происходит и в классических версиях ZFC и NBG, их модели интерпретируют арифметику Пеано PA, но не определяются ей. Совокупность всех натуральных чисел арифметики Пеано не является множеством (если конечно, мы явно не постулируем этого, что в корне меняет ситуацию), поэтому эту совокупность не получится экстенсионально сравнить с множеством $\omega$. Множество $\omega$ представляет собой нечто большее - объект другой теории, не принадлежащий области значений отношения следования, на него не распространяется внутренняя индукция по конечным объектам. Поэтому мы не можем сделать заключение, что $\bigcup \omega$ вообще является множеством, лишь дополнительные предположения, позволяют отнести такой объект к мнимым множествам - $\omega - 1$, или к множествам - $\omega$. Однако, в современных теориях мы делаем такие предположения неявно, распространяя внутреннюю индукцию на весь объект, в то время как бесконечное объединение множеств и бесконечное пересечение множеств должны рассматриваться отдельно, как имеющие отношение к точкам предикативного разрыва. К сожалению, нам придётся оставить эту примечательную теорию, так как в контексте данного исследования она призвана лишь продемонстрировать то, насколько серьезными оказываются последствия отсутствия критического анализа формы аксиомы бесконечности. Следует также сказать, что данные проблемы вовсе не критичны, они как мы выясним в дальнейшем вовсе не приводят к противоречивости теории или каким-то парадоксам, лишь затрудняют её анализ и препятствуют воссоединению с альтернативными взглядами в единую рамку, для чего окажется достаточным просто формально закрепить наши убеждения в аксиомах. 

Можно также взглянуть на проблему непротиворечивости с другой стороны, так Курт Гёдель показал, что логика первого порядка финитарно непротиворечива \cite{godel-compactness-logic}, что уже даёт нам мощный инструментарий для работы с множествами. Мы просто делаем ряд синтаксических дополнений. Всякий предикат $P(x)$ представляет домен $P$, где он истинен, то есть предикат и домен - это синонимы. Далее декларируем эквивалентность записи $P(x)$ и $x \in P$, то есть $P(x) \Leftrightarrow x \in P$ и эквивалентность операций $P(x) \wedge Q(x) \Leftrightarrow P \cap Q$, $P(x) \vee Q(x) \Leftrightarrow P \cup Q$. Тем самым получаем некоторую теорию множеств - доменов, которая при этом финитарно непротиворечива. В ней есть экстенсиональность $\forall x (P(x) \Leftrightarrow Q(x)) \Leftrightarrow \forall x (x \in P \Leftrightarrow x \in Q)$. Пара определяется предикатом $Pair_{a, b}(x) = (x = a \vee x = b)$ для любых заданных $a$ и $b$, пустое множество ложным предикатом, объединение и пересечение задано коньюнкцией и дизьюнкцией. Выделение - реализуется через пересечение, если задан предикат $P(x)$, его домен $P$, то каким бы ни был предикат $Q(x)$, выражение $P(x) \wedge Q(x)$ является предикатом с доменом $P \cap Q$. Если задан предикат $P(x)$ и его домен $P$, каким бы ни был предикат $Q(x, y)$, выражение $R(y) = \exists x (P(x) \wedge Q(x, y))$ является предикатом задающим домен $R$. Если при этом $\forall x \in P \exists! y Q(x, y)$, то имеем полный аналог схемы преобразования. Если сравнивать с ZFC, не хватает степени, регулярности, выбора и бесконечности. В финитарном контексте, если мы встраиваемся в доказательство Гёделя о полноте, то имеем дело с конечными фрагментами логики первого порядка и конечными моделями. Добавление аксиомы регулярности сохраняет непротиворечивость любой конечной модели в этом доказательстве, на любом шаге. Но ровно тоже самое справедливо и для выбора, если предикат не пустой, то существует элемент на котором он истинен, что и означает возможность выбора в финитарном контексте. Более того, аксиома бесконечности Сколема также встраивается без противоречий в контекст конечных моделей, что даёт нам финитарно непротиворечивую теорию множеств, в которой не хватает лишь аксиомы степени. В чистом виде, аксиому степени добавить нельзя, так как она несовместима с конечными моделями, тем не менее, можно заменить её на схему, например, $\forall x \exists y \forall z (z \in y \leftrightarrow z \subseteq x \wedge \phi(z))$, которая сохраняет финитарную непротиворечивость. Однако, с точки зрения современных представлений, данная схема не эквивалентна аксиоме степени, поскольку интерпретация последней подразумевает существование невыразимых множеств, для которых невозможно записать $\phi(z)$, при чём они возникают даже в счётной модели. Иными словами, мы верим в то, что при замене схемы аксиом на единственную аксиому степени в счётной модели, бесконечное множество $\omega$, которое является конечным, с внешней точки зрения, где описывается конечным множеством формул, происходит некоторая разблокировка части формул (причём большей). Однако, какой-либо финитарный механизм ответственный за такую блокировку отсутствует. Отметим, что выйдя за пределы современных убеждений финитарную непротиворечивость можно показать и для полной аксиомы степени, для чего можно воспользоваться идеей Кантора о перемежающихся последовательностях. Аксиома степени на каждом шаге увеличивает конечную модель до множества всех её подмножеств, непротиворечивость которых устанавливается на следующем шаге доказательства. Таким образом, действия аксиомы степени на конечные модели в финитарном контексте перемежаются с шагами доказательства, но перемежающиеся последовательности в соответствие с Кантором имеют один предел, которым является модель теории. 

Давайте воспользуемся ещё более простой и понятной конструкцией, чем новые основания Куайна, связанные с теорией простых типов Рассела, в конце концов, остаётся не очень понятным смысл, который несёт в себе стратификация, разделяющая формулы на два класса, может быть нестратифицируемость вовсе не свидетельствует о некоторой проблеме в формулах. Если мы внимательно взглянем на критику теории множеств со стороны интуиционистов и конструктивистов, то обнаружим, что основной претензией является вовсе не работа с бесконечными объектами как таковыми, особенно в контексте интуиционизма, а использование построений допускающих неодназначную интерпретацию, без достаточного обоснования предпочтения одной интерпретации над другой. Например, область определения предиката неприменимого к самому себе в парадоксе Рассела не оказывает никакого влияния на сам парадокс, даже если речь идёт о конечной области, парадокс всё равно возникает, что свидетельствует о том, что проблема заключается в самом предикате, а не в область его определения. Если быть более точными, то связь с бесконечностью всё же можно проследить через понятие предиката, поскольку концепции отображения или преобразования, наряду с концепцией множества, являются фундаментальными и очень тесно связанными между собой, так что идея бесконечности возникает из размышлений о любой из данных сущностей. То же самое можно сказать о парадоксе Рассела в теории множеств, наблюдаемая вселенная конечна, она состоит из фиксированного числа атомов (порядка $10^{80}$), конечно число связей, комбинаций и, следовательно, конечно число всех возможных множеств, включая все множества когда-либо сформулированные, так как ресурсы видимой вселенной ограничены, равно как и её информационная ёмкость (предел Бекенштейна). Среди полученной совокупности множеств есть как множества содержащие себя, например, $x = \{ x \}$ или множество всех идей, включающее себя как идею, так и множества не содержащие себя, например, множество всех атомов, не являясь атомом, оно не является собственным элементом. Но давайте рассмотрим, в этой конечной совокупности множеств, множество Рассела, состоящее из всех множеств, не содержащих себя: $R = \{ x \mid x \notin x \}$, которое так же конечно, как часть совокупности всех множеств. Легко видеть, что вопрос о принадлежности множества Рассела самому себе не имеет ответа $R \in R \Leftrightarrow R \notin R$, иными словами, для парадокса Рассела бесконечность вовсе не требуется. Нам достаточно рассмотреть множество всех когда-либо зафиксированных множеств (даже множеств зафиксированных в данном параграфе), которое безусловно конечно, и мы уже приходим к противоречию. Конечного множества удовлетворяющего формуле Рассела не может существовать, равно как и конечного множества удовлетворяющего аксиоме бесконечности, то есть обе формулы описывают нечто, невозможное в конечной реальности, в единственной доступной нам достоверной реальности, где критерий фальсифицируемости, позволяет считать некоторые предположения научными, отделяя их от убеждений и догм, пока они согласуются с наблюдениями (теории) или потенциально допускают проверку (гипотезы).

Однако, по какой-то причине, мы накладываем ограничения на множества - области определения, вместо того, чтобы как поступил Куайн, изолировать источник возникновения противоречий - формулы и исследовать механизмы связанные с возникновением парадоксов. Вооружившись данным наблюдением, не сложно исключить все конструкции, сомнительность которых связана с понятием пропозиции, из теории множеств. В самом деле, будем называть универсум {\it конечным} если он допускает потенциальное перечисление всех элементов, то есть доказательство того, что процесс перечисления элементов рано или поздно заканчивается, иными словами, конечный универсум представляет собой конечную коллекцию элементов - множеств: $V = \{ x_1, x_2, \dots, x_n \}$. Будем называть формулу $\varphi(x)$ {\it результативной} в универсуме $V$, если существуют множества $Z = \{ x \in V \mid \varphi(x) \}$ и $\neg Z = \{ x \in V \mid \neg\varphi(x) \}$ такие, что $Z \in V$ и $\neg Z \in V$, формально, мы пока не знаем является ли универсум множеством, поэтому запись $V \setminus Z$ была бы некорректной. Формула $\varphi(x)$ называется {\it конструктивной}, если существует конечный универсум $V$, в котором данная формула результативна, то есть определяет некоторое множество. Если существует универсум, в котором формула результативна для всех элементов, или не результативна для всех, то формула называется {\it условно конструктивной}. Данные простые определения дают нам возможность сформулировать схему неограниченного усложнения, вместе с аксиомой экстенсиональности составляющими конструктивную теории множеств:
\begin{align}
\begin{split}
\exists X \; \forall x \; &(x \in X \Leftrightarrow \check{\varphi}(x))\\
\end{split}
\end{align}

Где, $\check{\varphi}$ - конструктивная формула. Или, для каждого конструктивного свойства существует множество, состоящее в точности из тех элементов, что удовлетворяют этому свойству. В рамках современных представлений данная конструкция даёт нам инструментарий для работы с множествами, не позволяя использовать построения, которые в конечном контексте являются противоречивыми. Мы не понимаем причин, вследствие которых некоторые формулы способны порождать противоречия в определённых контекстах, из-за чего не способны ограничить аксиомами допустимые контексты, поэтому просто исключаем все такие формулы. В отличии от новых оснований Куайна, где нестратифицируемость формулы означает полную невозможность её использовать, в том числе в комбинациях с другими формулами, что можно трактовать как бессмысленность свойства выражаемого формулой в контексте универсума NF, в конструктивной теории множеств комбинирование неконструктивной формулы с другой способно сделать её конструктивной. Например, мы можем утверждать, что:

Теорема 8: {\sl Если $\check{\varphi}(x)$ конструктивна и $\check{\varphi}(x) \Rightarrow \psi(x)$, то $\check{\varphi}(x) \wedge \psi(x)$ конструктивна.}

Доказательство: 
Для формулы $\check{\varphi}(x)$ существует конечный универсум $V$, где она определяет конечное множество, в котором формула $\psi(x)$ выполнена для всех элементов и не задаёт никакой дополнительной структуры, таким образом $\check{\varphi}(x) \wedge \psi(x)$ результативна в данном универсуме. Теорема доказана.

Обратим внимание на весьма любопытное следствие, в случае $\check{\varphi}(x) = \bot$ любая формула конструктивна, иными словами, в пустом множестве никакие формулы не приводят к противоречию. Данное наблюдение помогает нам понять источник проблем возникающих в современных теориях множеств. Единственным действительно рефлексивным объектом является пустота. Только пустое множество можно отобразить на любую свою часть. Мы вводим пустое множество основываясь на интуитивном представлении, вместо того чтобы зафиксировать его через свойство рефлексивности, чтобы не иметь соблазна распространить данное свойство на другие объекты. Бесконечность в отличии от пустоты, обладает структурой целого подобной структуре части. Мы используем импредикативность - исключаем часть информации о целостной конструкции, выделяя его внутреннюю структуру, чтобы добиться подобия, которое всегда должно быть поддержано внешними операциями сохраняющими целостность объекта. В противном случае мы можем ожидать, что отсутствие информации проявит себя через рефлексивность точно также как $0$ без осознания его природы проявляет себя в арифметике, позволяя доказать $1 = n \forall n \in \mathbb{N}$. Это именно то, что происходит в ZFC и NBG, мы теряем внешнюю поддержку, которая была у Кантора, и пытаемся обойтись единой операцией сложения, вследствие чего считаем объекты теории действительно рефлексивными. Отметим также, что в конструктивной теории множеств универсум не является множеством, поскольку формула $\check{\varphi}(x) = \top$ неконструктивна. По сути данная теорема позволяет разблокировать неконструктивные формулы в рамках уже заданных множеств. В частности, можно определить множество регулярных множеств $\{ x \mid x = \emptyset \vee \exists y \; (y \in x \wedge y \cap x = \emptyset) \}$, его конечный универсум $V = \{ \emptyset, x = \{ x, \emptyset \} \}$. Если мы требуем данного свойства от всего универсума, любые самореферентные формулы становятся конструктивными в таком универсуме, так как они эквивалентны константе истине или лжи. По сути именно, это мы делаем в рамках ZFC и NBG, мы стремимся избежать противоречий вызванных самореференцией. Тривиально также убедится в справедливости теорем соответствующих аксиомам пустого множества, пары, объединения и степени. Более того, аксиома бесконечности Сколема, обеспечивает существование бесконечных множеств в такой конструктивной теории. К особенностям данной теории можно отнести, то что стратификация не блокирует часть формулы, а позволяет исправить её с помощью дополнительных условий. Например, формула Рассела $x \notin x$ неконструктивна, но формула $x \notin x \wedge \exists y \; (x \in y)$ уже конструктивна. 

Такая простенькая конструкция делает некоторые вещи невероятно наглядными, помогая разобраться в механизмах ответственных за возникновение парадоксов. Так если мы пытаемся сконструировать конечный универсум удовлетворяющий формуле Рассела, то сталкиваемся с проблемой, нам требуется множество которое бы содержало все множества не содержащие себя, но если такое множество есть в универсуме и оно не содержит себя, то оно должно быть в нашей коллекции, а если оно содержит себя, то оно не подходит под формулу. Даже в универсуме, где все множества регулярны, формула Рассела определяет сам универсум, из-за чего не может быть конструктивной. Таким образом, парадокс Рассела возникает вовсе не из-за само-принадлежности данного множества, а из-за того, что определяемое формулой множество заведомо отсутствует в универсуме. Иными словами, формула требует выхода за пределы универсума, описывая некоторую сущность вне его, вопрос о само-принадлежности данного множества не имеет смысла, поскольку оно не принадлежит универсуму. Это именно то, что делает Куайн в новых основаниях, он лишает сам вопрос о само-принадлежности смысла, его невозможно задать в рамках NF, вопрос о принадлежности универсума NF самому себе невозможно сформулировать в рамках теории, а тем более получить на него какой-либо ответ. Поэтому, мы используем финитарный фиксатор $\exists y \; (x \in y)$, позаимствованный в исходных версиях NBG Джона фон Неймана, в аксиоме бесконечности Сколема, в противном случае в ZFS возникают проблемы, когда универсум становится множеством, но в рамках ZFS, где справедлива теорема Кантора, он не может быть тождественен своему множеству степени. Многие парадоксы теории множеств связаны именно с таким поведением формул, например, парадокс Рассела \cite{russell-paradox}, парадокс Бурали-Форти \cite{burali-forti} и другие, мы называем подобные формулы {\it внешними}, поскольку они требует выхода за границы имеющихся абстракций, в свою очередь формулы замкнутые в универсуме называются {\it внутренними}. Несложно видеть, что одной из таких проблемных внешних формул является аксиома бесконечности, тем не менее, мы помещаем её в самое сердце оснований математики, совершенно не удивительно, что наши теории начинают выходить за свои рамки, требуя существования недоказуемых истин и объектов невыразимых в них.

Итак, аксиома бесконечности, не является атомарной (в общепринятой интерпретации влечёт отсутствие максимальных элементов), не является чистой (требует присутствия бесконечного замыкания в универсуме как объекта и задаёт его структуру), не является однозначное (неясны свойства сингулярности, разделяющей области регулярности), и имеет проблемы с независимостью (позволяет доказать рефлексивность в неочевидном сочетании с другими аксиомами). Но по задумке она должна лишь постулировать свойство замкнутости по отношению следования. Она не стратифицирована и неконструктивна, то есть потенциально может приводить к противоречиям в некоторых контекстах, а в её основе лежат именно те принципы, которые приводили к парадоксам в логике Готлоба Фреге и наивной теории множеств. Фундаментальным ограничением логики первого порядка стало ограничение само-референции, в том числе опосредованной, ведь, само-применимость не была явной ни в теории Фреге ни в наивной теории множеств, она реализовывалась через вспомогательные объекты - области значений концептов и множества элементов, на которых предикат истинен \cite{frege-logic}. Но именно это мы делаем в аксиоме бесконечности, мы применяем формулу к самой себе опосредованно через множества и универсум, мы квантифицируем по домену теории, в рамках которого определяем объект, который доступен для квантификации. Если рассмотреть аксиому бесконечности саму по себе, вне контекста современных аксиоматик, то квантификация $\forall x$ происходит в том числе и по определяемому объекту $I$ и ничто не препятствует $x = I$, в таком случае бесконечное множество становится неподвижной точкой отношения следования $I \in I \Rightarrow I \cup \{ I \} \in I$. Это не обязательно приводит к противоречию, однако, является именно зацикливанием. По какой-то причине, мы стремимся избежать синтаксиса неподвижных точек в теории и принимаем неявные соглашения, помогающие замаскировать их, которые пытаемся закрепить аксиомами цензуры. Мы выносим явную форму такой рефлексии за пределы описания теории, мы избавляемся от случая $x = I$, но стремимся запретить также возможность того, что бесконечности для возникновения требуются вспомогательные объекты, которые из-за этого порочного круга не могут обрести смысл $\omega - 1 = \omega$ естественный для классов эквивалентности - порядковых типов. Более того, мы стремимся закрепить в теории именно такую интерпретацию порочного круга и само-референции, в той форме, с которой мы так страстно боролись, через универсальность свойства рефлексивности. И это при том, что аксиома бесконечности в принципе является фундаментально не верифицируемой, указывая на то, что с ней нужно быть особенно осторожными и внимательными. Это уже можно охарактеризовать как грубую формальную ошибку в аксиоматиках теории множеств. Однако из-за всех этих хитросплетений мы упустили ключевой смысл заложенный в аксиому бесконечности: она является единичным свидетельством, а не универсальным правилом. Мы совершаем одну и ту же формальную логическую ошибку преждевременного обобщения дважды: один раз, принимая определённый нами класс множества пар за единственную и универсальную возможность реализации функциональных отношений между объектами, несмотря на прямо противоположное мнение высказанное Кантором и рудименты исходного разделяния непосредственно в нашей теории. Во второй раз, мы совершаем более специфичную разновидность - формальную ошибку ложной дихотомии, полагая будто засвидетельствованная в аксиоме бесконечности индуктивная структура - единственная возможная структура бесконечных множеств в теории, однозначно разделяющая объекты на два типа: конечные и индуктивно бесконечные. Но аксиома бесконечности не утверждает, что всякое бесконечное множество индуктивно, она утверждает лишь, что в теории среди прочих альтернатив гарантированно присутствует описание бесконечности через замкнутую индуктивную структуру. Это возвращает нас к открытию Сколема и вопросу однозначности отношения принадлежности. Аксиома бесконечности говорит, что среди классов пар задающих функциональное отношение принадлежности есть подкласс, в котором участвуют только объекты индуктивных структур. Но она никоим образом не запрещает тому же отношению принадлежности выражать в теории альтернативные морфологиями бесконечности, связанные с классами пар содержащие объекты не участвующие в индуктивной структуре, в частности, те что мы называли мнимыми. Это связано с тем, что в самой по себе FOL мы можем гарантировать изоморфность синтаксиса и семантики, благодаря полноте. Однако, FOL + ZFC(NBG) не обладают данным свойством, здесь изоморфизм не доказан и, по всей вероятности, недоказуем вследствие неполноты. Это означает, что приняв все аксиомы логики и все аксиомы ZFC в строго формальном смысле мы не можем обосновать все сделанные на сегодняшний день выводы, нам требуется также явно связать синтаксис и семантику, в противном случае, наш синтаксис может выражать параллельные семантические значения. Теория может быть синтаксически противоречивой, но семантически непротиворечивой, имея модель, которая разделяет семантические значения в модельном синтаксисе. Следует также сказать, что данные формальные ошибки являются не ошибками самой теории, а методологическими ошибками некорректного мета-позиционирования теории, если принять мета-установку однозначной согласованности синтаксиса и семантики, практически все выводы остаются идеально логически обоснованными. В то же самое время, такого рода методологические ошибки существенно затрудняют анализ консистентности теорий. Так или иначе, формула аксиомы бесконечности - это обычная формула, она просто оказывается крайне неудачной для того контекста, в котором мы её используем. Она приглашает нас к размышлениям, к анализу того, что может означать выход за пределы конечности, она не придаёт смысла понятию бесконечность, она требует, чтобы мы его предоставили. Неоднозначность синтаксиса теории может приводить к нарушению закона исключённого третьего (множество может оказаться конечным, индуктивно бесконечным и бесконечным в некотором альтернативном смысле), при этом если модель остаётся консистентной, различая типы функциональных отношений и структур бесконечных множеств, то теория сохраняет непротиворечивость. Архитектурные недостатки аксиомы бесконечности, действительно чрезвычайно сложно обнаружить, необходимо иметь представление об объекте поиска, авторам пришлось построить альтернативную теорию множеств, чтобы понять причины парадоксальности выводов в основаниях математики, благодаря чему удалось обнаружить ряд фундаментальных математических принципов, в частности, независимость рефлексивности бесконечных множеств от современных аксиоматик, о чём авторы надеются поведать в следующих работах.

В современных теориях множеств присутствует и другой недостаток, менее критический и относящаяся скорее к философии чем к математике - это ошибка реификации, выражающаяся в наделении алгебраической структуры объектностью. Тем не менее, данная онтологическая ошибка существенно затрудняла поиск проблем в теории, предлагая для исследований абстрактную иллюзию вместо фрагмента реальности. Мы стали воспринимать аннотацию через классы эквивалентности множеств - порядковые типы, как сами множества, несмотря на критику Сколема, Витгетштейна, Брауэра и других. Данная ошибка была растянута во времени и её причиной стала целая совокупность факторов, от вытеснения философии прикладными дисциплинами в образовании в процессе реформирования, до отсутствия практики формулировать систему методологических предположений. Исходно аксиоматики теории множеств работали с произвольными объектами, а не с самими множествами, где порядковые типы были классами эквивалентности множеств, сохраняя онтологическую чистоту Кантора. Но благодаря работе Джона фон Неймана оказалось, что произвольные объекты можно моделировать ординалами - множествами внутри самой теории. И хотя сам автор замечает, что в строгом смысле корректно говорить об ординалах как о типах множеств (особенно когда мы записываем формулы бесконечных ординалов), и лишь в контексте уже заданной трансфинитной индукции допустимо сопоставить порядковым типам некоторые ординалы в качестве канонических представителей, наглядно моделирующие иерархию порядковых типов \cite{von-neumann-ordinals}, его модель становится центральной в контексте аксиоматических теорий. А то, что $\omega = \{ 0, \{ 0 \}, \{ 0, \{ 0 \}\}, \dots \}$ - это тип бесконечных множеств, со структурой отраженной в формуле, который можно считать множеством, только если уже определён соответствующий порядковый тип $\omega$ (в частности, в терминах самих множеств), очень быстро забывается. Принимая во внимание всё сказанное выше, давайте попробуем временно дисквалифицировать неоднозначные интерпретации теорий множеств и вернутся к оригинальным построениям Кантора, уделявшего существенное внимание онтологической и эпистемологической согласованности, обеспечивающей простоту анализа и развития. Нам необходимо понять, где именно логически безупречные математические структуры Кантора начинают расходится со структурами современных теорий.

Для начала давайте снова обратимся к новым основаниям Куайна. В NF универсум $V = \{ x \mid x = x \}$ является множеством, при этом теорема Кантора $|X| < |\mathcal{P}(X)|$ недоказуема, так как построение диагонального множества нестратифицировано, в качестве примера множества нарушающего теорему Кантора, можно привести сам универсум $V$ поскольку $V = \mathcal{P}(V)$ и тождественное преобразование обеспечивает биекцию между ними \cite{rosser-nf} (несмотря на то, что не для всех комбинаций $V$ и $\mathcal{P}(V)$ стратификация позволяет построить такое преобразование, факт его существования и тождества $V = \mathcal{P}(V)$ делают его доступным для всех). В этой теории, как мы уже отмечали, понятие само-принадлежности не определено, мы просто исключили из нашего аппарата возможность выражать данное свойство, нельзя утверждать, что универсум NF $V$ содержит сам себя или не содержит сам себя, как и любое другое множество, все подобные вопросы являются не более чем спекуляциями внешнего наблюдателя. В этом состоит существенная трудность работы с NF, так как нам сложно заставить себя осознать, что некоторые высказывания не просто не имеют смысла, а не существуют. Например, Эрнст Спекер показал, что полная аксиома выбора несовместима с NF \cite{specker}, что является отражением того факта, что в NF существуют неупорядоченные множества, то есть множества для которых порядок невозможно задать или упорядоченная версия является другим множеством. Необходимо сделать уточнение, из аксиома выбора выводится центральное для современных аксиоматик следствие, что для всякого множества существует описание в терминах алгебраической структуры вполне упорядоченных множеств, но NF определяет свою алгебру отличную от ZFC и NBG, в которую, таким образом, невозможно вложить структуру вполне упорядоченных множеств, однако, в то же самое время, всякая алгебра претендующая на полноту описания имеет некоторый аналог данного утверждения, позволяющий описывать необходимые множества как объекты упорядоченной алгебры. Всё дело в том, что условие стратификации, унаследованное из теории типов, - существенно более строгая конструкция в отношении ликвидации самореференции. Как мы видели в аксиоме бесконечности, логика первого порядка позволяет обойти такое ограничение опосредовано через квантификацию по всем объектам универсума включая те, что возникают в аксиоме. В строгом смысле это означает, что сам универсум теорий использующих логику первого порядка может быть импредикативным, то есть не допускающем описания в терминах статичного существования в платоновском смысле, что является необходимым компромиссом для того, чтобы в теории вообще мог возникнуть такой объект как континуум. Но в NF несмотря на квантификацию по всему универсуму, часть формул оказывается нестратифицированными, означая отсутствие выразимости некоторых свойств в теории. Например, по схеме неограниченного усложнения $\forall x (x \in V \Leftrightarrow x = x)$, но в случае $x = V$ формула $V = V$ становится нестратифицированной в силу того, что левая часть $V \in V$ требует различия типов, которого формула $V = V$ не в состоянии удовлетворить. На семантическом уровне данный факт можно интерпретировать как ограничение квантификации, словно мы можем квантифицировать не по всему чему хотим, а только по тому, что выразимо в теории и не приводит к противоречию, обуславливая невозможность задать полный порядок на любом множестве. Иными словами, полный порядок в таких теориях является импредикативным и не имеет синтаксического выражения. В современной интерпретации теории множеств, аксиома выбора понимается как аксиома воспроизводимого или предикативного выбора. Которая всякий раз дублирует один и тот же результат, таким образом обеспечивается ещё одно ключевое свойство необходимое для доказательства теоремы Кантора - свойство тотальности. Доказательство Спекера означает существование бесконечных множеств в NF, поскольку для любого конечного множества аксиома выбора выполняется. Например, аксиома бесконечности Сколема: $I = \{ x \mid \bigcup x \in x \}$ задаёт такое множество. Если мы выходим на уровень метатеории, то бесконечность можно показать существенно проще, воспользовавшись аргументом Евклида, если в NF отсутствует бесконечность, то $I$ - конечно, но на уровне мета-теории (достаточно конструктивной теории) мы можем выбрать из $I$ все конечные ординалы, взять максимальный $m$ и определить $m + 1 = m \cup \{ m \}$, после чего проверить, что новое множество удовлетворяет формуле, получая противоречие с тем, что $\dot{\omega}$ содержит все множества описываемые формулой. Однако, для наших целей первостепенное значение имеет следующее. Если NF непротиворечива, а у нас есть все основания предполагать, что так оно и есть \cite{holmes-nf-consistent}, то по теореме Лёвенгейма-Сколема она имеет счётные модели \cite{skolem-set-theory}. Ключевую роль для нас играет то, что даже на счётном множестве можно определить алгебраическую структуру (аксиоматикой NF) позволяющую отождествить (взаимно однозначно отобразить) множество и множество его подмножеств. В частности, у нас нет причин полагать, что универсумом NF не является, например, множество $P$, построенное выше и связанное с иерархией фон Неймана, с тем лишь отличием, что в NF не возможно выразить его иерархическую структуру полностью как в ZFC и NBG. Это позволяет ему быть тождественным множеству своих подмножеств, а стратификация блокирует парадокс Рассела, который воспроизводился как диагональная конструкция включающая все элементы. Здесь лучше всего видно как работает аксиоматическая цензура, мы просто вырезаем из теории синтаксические конструкции, которые приводили к противоречию, точно также как мы делали в ZFC и NBG. Так или иначе, мы снова приходим к простому наблюдению, парадоксы теории множеств не связаны с размером множеств, они связаны с их алгебраической структурой, понятия мощности и бесконечности относительны как отметил Сколем, а значит слишком большое множество - относительное наблюдение. Обратим внимание, на то, что относительность понятий установленная Сколемом и подобие алгебраических структур исследуемых Кантором и Дедекиндом тесно взаимосвязаны между собой, классы эквивалентности описывающие бесконечность на разных уровнях имеют подобную алгебраическую структуру, что и позволяет говорить об их относительности.

Принимая во внимание ошибку реификации мы можем построить теорию множеств которая комбинирует аксиоматики ZFC или NBG с аксиоматикой NF, всё что для этого требуется это ввести дополнительный символ $\dot{\in}$ и соответствующее отношение, тогда ZFC и NGB будут задавать алгебраическую структуру аннотирующую множества связанную с отношение принадлежности $\in$, в то время как NF будет задавать алгебраическую структуру аннотирующие те же самые множества, но связанную с отношением {\it включения} $\dot{\in}$. Для этого нам не требуется усиление логики, мы просто вводим в теорию новый символ представляющий альтернативное функциональное отношение включения, получая тем самым консервативное расширение, при чём в обе стороны, либо NF расширяет ZFC (NBG), либо наоборот. Это позволяет нам явно проиллюстрировать логическую ошибку преждевременного обобщения стандартной интерпретации, теоремы связанные с аксиомой бесконечности не являются универсальными теоремами теории, они справедливы только в индуктивном фрагменте теории связанным с данной аксиомой, в то время как NF высвечивает альтернативные морфологии никак не конфликтующие с индуктивной. Но именно такой была онтологическая картина Кантора, он работал непосредственно с объектами нашего мышления, а не с абстрактными алгебраическими объектами некоторого формализма. Кантор конструирует алгебры кардинальных и порядковых чисел для аннотации более базовых объектов - множеств. Именно поэтому он вводит различную терминология для типов эквивалентности и типов подобия, и первая предшествует второй, но мы полностью игнорируем это в современных интерпретациях ZFC и NBG. Здесь кардиналы - это просто вид ординалов, не более того, они не связаны со своей алгебраической структурой, они являются частью алгебраической структуры ординалов. Однако, сам Кантор категорично отстаивал убеждение, что кардинальная и ординальная алгебры - это разные системы описания не сводимые друг к другу. Причина путаницы заключается в том, что мы определяем мощность через неупорядоченные биекции в упорядоченном универсуме. Но в упорядоченном универсуме любая биекция упорядочена, она является частью строгой кумулятивной иерархии, где даже произвольный порядок требует структуры подтверждающей его произвольность, нам нужна уверенность, что абсолютно все элементы связаны данным отношением, и как следствие любая биекция индуцирует определённый порядок на областях определения и значений. Например, биекция между множество натуральных чисел и множеством рациональных чисел, индуцирует всюду плотный порядок на множестве натуральных чисел. Биекция - это часть алгебраической структуры, в алгебрах ZFC и NBG не существует биекции между множеством и его множеством степени, в то время как в алгебре NF и других системах, например, позитивной теории множеств она есть. У Кантора было две алгебраических структуры, каждая со своим классом функциональных отношений, включающим взаимно однозначное соответствие. В кардинальной алгебре соответствие не подразумевало никакого порядка, оно было классом эквивалентности множеств пар элементов неупорядоченных множеств. В ординальной алгебре соответствие подразумевало сохранение порядка, точно также являясь классом эквивалентности множеств пар, но уже упорядоченных множеств. Таким образом, наше определение мощности в рамках общепринятых аксиоматических теории становится просто свойством порядка - мощность - это отображаемые друг на друга порядки. Некоторые множества пар нарушающие структуру порядка универсума становятся классами за его пределами, например, взаимно однозначное соответствие множества натуральных чисел в естественном порядке (изоморфного $\omega$) с множеством натуральных чисел упорядоченных по чётности и по возрастанию (изоморфного $2\omega$). То, что порядок вполне упорядоченного множества нельзя отобразить на порядок его множества степени - это просто свойство нашей ординальной алгебраической структуры, которой мы описываем множества (биекция возможная в NF просто является классом и вынесена из описания). К сожалению, даже несмотря на то, что Сколем в своей критике развеял миф о сакральности взаимно однозначного соответствия, показав, что понятие биекции относительно и замкнуто в алгебраической структуре, то есть в системе определённых нами правил и операций, даже в настоящее время можно встретить работы рассматривающие взаимно однозначное соответствие в универсальном контексте. В контексте современных теорий, мощность более точно по смыслу можно определить как существование произвольно упорядоченной биекции. Стандартное определение касается именно неупорядоченной биекции, не описываемой никаким порядком, к которой, в том числе, относится класс пар взаимно однозначного соответствия универсума NF, не являющийся множеством в ZFC и NBG. Это ровно те ограничения доказательств несчётности вещественных чисел (теорем Кантора), которые мы наблюдали в предыдущих статьях серии. На данном историческом этапе понятие мощности закреплено за интерпретациями ZFC и NBG, поэтому мы вернём в математику классическое понятие сформулированное Кантором, но назовём его {\it ёмкостью}, чтобы избежать путаницы. Таким образом, мощность подразумевает произвольный порядок, а ёмкость отсутствие каких-либо требований к порядку. В порядковой алгебре мощность множества строго меньше мощности его множества степени, в то время как в кардинальной алгебре, теорема Кантора недоказуема и емкости этих множеств равны, благодаря биекции в тотально неупорядоченных системах.

При чтении фундаментальной работы Георга Кантор "К обоснованию учения о трансфинитных множествах", невозможно не заметить его исключительную осторожность: автор избегает любых утверждений, в истинности которых он не уверен. В этом объёмном и скрупулёзно выстроенном труде, который принято считать кульминацией его теории множеств, Кантор ни разу не формулирует неравенство $\aleph_0 < 2^{\aleph_0}$, более того, это неравенство отсутствует и в его последующей переписке. Кантор возвращается к своему результату 1878 года о том, что $\mathbb{R}^n$ и $\mathbb{R}$ имеют одну и ту же мощность, и показывает, как этот факт теперь элегантно получается в рамках арифметики кардинальных чисел, для чего он использует обозначение мощности континуум $2^{\aleph_0}$ - это прямое следствие определения степени множества как покрытие одного множества другим. Однако, даже в этом контексте Кантор сознательно не упоминает ни диагональный метод, ни само неравенство $\aleph_0 < 2^{\aleph_0}$ \cite{cantor-continuum-cardinality}, равно как и в своём исследовании порядкового типа линейного континуума \cite{cantor-continuum-order}. Единственное известное место, где Кантор применяет диагональный аргумент именно к кардинальной арифметике и, по-видимому, склоняется к тому, что $2^{\aleph_0} > \aleph_0$, - это письмо Рихарду Дедекинду от 31 августа 1899 года из санатория в Ханенклее \cite{cantor-strange-letter}. Однако этот период совпадает с одним из тяжёлых срывов Кантора \cite{dauben-cantor-1899, grattan-guinness-1899}, что не позволяет с уверенностью утверждать, что данное письмо выражает его окончательную позицию и что он действительно считал возможным объединять кардинальную и ординальную арифметику в единую алгебраическую структуру. Дело в том, что понятие мощности и трансфинитного числа постоянно эволюционировали в творчестве Кантора, но только в работе "К обоснованию учения о трансфинитных множествах" они впервые получают свою окончательную форму. Именно здесь мощность перестаёт быть просто признаком для сравнения двух множеств и становится полноценным объектом арифметики - кардинальным числом. Кантор предлагает структурное построение алгебры, через систему операций, в первую очередь дизъюнктного объединения, воплощающуюся в сложении кардинальных чисел, из которой комбинаторно получается декартово произведение - умножение, и покрытие - возведение в степень. Эти операции абстрагированы от конкретной природы элементов и от любого порядка их расположения. Два множества имеют равные кардинальные числа, если между ними существует взаимно однозначное соответствие, независимо от какого-либо упорядочения \cite{cantor-cardinal}. 

Только после этого Кантор переходит к введению порядка на множествах и особое внимание уделяет вполне упорядоченным множествам, для которых вводит понятие порядкового типа. Порядковые числа образуют самостоятельную алгебраическую структуру с собственной операцией сложения - конкатенацией вполне упорядоченных отрезков, из которой выводятся сохраняющее порядок декартово произведение - умножение и возведением в степень. Два вполне упорядоченных множества имеют один и тот же порядковый тип, если они изоморфны как упорядоченные множества, то есть существует сохраняющий порядок взаимно однозначное соответствие \cite{cantor-ordinal-algebra}. Операции сложения и умножения двойственны, левая операция рефлексивна и связана с точкой разрыва предикативности, в то время как правая - нерефлкесивна и обеспечивает сквозную индукцию. Порядковая арифметика позволяет ввести понятие последовательности и связанные с ней понятия предела и ряда, которые строго связаны с её структурной доминантой и правыми операциями сложения и умножения. Вследствие чего в порядковых типах существует единственная операция возведения в степень выражаемая через нерефлексивное правое умножение. Кантор пишет, что понятие порядкового типа перенесённое на ''кратно упорядоченные множества'' и понятие мощности охватывают всё числоподобное, которое мыслимо вообще, однако, он несколько преувеличивает универсальность такого подхода. В 20 веке Джон Конвей, развивая аналогичный структурный принцип (но на основе рекурсивного задания пар множеств левых и правых элементов $\{L \mid R\}$), построил класс сюрреальных чисел $No$. Этот класс содержит все порядковые числа Кантора, все вещественные числа, а также всевозможные инфинитезимальные и бесконечно большие величины. В отличии от Кантора, у Конвея равенство и порядок определяются первично и чисто арифметически, а изоморфизм в глобальном линейном порядке оказывается уже теоремой (что подчёркивает связь взаимно однозначного соответствия с алгебраической структурой). Ключевая конструктивная операция у Конвея - не линейная конкатенация, а более сложное слияние древовидных структур рождения \cite{conway-surreal-numbers}. Поэтому, хотя любой сюрреальный порядок индуцирует обычное отношение порядка, обратное восстановление полной древовидной структуры рождения только по линейному порядку невозможно. Таким образом, класс всех мыслимых числоподобных объектов шире, чем предполагал Кантор, и требует более богатой базовой информации, чем просто отношение порядка на множестве.

По-видимому, Кантор к этому моменту уже осознаёт ограниченность всех своих прежних доказательств несчётности континуума, каждое из них - будь то метод вложенных отрезков или диагональный аргумент - опиралось на конкретную структуру порядка, которой в чистой теории мощностей быть не должно (что мы отмечали в предыдущих статьях серии). В ключевой работе Кантор переопределяет мощность, кардинал - это первичный, структурно-нейтральный объект, определяемый исключительно абстрактными операциями. Никакого внутреннего порядка, никакой иерархии начальных ординалов, никакой привилегии хорошо упорядочиваемых множеств в этой алгебре нет. Взаимно однозначное соответствие - это просто факт существования отображения, лишённый любой дополнительной морфологии. В такой алгебре диагональный аргумент в его классической форме невозможен: он требует дополнительной структуры, которая в чистой кардинальной арифметике просто невыразима, что оставляет открытым вопрос о соотношении между $\mathfrak{a}$ и $2^{\mathfrak{a}}$, на который Кантор не даёт никакого ответа. Он вводит возведение в степень $\mathfrak{a}^{\mathfrak{b}}$, доказывает алгебраические равенства вроде $\mathfrak{a}^{\mathfrak{b}}*\mathfrak{a}^{\mathfrak{c}} = \mathfrak{a}^{\mathfrak{b} + \mathfrak{c}}$ или $\mathfrak{a}^{\mathfrak{c}}*\mathfrak{b}^{\mathfrak{c}} = (\mathfrak{a}*\mathfrak{b})^{\mathfrak{c}}$, но нигде не приводит $\mathfrak{a} \neq 2^{\mathfrak{a}}$ как теорему кардинальной арифметики. В предыдущих статьях серии мы неоднократно наблюдали как предположения об гомогенности или единственности бесконечности, об отсутствии её форм и структуры, блокируют возможность построений доказательств несчётности, в силу несовместимости их условий, что именно предположение о структуре бесконечности обеспечивает частичную консистентность достаточную для доказательства. Но в кардинальной арифметике Кантора $1 + \aleph = \aleph + 1 = \aleph$, $\aleph + \aleph = \aleph$, $\aleph*\aleph = \aleph$. Здесь не существует никакой аддитивно-мультипликативной структуры, на которую могли бы опереться построения этих доказательств, все комбинаторные асимптотики схлопываются в один объект. Только спустя сорок лет, Куайн в своих Новых основаниях показал, что алгебраическая картина, очень близкая к той, что построил Кантор, допускает последовательную аксиоматизацию, в которой для универсума $V$ действительно существует биекция $V \leftrightarrow \mathcal{P}(V)$. Тем самым Куайн реализовал одну из возможностей, которые Кантор оставил открытыми, отказавшись навязывать своей кардинальной алгебре дополнительные структурные ограничения. Современные аксиоматики теории множеств, подчиняют кардинальные числа ординальной иерархии и принудительно сливаю две независимые алгебры, являясь тем самым, лишь одной частной, жёсткой, теории упорядоченных множеств Кантора. Именно поэтому нам пришлось временно отказаться от них. Исходная теория Кантора оказывается достаточно гибкой, чтобы служить общей платформой для объединения трёх числоподобных структур - собственно кардинальных чисел, ординальных чисел и нестандартных натуральных чисел, сюрреальных чисел или гросс-чисел - без принудительного сведения одних к другим. Именно в этой нейтральности и заключается подлинная глубина и современность позднего Кантора.

Итак, в контексте фундаментальной работы Кантора по теории множеств, множество степени счётного множества оказывается счётным, то есть имеет ёмкость просто бесконечной последовательности. Данный вывод поддержанный Новыми основаниями Куайна и теорией Кантора, проливает свет на то почему так много работ, которые незаслуженно считают маргинальными опровергают современную интерпретацию о несчётности континуума. Поскольку мы упразднили современный статус кардиналов связанный с некоторыми порядковыми числами, необходимо немного погрузится в алгебру этих объектов, чтобы сохранить возможность описывать ключевые структуры. Давайте немного углубимся в онтологию Кантора, чтобы понять как порядковые числа функционировали, до внедрения аксиоматик. Совокупность всех конечных кардинальных чисел $\{ \mathfrak{n} \}$ мы определяем как трансфинитное множество, которое описывается первым трансфинитным кардинальным числом $\aleph_0$. Нам не известно каким именно образом реализуется бесконечное множество всех конечных кардинальных чисел, мы не можем утверждать, что существует только одно такое множество. Однако, описывая его кардинальным числом, мы фиксируем класс эквивалентности таких множеств, которые отображаются друг на друга каким-либо взаимно однозначным соответствием без порядка и содержат при этом все конечные кардинальные числа. Если у нас есть аксиома экстенсиональности, то мы можем доказать, что таких множеств существует много, более того, Кантор сам даёт такое доказательство рассматривая множество $(\{ \mathfrak{n} \}, \mathfrak{e}_0) = \{ \mathfrak{n} \mid \mathfrak{n} - \text{конечный кардинал} \} \cup \{ \mathfrak{e}_0 \}$, оно удовлетворяет нашим требованиям: содержит все конечные кардинальные числа и описывается $\aleph_0$, таким образом, принадлежит обозначенному классу эквивалентности и не равно исходному $\{ \mathfrak{n} \}$ \cite{cantor-aleph-zero}. Существенно то, что в рамках такой онтологии мы оставляем вопрос конструкции трансфинитных множеств за скобками, мы не знаем, требуются ли для существования таких множеств какие-либо вспомогательные структурные элементы или нет, мы конструируем алгебру для объектов, которые каким-то образом аннотируют исходные объекты, независимо от возможной структуры последних. В итоге, мы получаем описание, объекты которого обладают рефлексивностью, в то время как, исходные множества таким свойством не обладают, в результате чего, экстенсиональность исходных множеств не запутывается с рефлексивностью аннотации, как это происходит в современных теориях. В то же время, кардинальные числа идентифицируемы, их тождество определяется эквивалентностью внутри алгебры дизъюнктивного объединения, с помощью которого, в частности, строятся взаимно однозначные соответствия.

Ровно тоже самое мы делаем для порядковых чисел: $\omega$ - это порядковый тип вполне упорядоченных множеств вида $N = (n_1, n_2, n_3, \dots, n_i, \dots)$, которые в отличии от $\{ \mathfrak{n} \}$ мы можем записать в форме последовательности. Мы точно также не делаем никаких предположений относительно конструктивных особенностей бесконечных множеств, $\omega$ - это просто определённая структура порядка, класс эквивалентности подобных множеств. Экстенсиональность здесь, также позволяет доказать многочисленность множеств в рамках одного класса эквивалентности, благодаря тому, что порядок слагаемых играет роль при конкатенации. Такой подход, позволяет показать, что $\omega$ и $\aleph_0$ - минимальные числа, своего рода \cite{cantor-omega}. Мы не делаем предположений относительно того, существует ли множество с минимальным количеством элементов, или нет, лишь позволяем нашей аннотации иметь такой минимальный объект. То есть, минимальные трансфинитные объекты существуют лишь в рамках конструируемой алгебраической структуры и не требуют существования исходных множеств обладающих схожими свойствами (минимальностью). Благодаря чистоте онтологии, Кантор вводит понятие разности: для двух произвольных порядковых чисел $\alpha < \beta$ всегда существует порядковое число $\beta - \alpha$, удовлетворяющее равенству $\alpha + (\beta - \alpha) = \beta$, фактически позволяя удалять элементы из начала последовательности \cite{cantor-subtraction}. В частности, для произвольных трансфинитного числа $\alpha$ и конечного числа $n$ выполняется $\alpha - n = \alpha$ \cite{cantor-limit-minus}. Когда мы утратили чистоту онтологии в современных теориях, равенство $\omega - 1 = \omega$ перестало иметь какой-либо смысл, поскольку в силу аксиомы экстенсиональности множество без первого элемента не может быть тождественно самому себе, означая, что ординалы, вообще говоря, не изоморфны порядковым типам. При этом данный факт позволяет нам доказать, что минимального множества порядкового типа $\omega$ не существует, в самом деле, если бы такое множество существовало, то удалив из него первый элемент, мы не могли бы получить множество порядкового типа $\omega$. То есть, отсутствие минимального экстенсионально сравнимого объекта является теоремой в теории Кантора, а в современной интерпретации мы пытаемся доказать ровно противоположное. Легко также видеть, почему Кантор настаивает на строгом неравенстве в определении вычитания $\alpha < \beta$, действительно, для класса эквивалентности порядкового типа $\omega$ удаление начального сегмента порядкового типа $\omega$ неоднозначно.

Авторы позволили себе этот экскурс в онтологию Кантора (и близкую онтологию Дедекинда) потому, что считают её единственно возможной основой для осмысленного продолжения его построений и вообще для описания бесконечного. Ещё один вопрос, который Кантор оставляет открытым - это строгое отношение между порядковыми числами $\omega$ и $2^{\omega}$, он доказывает лишь теорему, из которой следует $\alpha \leq 2^{\alpha}$ для любых порядковых чисел, но ни слова не говорит о том, как соотносятся порядковые типы $\omega$ и $2^{\omega}$. В теории Кантора невозможно вывести утверждение современных теорий множеств $2^{\omega} = \omega$. Иными словами, причина по которой Кантор ни разу не упоминает теорему названую его именем заключается в том, что он не может уверенно отнести её ни к бесструктурной кардинальной алгебре, ни к структурированной ординальной. В контексте накопленных на данный момент знаний, данный вопрос решается однозначно, теорема Кантора требует некоторой дополнительной морфологии, отсутствующей в чистой алгебре мощностей. Давайте убедимся в том, что в теории множеств Кантора легко доказывается тот факт, что $\omega < 2^{\omega}$, означающий, что с ростом сложности порядковые типы перестают отображаться не предшествующие порядковые типы, что в свою очередь означает, что множество порядкового типа $2^{\omega}$ нельзя переупорядочить так, чтобы получить множество порядкового типа $\omega$. Это практически невозможно заметить в рамках современных аксиоматических теорий, поскольку в них мы подменили понятие предела Кантора, связав его с супремумом $\alpha = \lim\limits_{\nu \to \beta} \gamma_{\nu} = \sup \{ \gamma_{\nu} \mid \nu < \gamma \} = \bigcup\limits_{\nu < \gamma} \gamma_{\nu}$, где $\gamma_{\nu}$ - фундаментальная последовательность ординалов. Поскольку мы не замечаем того, что $\omega$ является точкой разрыва предикативного описания соотношение $2^{\omega} = \omega$ кажется вполне естественным, ведь нет такой натуральной степени, которая бы позволяла выйти за пределы $\omega$. Но это рассуждение внутренней замкнутой арифметики, кардинальной арифметики, или арифметики Пеано PA которая вместе с индукцией заперта внутри $\omega$, что не позволяет делать выводы о самом $\omega$. В результате в современных теориях мы вводим некоторую функцию ординального возведения в степень $\alpha^{\beta} = \lim\limits_{\nu \to \beta} \alpha^{\nu} = \sup \{ \alpha^{\nu} \mid \nu < \beta \} = \bigcup\limits_{\nu < \beta} \alpha^{\nu}$, она действительно однозначно определена и совпадает с функцией степени для натуральных чисел. Однако, такой подход не соответствует тому, что подразумевает под пределом фундаментальной последовательности сам Кантор. Он совершенно не случайно не разделяет возведение в степень на две операции, хотя технически это вполне возможно: мы можем записать бесконечное произведение влево и получить альтернативную операцию, которая будет соответствовать современной. Однако Кантор сохраняет за возведением в степень $\alpha^{\beta}$ исключительно смысл сквозной индукции, связанной с правыми операциями: тип лексикографического порядка бесконечных последовательностей из элементов множества $\alpha$ длины $\beta$. В современной интерпретации мы используем супремум и не разделяем внутренние и внешние операции, вследствие чего теряем возможность определить какой конкретно объект является пределом последовательности. Так супремумом последовательности $0, 1, 2, 3, \dots$ может быть как ординал $\omega$ так и кардинал $\aleph_0$, так как несмотря на то, что сама последовательность упорядочена, множества являющиеся её элементами могут быть неупорядоченны. Поэтому в качестве предела мы можем получить два совершенно разных объекта с разными свойствами, из-за чего и возникают сложности (как мы знаем, оба объекта являются классами эквивалентности множеств, а не одним множеством). 

При определении через супремум преобладает {\it внутреннее} восприятие свойственное кардинальной алгебре, в то время как, для алгебры порядковых чисел подходит больше концепция Кантора. Он очень точно подметил, что для описания бесконечных множеств как целого нам требуются не внутренние операции, которые не позволяют говорить о свойстве всей совокупности, а какие-то операции, позволяющие взглянуть на множества со стороны, как на законченные объекты, которые мы называем внешними. Кантор связывает своё понятие предела с бесконечной суммой: $\alpha = \lim \gamma_{\nu} = \gamma_1 + (\gamma_2 - \gamma_1) + (\gamma_3 - \gamma_2) + \dots = \sum (\gamma_{\nu} - \gamma_{\nu - 1})$, где $\gamma_0 = 0$ и поскольку последовательность фундаментальна, каждое слагаемое определено, например, $\omega = \lim \nu = \sum 1$ \cite{cantor-limit}. Вспомним, наши рассуждения о том, что сумма легче ассоциируется с актуальной бесконечность из статьи ''Несчётность множества континуум, шах'', а абсолютная бесконечность в отличии от потенциальной как раз и связана с целостностью и законченностью. То есть, Кантор очень тонко чувствует индукцию, бесконечная сумма - это класс эквивалентности сумм, индексное множество которых бесконечно, таким образом, разные классы эквивалентности этого индексного множества дают совершенно разные суммы, которые определяют разные классы эквивалентности того, что называется пределом фундаментальной последовательности. Фактически, он предлагает подняться на мета-уровень и рассматривать предел оттуда, арифметика конечных чисел замкнута внутри кольца, но для внешнего наблюдателя мы представляем её как открытую структуру, что позволяет определить элемент отличный от всех заданных, в точности как в диагональном построении. Отказавшись от онтологии Кантора, мы не смогли предложить ничего подобного, его ключевое наблюдение о том, что числа - это артефакты структуры и разные структуры порождают разные числоподобные объекты, то которое он так яростно защищал и продвигал, осталось не востребованным. Теперь давайте рассмотрим функцию возведения в степень, определённую Кантором, он налагает на неё ключевое условие, если $\lim \gamma_{\nu} = \alpha$, то $\beta^{\alpha} = \lim \beta^{\gamma_{\nu}}$, иными словами, он связывает аддитивный и мультипликативный монойды в единую целостную структуру с общим понятием предела: $\beta^{\alpha} = \beta^{\lim \gamma_{\nu}} = \beta^{\sum (\gamma_{\nu} - \gamma_{\nu - 1})} = \prod \beta^{\gamma_{\nu} - \gamma_{\nu - 1}} = \lim \beta^{\gamma_{\nu}} = \sum (\beta^{\gamma_{\nu}} - \beta^{\gamma_{\nu - 1}})$, где $\beta^{\gamma_0} = 0$ \cite{cantor-power}. Это довольно естественно, ведь Кантору нужно, чтобы показатели степени были объектами порядковой алгебры, а не просто формальными обозначениями, что обеспечило бы возможность строить степенные башни далее вверх, на всех уровнях воспроизводства аддитивной структуры. Если мы обратимся к современным теориям, то там ординальное возведение в степень не имеет такой строгости, например, для $2^{\omega}$ объект $\omega$ в показателе степени не является бесконечным ординалом. Это какой-то промежуточных объект обозначающий сколь угодно большое конечное число, который необходим для того, чтобы поддержать интерпретацию данного ординала как множества функций с конечным носителем. В следующей статье серии мы увидим, почему такая интерпретация связана именно с кардинальной (внутренней) алгеброй. Отсюда непосредственно следует теорема:

Теорема 9: {\sl Никакое вполне упорядоченное множество порядкового типа $2^{\omega}$ нельзя переупорядочить так, чтобы получить порядковый тип $\omega$.}

Иными словами, в алгебре порядковых чисел не существует взаимно однозначного соответствия между множеством типа $\omega$ и множеством типа $2^{\omega}$.

Доказательство: 
Имеем $2^{\omega} = 2*2*2*\dots = 2 + 2 + 4 + \dots + 2^{\nu} + \dots$, данное произведение и данная сумма в точности соответствует лексикографическому порядку бесконечных бинарных последовательностей, которое в свою очередь выражает определение понятия покрытия множества $N$ порядкового типа $\omega$ множеством $M$ из двух элементов $m$ и $w$, задающего функцию $f : N \to M$. Действительно, каждое $\nu$-ое слагаемое бесконечной суммы представляет собой порядковое число множества функций, где варьируются значения только от первых $\nu - 1$ элементов множества $N$, значение от $n_{\nu}$ зафиксировано и равно $m$, а значения от всех остальных элементов зафиксированы и равны $w$. Иными словами, для первого элемента $n_1$ имеем две функции: $f_1: \forall \nu f(n_{\nu}) = w$ и $f_2: f_2(n_1) = m \wedge \forall \nu \neq 1 f_2(n_{\nu}) = w$, для второго элемента $n_2$ добавляется ещё две функции $f_3: f_3(n_2) = m \wedge \forall \nu \neq 2 f_3(n_{\nu}) = w$ и $f_4: f_4(n_1) = m \wedge f_4(n_2) = m \wedge \forall \nu > 2 f_4(n_{\nu}) = w$ и так далее. Таким образом, когда индекс $\nu$ пробегает по всему множеству $\omega$ мы получаем порядковый тип $2^{\omega}$, учитывающий вклад каждого элемента $n_{\nu}$. Но тогда, применяя теорему Кантора к множеству $F$ всех функций $f: N \to M$ приходим к невозможности установить взаимно однозначное соответствие между элементами $N$ и $F$. Теорема доказана.

Несложно видеть, что в рамках порядковой арифметики невозможно иметь какую-либо иную интерпретацию порядкового типа $2^{\omega}$. Каждое слагаемое суммы представляет собой вклад соответствующего элемент, а промежуточная сумма даёт порядковый тип покрытия множества первых элементов (до данного включительно). Для всякого порядкового числа $\alpha < \omega$ порядковый тип $2^{\alpha}$ имеет конечное количество слагаемых и, следовательно, является конечным. Стоит отметить также, насколько отчётливым данный факт становится в свете $\omega$-непротиворечивости, введённой Куртом Гёделем, что связано с тем, что Гёдель прекрасно осознаёт разницу между внутренней индукцией и внешним обобщением всех элементов. В самом деле, в этом случае мы можем утверждать, что индуктивное построение диагонального элемента производимое для конечных чисел, допускает обобщения для всех чисел. Полученный диагональный элемент будет отличен от всех элементов во всех позициях. Если бы диагональный элемент не допускал построения, это означало бы существование формулы $d_i \neq f_i(n_i)$, истинной для каждого конкретного $i \in I$, где $I$ - индексирующая последовательность слагаемых, но ложной для квантора $\forall i \in I \; d_i \neq f_i(n_i)$ (что есть явное нарушение $\omega$-непротиворечивости). Каждый элемент суммы представляет собой вклад нового элемента в комбинаторику комбинаций, а в силу $\omega$-непротиворечивости не существует элемента, который избегает включения в эту систему, таким образом всё бесконечное множество участвует в комбинациях. Более того, этот факт находит подтверждение в самой теории Кантора. Так, Кантор строго формулирует правило: всякое получившееся из элементарных операций равенство между порядковыми типами остаётся справедливым, если в нём все типы заменить мощностями \cite{cantor-ordinal-to-cardinal}. Таким образом, операции над вполне упорядоченными множествами сохраняют комбинаторный смысл операций в кардинальной алгебре, откуда $2^{\omega} = \prod 2 = \sum (2^{\nu} - 2^{\nu - 1})$ (принимая $2^0 = 0$) не допускает иного смысла, кроме покрытия одного множества другим. Данный примечательный факт свидетельствует о глубокой связи между тремя независимыми алгебрами: кардинальных чисел, порядковых чисел и {\it индексных чисел} (к которым можно отнести нестандартные натуральные числа, сюрреальные числа Конвея и гросс-объекты). Кантор указывает на связь между кардинальной и ординальной арифметиками через аннотируемые ими объекты, но он также вплотную подходит к их переплетению с индексной арифметикой - через нормальную форму ординалов, структура которой определяется правыми операциями. Последняя позже обретает полную реализацию в работах Хессенберга с его натуральными операциями \cite{hessenberg-natural}. Не составляет большого труда структурно определить индексные числа, аналогично тому как Кантор вводит кардинальные и порядковые: в качестве базовой операции над вполне упорядоченными множествами взять не конкатенацию, а замещение, когда элементы правого операнда замещают элементы левого в тех же позициях, а замещённые элементы добавляются в конец. Что в сущности даёт онтологический эквивалент нормальных операций с ординалами. Алгебра индексных чисел в их различных проявлениях также поддерживает описанную выше интерпретацию экспоненты. Так Конвей вводит аналитическое возведение в степень, в точности соответствующее порядковой операции $\alpha^\beta$ как типу всех функций $f : \beta \to \alpha$, поскольку класс порядковых чисел - подструктура омнифических целых чисел \cite{conway-omnific-integers}. Это совпадает с экспонентой гросс-объектов Сергеева, дающей точное число элементов континуума: $2^{\raisebox{.5pt}{\textcircled{\scalebox{0.8}{$\omega$}}}}$ \cite{grossone} - и, судя по всему, является стандартным для индексной арифметики в целом.

Стоит отметить, что описанный смысл операции возведения в степень в целом никак не влияет на теорию Кантора в "К обоснованию учения о трансфинитных множествах". Действительно, так как по теореме Кантора $2^{\omega} > \omega^{\nu}$ для любого конечного $\nu$, то $2^{\omega} \geq \omega^{\omega}$, в силу того, что предел фундаментальной последовательности порядковых чисел есть число, непосредственно за ними следующее. Это оставляет единственную возможность $2^{\omega} = \omega^{\omega}$, что полностью сохраняет нормальную форму порядковых чисел. Онтологическая картина Кантора столь изящна, что позволяет вводить не только разность порядковых чисел, но и обратные операции к умножению и, с осторожностью, даже к экспоненте. Так, пары ординалов $(\alpha, \beta)$ образуют кольцо частных с эквивалентностью по натуральным операциям $\frac{\alpha_1}{\beta_1} = \frac{\alpha_2}{\beta_2} \Leftrightarrow \alpha_1\otimes\beta_2 = \beta_1\otimes\alpha_2$ и операциями сложения и умножения: $\frac{\alpha_1}{\beta_1} + \frac{\alpha_2}{\beta_2} = \frac{\alpha_1\beta_2 \oplus \beta_1\alpha_2}{\beta_1\beta_2}$, $\frac{\alpha_1}{\beta_1} * \frac{\alpha_2}{\beta_2} = \frac{\alpha_1\otimes\alpha_2}{\beta_1\otimes\beta_2}$. Используя индексные числа, мы можем придать им вполне определённый смысл. Поскольку индексные числа описывают те же исходные множества, утверждение о том, что трансфинитные ординалы - классы эквивалентности множеств, обретает интерпретацию. Индексному числу $\dot{\omega} - 1$ соответствует совокупность множеств (не одно, поскольку трансфинитные индексные числа - классы эквивалентности, в силу того, что мы строим описание, отвлекаясь от реализации), но их порядковый тип будет равен $\omega$, что и объясняет равенство Кантора $\omega - 1 = \omega$. В каком-то смысле ординалы - классы эквивалентности индексных чисел: индексное число входит в класс ординала, если все его множества входят в класс эквивалентности ординала. Тогда тип индексных чисел $\frac{\dot{\omega}}{2}$ будет либо $\omega$, либо $\omega + \frac{1}{2}$. Здесь проявляется преимущество гросс-объектов: их классы эквивалентности включают только множества, число элементов которых делится нацело на любое конечное число, поэтому $\frac{{\raisebox{.5pt}{\textcircled{\scalebox{0.8}{$\omega$}}}}}{2}$ всегда имеет порядковый тип $\omega$. Аналогично можно ввести логарифм, как наименьший ординал $\gamma$, такой что $\alpha^{\gamma} \geq \beta$ и $\forall \delta < \gamma : \alpha^{\delta} < \beta$, с однозначностью по фундированности. Однако здесь требуется осторожность, строение монады $\omega$ связано с внутренними операциями, которые теорема Кантора ограничивает экспонентой. То есть, внутренняя предикативность $\omega$ полиномиально ограничена, вследствие чего теория не допускает выражения $\log_2 \omega = \omega$ в порядковой алгебре. Оно здесь не имеет смысла, так как данное порядковое число не может быть конечным, но и бесконечным оно быть не может, так как по теореме Кантора оно должно быть строго меньше $\omega$, которое минимально. Дело в том, что операция логарифма нарушает баланс импредикативности и предикативности, который Кантор вложил в определение порядковых чисел. Левые (внутренние) операции предикативны внутри упорядоченных отрезков и импредикативны снаружи, в то время как правые операции предикативны снаружи, что конфликтует с импредикативностью левых операций и разрешается в симметричной импредикативности правых операций внутри. То есть, акт выделения монады $\omega$ скрывает внутреннюю структуру данного объекта позволяя избежать противоречий. В то же время, в индексной арифметике связанной исключительно с правыми (внешними) операциями логарифмы становятся абсолютно законными. Иными словами, внутренняя структура бесконечных множеств де-факто содержит бесконечно большие и бесконечно малые объекты, элегантно скрытые за классами эквивалентности множеств реализующими разные пропорции импредикативности и предикативности. Кантор демонстрирует владение этими инструментами математики, которое в свете современной интерпретации кажется абсолютно недосягаемым. Именно вследствие необходимости соблюдать баланс импредикативности и предикативности Кантор прибегает к формальному подходу при определении функции степени. Ему требуется чистое существование функции, которая отображает область определения - множество порядковых чисел, меньших либо равных $\omega$, - в подмножество элементов вида $2^\nu$ множества порядковых чисел, меньших либо равных $2^\omega$. Структура этих двух разных множеств $\omega$ и $2^\omega$ не позволяют конструктивно указать значения такой функции, превосходящие $\omega$. Отображение порядков этих множеств не имеют выражения внутри порядковой алгебры, это в точности проблема вполне упорядочения континуум. Тем не менее значения функции $2^\nu$ превосходящие $\omega$ непременно существуют, иначе множества совпадали бы, просто они попадают в непредикативные области вполне упорядоченной структуры. Чтобы применить внутреннюю индукцию (в частности, теоремы о перемежающихся последовательностях) к значениям данной функции, требуется предварительно доказать обратное утверждение - что образ $\omega$ замкнут в $\omega$. Кантор явно фиксирует, что понятие перемежающихся последовательностей определено лишь внутри конкретного множества \cite{cantor-alternate}. Следовательно, сначала необходимо установить равенство множеств и лишь затем воспользоваться фактом чередования. Однако равенство множеств невозможно доказать в силу теоремы Кантора. Таким образом, Кантор ясно осознаёт границы применимости внутренней индукции, где внутренние функции в арифметике ординалов работают, а где вследствие теоремы Кантора остаются только внешние. Когда мы выносим монаду - класс эквивалентности $\omega$ - в показатель степени, её внутренняя структура, которая не может быть явно задана в рамках наших построений, тем не менее приводит к тому, что множество значений не содержится во множестве определения. Конечные порядковые числа завершаются индуктивной сингулярностью, после которой начинается класс эквивалентности $\omega$. При помещении данного объекта в показатель степени мультипликативная картина накладывается на аддитивную, мы видим это наложение в конечной области. При этом области регулярности от $\omega$ до $2^\omega$ разделены достаточным количеством сингулярностей, способных скрыть мультипликативные объекты, которые представляют собой элементы внутренней структуры монады $\omega$, невыразимые в нашей теории. Иными словами, импредикативности связанные с вполне упорядочением скрывают элементы связанные с вполне упорядочением в показателе степени. В дальнейшем мы убедимся, что факты такого рода доказываются в достаточно сильных теориях.

Углубляясь в творчество Кантора, невозможно обойти вниманием его представление о континууме, где его интерпретации тяготеют к объекту, который оказывается несчётным в принятом нами смысле. К сожалению, современные исследования редко уделяют должное внимание работам Кантора, посвящённым континууму и его свойствам. В современном представлении множество вещественных чисел в двоичной позиционной записи обладает порядковым типом $2^\omega$, понимаемым как внешний предел всех возможных комбинаций двоичных последовательностей. Это существенно различается со структурой континуума Дедекинда: его сечения - это разбиения множества рациональных чисел $\mathbb{Q}$ на два класса, которые впрочем также можно ограничить функциями $f: \mathbb{Q} \to \{0,1\}$. Однако Кантор развивает более общую концепцию континуума, которая ныне может быть ассоциирована с сюрреальными числами Конвея. Действительно, в своих ранних работах Кантор подробно изучает точечные множества и вводит их производные (множества предельных точек), включая множественные и даже трансфинитные, получая замыкания итеративным включением всех предельных точек \cite{cantor-derivatives}. Сечения Дедекинда соответствуют лишь первой такой производной, в то время как подход Кантора позволяет строить иерархию производных произвольной трансфинитной глубины. Аналогичным образом Кантор подходит к фундаментальным последовательностям, вводя их порядки, избегая тем самым неопределённостей в пределах. Здесь ему приходится оправдываться перед консервативным математическим сообществом, подчёркивая, что тем самым он не вводит новых чисел, а лишь добавляет иное понятийное описание существующим \cite{cantor-sequence-levels}. Мы уже обращали внимание на то, что формулируя диагональный аргумент Кантор не ассоциирует бинарные последовательности с континуумом, он утверждает лишь, что множество таких последовательностей как и континуум не представимо в форме последовательности, явно подразумевая под континуумом нечто большее, чем современная теория. В "К обоснованию учения о трансфинитных множествах" эти идеи обретают законченную форму, здесь он решает ключевую проблему классического анализа - несогласованность аддитивной и мультипликативной структур в определении предела. В традиционном анализе предел определяется чисто аддитивно: $\lim x_n = x \Leftrightarrow \forall \epsilon > 0 \exists N : \forall n > N |x - x_n| < \epsilon$, и ничего не знает о мультипликативной структуре. В результате мультипликативные инварианты не сохраняются: например, $a = \lim \frac{1}{n}$, $b = a + a = \lim (\frac{1}{n} + \frac{1}{n})$, но $3a = 5b$, хотя для всех конечных $n$ такого равенства нет. Это приводит к необходимости введения идеалов в алгебраических структурах, порождающих целый спектр инфинитезимальных описаний. Внешние операции ординалов у Кантора позволяют избежать этой проблемы: $a = \lim n$, $b = \lim n + \lim n$, и мультипликативные инварианты сохраняются - $a*2 = b$. В кольце частных ординалов второго числового класса (т.е. трансфинитных ординалов счётной мощности) $\frac{1}{\omega}$ и $\frac{1}{\omega*2}$ представляют разные порядки малости: они принадлежат идеалу нуля вещественных чисел, но не равны друг другу и не вызывают противоречий. Иными словами, если рассматривать это кольцо как континуум, который искал Кантор, то он действительно имеет несчётную не только мощность, но и ёмкость, при этом замкнут относительно обеих операций - сложения и умножения - без необходимости вводить идеалы. Это даёт более согласованную структуру, чем классический континуум, и идеально согласуется с более детализированной индексной структурой сюрреальных чисел Конвея, где бесконечно малые и бесконечно большие величины интегрированы естественным образом.

Рассмотрим возможность реконструкции современной интерпретации аксиоматических теории множеств на основе подхода Кантора к трансфинитным числам. Мы сознательно отвлечёмся от архитектурного несовершенства связанного с аксиомой бесконечности в её стандартной формулировке; исторически именно эта форма была принята в формальных системах. Следуя Джону фон Нейману, введём ординалы как канонические представители порядковых типов. Хотя такое отождествление является реификацией, оно обеспечивает значительное удобство и наглядность описания. Допустим ряд методологических послаблений, позволяющих установить рефлексивность бесконечных множеств, в частности, внутреннее биективное соответствие между множеством $\omega$ и его собственным подмножеством. Определим кардиналы как минимальные ординалы, равномощные данному множеству (в смысле существования внутренней биекции). Тем не менее, возникает концептуальная трудность при рассмотрении пределов порядковых чисел в контексте операции возведения в степень. В современной теории множеств предел $2^{\omega}$ при ординальном возведении в степень записывается как $2^{\omega} = \lim\limits_{\nu \to \omega} 2^{\nu} = \sup \{ 2^{\nu} \mid \nu < \omega \} = \bigcup\limits_{\nu < \beta} \alpha^{\nu}$, где супремум интерпретируется внутренним образом: $\sup \{ 2^{\nu} \mid 2^{\nu} < \omega \}$. Это означает, что в рамках внутренней алгебры ординалов невозможно указать элемент $x \in \omega$ такой, что $2^{x} > \omega$, вследствие чего $2^{\omega} = \omega$. Напротив, во внешней алгебре, соответствующей исходному подходу Кантора, предел $2^\omega$, определяемый через бесконечный ряд, превосходит $\omega$. Корректный супремум в этом случае должен удовлетворять теореме Кантора: $2^{\omega} = \sup \{ 2^{\nu} \mid \nu < \omega \} > \omega$, ведь последовательность $2^{\nu}$ принадлежит множеству $2^{\omega}$, несмотря на то, что её предикативная часть полностью заключена в $\omega$. Впрочем, данная коллизия не представляет существенной проблемы. Мы просто определяем операцию возведения в степень не через внешние суммы (как у Кантора), а через внутренний супремум. Речь идёт о двух различных операциях, обозначаемых одинаково. Одна из них соответствует комбинаторной (внешней) алгебре, другая - внутреннему способу упорядочения в рамках ординальной структуры. Впрочем, здесь мы должны решить, когда правомерно использовать одну интерпретацию, а когда другую. Следуя за Кантором мы объявляем ординал $2^{\omega}$ счётным, но у Кантора это обусловлено кардинальной алгеброй в которой биекция между множествами существует, в то время как в порядковой алгебре структуры структуры этих множеств не могут быть изоморфно сопоставлены друг другу. В оригинальной теории это работает для объекта $2^{\omega}$ как множества бесконечных последовательностей бинарных строк. Но вследствие того, что в современной интерпретации мы слили обе алгебры в одну, множество степени счётного множества не может быть счётным. Выходом из данной ситуации становится концепция множества функций с конечным носителем, для которой нет никакого аналога в теории Кантора. В рамках современной интерпретации существование таких множеств обосновывается именно через теорему о счётности счётного объединения счётных множеств $\aleph_0 = \aleph_0*\aleph_0$, которую Кантор формулирует в кардинальной алгебре. Но если эта теорема не верна в порядковой алгебре в общем случае, факт существования этих промежуточных объектов повисает в воздухе. Простого определения множеств функций с конечным носителем явно не достаточно для того, чтобы они стали объективной реальностью теории. Так, например, определение максимального натурального числа не гарантирует его существование. Более того, разделение на $2^{\omega}$ и $2^{\aleph_0}$ имеет смысл в теории Кантора, где эти выражения относятся к двум различным алгебрам, но в современной интерпретации $\aleph_0$ - это в точности $\omega$ - первый счётный ординал, что делает два этих выражения синтаксически тождественными. В следующей статье серии мы разберёмся, какая именно интуиция стоит за идей множества функций с конечным носителем - это идея потенциальной бесконечности, обладающей внешней импредикативностью, что позволяет ассоциировать эту интуицию именно с выражением $2^{\aleph_0}$ в кардинальной алгебре Кантора. 

Однако ситуация становится ещё сложнее в случае с ординалом $\omega^{\omega} = \sup \{ \omega^{\nu} \mid \nu < \omega \}$, являющимся непосредственным преемником всех $\omega^{\nu}$, что совпадает с внешним пределом и не допускает биекции с $\omega$. По смыслу это именно тот ординал, который превосходит $\omega$, возведённое в любую конечную степень. Интуитивно $\omega^{\nu}$ можно представить как множество всех векторов длины $\nu$ с координатами из $\omega$. Тогда пространство с бесконечным числом измерений естественно оказывается строго больше любого пространства с конечным числом измерений, и между областью конечномерных пространств и пространством бесконечной размерности не существует промежуточных объектов. Но это именно то, что Кантор называет покрытием одного множества другим, следовательно теорема Кантора работает для множества $\omega^{\omega}$, не позволяя упорядочить его в множество изоморфное $\omega$. Нет никакого пространства с произвольным конечным но не бесконечным количеством измерений, которое могло бы поддержать идею множества функций с конечным носителем. Путаница с этим понятием, заключающаяся в попытке отождествить внешнюю импредикативность объекта $\aleph_0*\aleph_0 = \aleph_0$, свойственную кардинальной алгебре, с объектами порядковой алгебры, которые обладают внешней предикативностью $\omega*\omega \neq \omega$, отражает неудовлетворительное владение инструментарием импредикативности и предикативности в рамках современной интерпретации теории множеств. При переходе к кардиналам (внутренней алгебре) ситуация вырождается: для бесконечных кардиналов $\aleph^{i}$ выполняется $\aleph^{i} = \aleph$, и супремум сводится к супремуму одноэлементного множества. Это плохо согласуется с тем, что ординал $\omega^{\omega}$ строго превосходит все ординалы $\omega^{\nu}$ в бесконечной возрастающей последовательности. Иными словами, если определять мощность множества через возможность биекции внутри упорядоченного универсума фон Неймана (т.е. в терминах ординалов), то естественно получить $\aleph_1 = \omega^\omega$. В то же время, в терминологии самого Кантора ординал $\omega^\omega$ остаётся счётным (в смысле ёмкости), поскольку существует класс пар задающих неупорядоченную биекцию в кардинальной арифметике. Здесь возникает ключевая развилка. Переопределяя предел через внутренний супремум, мы вводим существенно иную операцию возведения в степень, для которой $i^{\omega} = \omega$ при любом конечном $i$. В этом нет принципиальной ошибки - это всего лишь определение. Однако такое определение не позволяет ассоциировать комбинаторные объекты с $i^{\omega}$; в частности, множество вещественных чисел в двоичной записи уже нельзя обозначать как $2^{\omega}$. Продолжая в том же духе, мы распространяем внутреннюю кардинальную алгебру (по сути, арифметику Пеано) на объект $\omega^{\omega}$ и доказываем, что $|\omega^{\omega}| = \aleph_0$. Но для доказательства данного факта, нам требуется теорема кардинальной алгебры о том, что счётное объединение счётных множеств счётно $\aleph_0\aleph_0 = \aleph_0$. Мы можем доказать данный факт в кардинальной алгебре, но в порядковой алгебре, конкатенация бесконечного числа счётных отрезков по теореме Кантора может не иметь отображения на $\omega$ (например, $\omega + \omega^2 + \dots = \omega^{\omega} > \omega$). Мы видим, что область применения кардинальных теорему существенно ограничена. Если мы используем подход Кантора с определением высоты многочлена, то количество многочленов определённой высоты растёт экспоненциально, и доказательство в порядковой алгебре приводит к порочному кругу: счётность экспоненциального предела мы доказываем через счётность экспоненциального предела. Но и матричное доказательство в порядковой алгебре не проходит, так как предел последовательности порядковых типов $\omega^i$, которая имеет порядковый тип $\omega$, может быть не изоморфен $\omega$, но тогда индексация столбцов уже требует не изоморфного типа. Тем самым мы пытаемся подражать Кантору, но в современной интерпретации у нас отсутствует даже язык для описания основных математических инструментов и манипуляций с ними, которые Кантор умело использует, - импредикативности и предикативности. 

Кардинальные и порядковые числа Кантора принадлежат двум различным алгебрам, мы же выражаем кардинальную алгебру через ординальную, и сливаем правое и левое сложение в единую внутреннюю операцию $x \cup \{ x \}$. При этом, не задумываясь, проецируем кардинальную логику Кантора на новое суррогатное понятие кардинального числа, возникшее в нашей теории вследствие такого упрощения. Но при строгом рассмотрении даже в контексте современной интерпретации мы сталкиваемся с проблемой, заключающейся в том, что $\omega$ является точкой прерывности индукции, связанной с аксиомой бесконечности, требующей особенно аккуратного доопределения операций. Мы находимся в упорядоченном универсуме, где свойства конечных множеств в $\bigcup\limits_{ x \in \omega } x$ автоматически не распространяются на предельное множество. Тогда что заставляет нас думать, будто свойство счётности всех элементов $\bigcup\limits_{ x \in \omega } \omega^x$ появляется и у предельного объекта $\omega^{\omega}$? Мы не можем переупорядочить каждый конечный элемент, чтобы он имел тип $\omega$, поскольку с точки зрения ординальной алгебры $\bigcup\limits_{ x \in \omega } \omega^x$ и $\bigcup\limits_{ x \in \omega } \omega$ являются совершенно разными объектами. Мы даже не можем утверждать, что трансфинитная рекурсия позволит нам индуктивно переупорядочить все элементы последовательности $\omega^x$ в порядковый тип $\omega$, поскольку термин все элементы в ординальной алгебре связана с пределами, которые не наследуют предикативных свойств элементов автоматически. Переупорядочение способно вернуть нас к кардинальной морфологии, где мы абстрагируемся от порядка, но мы не выделяем кардинальную алгебру внутри упорядоченного универсума, как это делал Кантор. Стандартное доказательство $\aleph_0\aleph_0 = \aleph_0$ использует внутреннюю индукцию и аксиому выбора. Однако использование аксиомы выбора с внутренней индукцией не конструктивно для доказательства счётности несчётного множества, теорема Кантора выступает здесь универсальным ограничителем границ использования внутренней индукции. Мы не можем сослаться на наше определение множества функций с конечным носителем, поскольку получим порочный круг, так как мощность таких множеств мы выводим именно из данной теоремы. Аксиома выделения не позволяет утверждать, что такие множества семантически отличны от полной степени в отношении мощности, вследствие разрыва предикативности при переходе к пределу. Чтобы объединение конечных чисел $\bigcup\limits_{ x \in \omega } x$ на уровне аддитивной структуры и объединение конечных носителей $\omega^{\bigcup\limits_{ x \in \omega } x}$ на первом уровне мультипликативного воспроизводства аддитивной структуры редуцировались в теории в семантически разные объекты, нам требуется дополнительное соглашение. Для Кантора принципиальное значение имеет согласование аддитивного и мультипликативного монойдов. Именно факт того, что синтаксическое выражение редуцируется в один и тот же семантический объект обеспечивает требуемое воспроизводство аддитивной структуры сквозь всю иерархию. При различной редукции синтаксического описания, мы должны доказать, что объект в показателе степени структурно изоморфен аддитивному. Однако это не так, поскольку добавление элемента во все наследственно конечные множества оставляет эти множества наследственно конечными, что свидетельствует об изоморфизме кардинальной структуре, а не порядковой. То же самое, мы наблюдали в случае множества $P_{\omega}$: добавление элементов в наследственно конечные множества замыкало построение, что также свидетельствует о путанице кардинальной и порядковой структур в теории. 

Аксиома выбора позволяет нам проиндексировать все элементы счётного объединения счётных множеств некоторой трансфинитной последовательностью ординалов, но она не способна указать на мощность данной последовательности. Это не вопрос рефлексивности, поскольку теорема Кантора в любом случае должна ограничивать рефлексивность, это вопрос целостности нашего описания. Объекты аннотации могут быть рефлексивными, но задача как раз в том и заключается, чтобы создать инструментарий позволяющий устанавливать пределы этой рефлексивности в контексте определённой алгебраической структуры. Кантор разработал специальный предельный аппарат, позволявший корректно и интуитивно понятно описать поведение индукции в трансфинитном случае. Мы же, не вводя никаких дополнительных предположений, по сути отказались от теоремы Кантора и применили внутреннюю кардинальную индукцию к внешней порядковой алгебре. Но даже с учётом принятия всех нюансов современной интерпретации у нас нет никаких оснований внутри теории с упорядоченным универсумом, чтобы сделать заключение о счётности любого счётного объединения счётных множеств. Такая кардинальная теорема требует дополнительных предположений, осознание чего было существенно затруднено перечисленными в данной работе архитектурными дефектами современных теорий, ошибкой реификации, не владением аппаратом сочетания импредикативности и предикативности и игнорирование методологических формальных ошибок преждевременного обобщения и ложной дихотомии. В результате мы утратили изоморфизм между порядковыми типами в понимании Кантора и ординалами современных теорий множеств. Это произошло несмотря на то, что сам Кантор неоднократно подчёркивал принципиальную невозможность сведения кардинальных и порядковых чисел к единому понятию числа и жестко критиковал других авторов за подобные попытки. С современных позиций мы даже можем сказать почему их невозможно объединить, они требуют двух различных вариантов бесконечности в пространстве альтернатив $\aleph_0$ и $\omega$, таким образом пространство альтернатив становится трёхзначным: конечное, кардинально бесконечное и порядково бесконечное. Ошибка ложной дихотомии становится фатальной, отрицание свойства порядковой бесконечности не обязательно означает конечность, а двойное отрицание не обязательно возвращает нас к порядковой бесконечности, так отрицание свойства кардинальной бесконечности, которая не является порядковой может указывать на конечность, а вовсе не указывать обратно на упорядоченную структуру. Бинарная логика перестаёт быть применимой в такой теории, мы обязаны использовать более слабые логики. Кантор не совершает формальной ошибки ложной дихотомии, использование классической логики внутри кардинальной алгебры абсолютно корректно равно как и внутри порядковой алгебры. Но их слияние уничтожает основное методологическое условие использования классической логики - двухзначность пространства альтернатив. В результате, мы видим конфликт внутренней замкнутости и внешней теоремы Кантора, но разрешаем его как нам хочется, а не через исследование вариантов. Бесконечное множество действительно замкнуто внутренне, но такая индуктивная замкнутость не определяет внешних свойств бесконечных объектов из-за разрыва предикативного описания. Кантор показал что в чистом упорядоченном универсуме, какими мы считаем универсумы ZFC и NBG, внешние операции в алгебре начинают доминировать при пересечении экспоненциального барьера. Но мы отказываемся от методологически строгих рассуждений Кантора, полагая что разрешение противоречия между внутренней индукцией и внешней теоремой Кантора в пользу первой - вопрос выбора, а не следствие фундаментальных законов. Однако, именно этот не зафиксированный формально выбор блокирует усиление и дальнейшее развитие наших теорий. Более того, в действительности, в наших теориях присутствует внешнее описание, ведь из-за того, что мы предпочитаем не замечать объекты, которые семантически могут быть интерпретированы как внешние функциональные отношения, такие объекты не перестают существовать. Мы пользуемся исключительно внутренними операциями, которые являются рефлексивными, но альтернативные классы объектов могут определять альтернативные функциональные отношения, допуская восстановление внешних операций Кантора. Которые также необходимо согласовывать со следствиями преждевременных обобщений современных теорий. Так или иначе, авторы полагают, что единственный путь заключается в фиксации неявных предположений современных теорий в форме дополнительных аксиом. Такие расширения ZFC и NBG станут их классическими интерпретациями, сохраняющими все полученные к настоящему моменту результаты на более строгом фундаменте. Альтернативные расширения позволят реализовать исходную интуицию Кантора, которая судя по всему, окажется совместимой практически со всеми маргинализированными подходами. Этот путь позволит объединить многие разрозненные на данный момент течения в единую теоретическую рамку.

Вопрос о существовании и правомерности альтернативной операции возведения в степень, ставшей неотъемлемой частью современных аксиоматических теорий множеств, авторы оставляют открытым. Дело в том, что подход Кантора предполагает, что класс эквивалентности предела должен начинаться с первого элемента, который не принадлежит ни одному из предшествующих классов эквивалентности в последовательности. Однако такая конструкция крайне трудно поддаётся строгой аксиоматической формализации в силу импредикативности предельного конца. В современной математике отсутствует адекватный аппарат для описания сингулярности, возникающей на конце бесконечной последовательности. По-видимому, подобная сингулярность, разделяющая области регулярности различных алгебр, носит фундаментальный характер и не может быть устранена в рамках стандартных аксиоматических средств. Мы видим, что редукция синтаксических выражений к определённым семантическим объектам - это предмет соглашения, которые иногда даже не фиксируются в явном виде. Это обстоятельство оставляет пространство для применения метода форсинга, который позволяет населить указанную сингулярность элементами с требуемыми свойствами и тем самым влиять на поведение специальных операций. В частности, сингулярность предельного перехода $\lim_{\nu \to \omega} \omega^\nu = \omega^\omega$ может в принципе вмещать всю иерархию счётных ординалов, рассматриваемых в современных теориях, в силу относительности самого понятия мощности. Иными словами, экспоненциальный скачок мощности на одном масштабе, способен вместить всю иерархию другого масштаба, поскольку само понятие масштаба точно также относительно. Без усиления теории функциональными отношениями, добавляющими внешней предикативности, данная относительность фундаментально неустранима и может оставаться легитимной внутри таких систем. В этом случае статус континуум-гипотезы в подобных построениях сохраняет свою актуальность как открытая проблема. В рамках согласованной картины трёх алгебр - кардинальной, ординальной и индексной - можно предложить следующую интерпретацию причин возникновения данной проблемы в стандартных теориях. Рассмотрим первое множество, мощность которого превосходит $(\frac{3}{2})^{\dot{\omega}}$. Оно располагается внутри сингулярности перехода между $\omega$ и $2^\omega$, которую мы не умеем адекватно описывать в аксиоматических терминах. Следовательно, если с помощью форсинга добавить элемент, соответствующий положению такого множества, то его невозможно будет ни счётно достичь из $\omega$, ни счётно достичь $2^\omega$ из него. Именно это обстоятельство, по-видимому, и объясняет независимость континуум-гипотезы от аксиоматик ZFC и NBG. Более того, согласно теореме Лёвенгейма–Скулема, внутри любой счётной модели мы способны видеть всю иерархию мощностей. Таким образом, наше внутреннее представление иерархии мощностей вполне может быть реализовано уже внутри $\omega$. Однако, как будет показано далее, в таком случае построенная теория оказывается фундаментально слабой, не допуская увеличения предикативности. Более того, воссоединение с альтернативными (в том числе более сильными) теориями становится проблематичным.

Таким образом, неоднозначность интерпретации в современной теории множеств напрямую связана с архитектурными дефектами, ошибкой реификации и утратой исходной онтологии Кантора. Эти ошибки приводит к тому, что взаимно однозначное соответствие интерпретируется не только как свидетельство подобия различных алгебраических структур, но и как указание на тождество самих множеств в случае, когда сопоставляемые элементы, полученные индуктивным построением, совпадают. От оригинального доказательства Кантора обсуждаемое здесь отличается прежде всего формализацией биекции в виде функции $\varphi(x)$, без учёта принадлежности этой функции конкретной алгебре. Такая функциональная трактовка существенно затрудняет интуитивное понимание того, каким образом взаимно однозначное соответствие может не означать равенства самих множеств. Мы привычно предполагаем, что функция, определённая на конечной области, тривиальным образом расширяется на бесконечную независимо от алгебраических структур. Однако парадокс Рассела в его интенсиональной форме показывает, что подобное расширение далеко не всегда тривиально. В этой связи представляется плодотворным обратиться к продолжению известного апории Зенона об Ахиллесе и черепахе. Допустим, что после того, как Ахиллес догнал черепаху, она заметила: ''Не спеши радоваться, друг мой; возможно, если ты продолжишь движение вперёд, а я поползу за тобой, то в итоге именно я догоню тебя''. Уязвлённому Ахиллесу не остаётся ничего иного, кроме как продолжить путь. Чтобы определить, сможет ли черепаха догнать Ахиллеса, воспользуемся аппаратом взаимно однозначного соответствия. Ахиллес и черепаха равномощны: каждый из них преодолевает любое расстояние, которое преодолевает другой, независимо от способа его измерения. Однако означает ли это, что в итоге они пройдут одно и то же расстояние и черепаха действительно догонит Ахиллеса? Если путь черепахи представить в виде последовательности столбцов, каждый из которых соответствует одному метру, а путь Ахиллеса - в виде последовательности строк той же длины, следует ли отсюда, что множество индексов столбцов равно множеству индексов строк? Сколько существует способов исключить время из данного построения? Должны ли мы считать какой-либо из них принципиально предпочтительным? И если один из таких способов приводит к выводу, что две непараллельные прямые могут иметь более одной точки пересечения, следует ли вообще рассматривать этот способ в качестве достоверного? Мы уже используем факт неравенства равномощных множеств при объяснении парадокса колеса Аристотеля. Действительно, если взаимно однозначное соответствие означало бы тождество множеств (оба множества состояли бы из одних и тех же точек), то понятие меры утрачивало бы всякий смысл: все атрибуты тождественных множеств должны были бы совпадать, ведь в противном случае множества не могут считаться равными. Рефлексивность бесконечных множеств также требует более осторожного истолкования. Когда мы утверждаем, что бесконечное множество содержит бесконечное подмножество того же типа (структурно эквивалентное исходному), это интуитивно приемлемо. Однако утверждение, что такое подмножество может быть преобразовано во множество, равное исходному, представляет собой нетривиальное допущение. Требует ли данный факт специального доказательства или, по крайней мере, принятия в качестве аксиомы некоторого более общего утверждения, из которого он логически следует? Инфинитезимальный анализ демонстрирует, что отбрасывание информации о способе построения объекта не всегда оправдано. В частности, сведения о скорости сходимости последовательности позволяют различать объекты внутри одного класса эквивалентности вещественных чисел и оперировать ими по-разному. В этой связи возникает вопрос: почему мы считаем безопасным отбрасывать информацию о построении множества (например, состоит ли оно из чётных чисел или из удвоенных натуральных чисел)? Осознаём ли мы, в какой момент отбрасываемая информация действительно несущественна для данного контекста, а в какой её утрата вносит принципиальную неопределённость в описание действительности? Наука в значительной степени представляет собой обмен контролируемым любопытством. Любопытство по своей природе наивно и инстинктивно; подвергая его чрезмерно жёсткой формализованной критике, мы рискуем подорвать основания собственного познания.

Основная проблема заключается не в самом доказательстве Кантора, а в неудачной интеграции его в контекст аксиоматических теорий, в результате чего оно стало применятся к объектам, к которым оно в действительности не применимо. В рамках формальных теорий мы оперируем абстрактными структурами, которые представляют собой лишь аннотации (или представления) объектов, а не сами объекты в их полноте. Различные аннотации одних и тех же объектов могут обладать существенно различными свойствами, и приписывание этих свойств самим исходным объектам является методологически некорректным. Множество всех подмножеств данного множества может быть эквивалентно исходному множеству в смысле идентичности структуры в рамках одной алгебры, в то время как в рамках иной алгебры структуры этих объектов оказываются принципиально различными. В свете данных наблюдений приходится констатировать, что разногласия между математиками относительно корректности теоремы Кантора обусловлены в значительно большей степени не математическими, а герменевтическими и методологическими причинами - прежде всего недостаточной готовностью внимательно рассматривать позицию оппонентов. Теорема Кантора безусловно верна в одних алгебрах (системах операций и отношений), но не верна в других. Более того, она верна в одних интерпретациях ZFC и NBG, в то время как не верна в других. С точки зрения булевой логики это делает её некорректной в контексте данных систем. Соответственно, те, кто стремился ограничить или опровергнуть универсальную применимость теоремы, в определённом смысле были ближе к истине, чем те, кто настаивал на её абсолютном и фундаментальном характере. Примечательно, что в оригинальных работах самого Георга Кантора наблюдается постоянная эволюция взглядов, а в поздний период - явные сомнения. Для него мощность множества представляла собой всего лишь некоторый класс эквивалентности, в рамках которого, как и в любом другом классе эквивалентности, неизбежно происходит отбрасывание части информации, несущественной для выбранного контекста. При этом Кантор последовательно воздерживался от выхода за рамки логически обоснованных допущений и, в частности, никогда не применял диагональный аргумент непосредственно в кардинальной арифметике. Если попытаться понять существо многочисленной критики теоремы Кантора, не отмахиваясь от неё и не прикрываясь формализмом, который не имеет прямого отношения к обоснованности той или иной интерпретации, то становится очевидным: под некорректностью теоремы большинство критиков подразумевали именно существование альтернативных способов аннотации множеств - альтернативных алгебр, которые уже являются частью современных формальных систем. Наиболее парадоксальным фактом, связанным с теоремой Кантора, остаётся то, что в конечном итоге её в той или иной форме отвергали как её противники, так и её сторонники. Противники, как правило, отрицали её тотальный характер и универсальную применимость к понятию мощности в неупорядоченном (кардинальном) описании, то есть утверждение $\aleph_0 < 2^{\aleph_0}$. Сторонники же, напротив, отвергали её именно в ординальной алгебре, где она проявляет свою подлинную силу и гарантирует невозможность для множества и множества всех его подмножеств обладать одной и той же упорядоченной структурой: $\omega < 2^\omega$. Впрочем, авторы не питают иллюзий относительно того, что подробный разбор методологических упрощений, лежащих в основе современной математической идеологии, оказался достаточно ясным и убедительным, чтобы сделать альтернативное восприятие теоремы Кантора широко доступным математическому сообществу. Научная среда остаётся в значительной степени идеологизированной и догматизированной, что затрудняет развитие.

Неопределённости в интерпретации теоремы Кантора разрешаются относительно просто, хотя последствия их систематического игнорирования оказываются чрезвычайно серьёзными. Проведённый нами анализ высвечивают ещё одну фундаментальную проблему логического характера, связанную с данной теоремой и проявляющуюся в антиномии Рассела и аналогичных парадоксах. Широко распространено убеждение, что антиномия Рассела обусловлена размером множества (его бесконечностью). Однако, в действительности, подобные антиномии возникают не вследствие размера множества, а вследствие особенностей его структуры и способа её интерпретации в рамках того или иного формального языка. Возможность сформулировать антиномию Рассела даже для конечных множеств указывает на то, что корень проблемы лежит в понятии абстракции и, в частности, в её предельном случае - абсолютной бесконечности, которые в математической практике используются весьма свободно, хотя и нуждаются в строгой концептуальной и формальной проработке. Примечательно, что финитарные (конструктивные) методы, напротив, оказываются свободными от подобных парадоксов и допускают неограниченное последовательное продолжение. Здесь мы сталкиваемся с классическим противопоставлением потенциальной и актуальной бесконечности, импредикативности и предикативности, а также внутренней замкнутой и внешней открытой структур. Вместо тщательного исследования этих противоречий современная аксиоматика предпочитает маскировать их введением соответствующих аксиом и идеологических табу. К сожалению, интерпретационные допущения, ответственные за указанные логические противоречия, лежат значительно глубже тех, что связаны с множественностью алгебр описания множеств. Для их полноценного анализа потребовалось бы радикальное переосмысление всего фундамента математики и теории множеств. Поскольку настоящее исследование уже погрузило нас в эту проблематику, обойти столь значимую математическую загадку стороной представляется невозможным, что вынуждает авторов не ограничиваться данной серией работ и продолжить изыскания. Расширение внутренней индукции на внешний экспоненциальный контекст может быть, может быть без труда постулировано в рамках современных систем ZFC и NBG. В то же время, альтернативные расширения, возвращающие нас к онтологии Кантора, допускают усиление теории не только через согласованность различных структур, но допускают усиление внешнего контура практически до соприкосновения с реальностью через предположения аналогичные аксиоме Евклида. В этой онтологии гибридные структуры классов эквивалентности множеств естественным образом становятся ординалами, сосуществующими с внешними индексными и внутренними кардинальными морфологиями. Однако теории, сохраняющие неконтролируемую рефлексивность, продолжают порождать картины, содержащие принципиально невыразимые объекты. Такие объекты требуют привлечения метатеоретических утверждений и истин, лежащих за пределами доказательной силы самих теорий. В частности, допустимо существование объектов теории, которые невыразимы средствами её собственного языка и истинных утверждений, которые не могут быть доказаны средствами самой теории (если утверждения ложны, то системы ZFC и NBG оказываются противоречивыми). Теории, содержащие подобные конструкции, но не допускающие предикативного усиления, мы будем называть недостоверными, воспользовавшись терминологией Лёйтзена Брауэра. Во времена Брауэра ещё не были разработаны современные понятия непротиворечивости и полноты формальных систем, поэтому он апеллировал непосредственно к фундаментальным законам логики. Обнаруженные нами дефекты теории проливают свет на природу открытого Брауэром нарушения закона исключённого третьего. Данный эффект обусловлен неоднозначностью формальной системы аксиом как единого целого, в частности в теории множеств неоднозначным оказывается само отношение принадлежности. Альтернативные интерпретации такой системы имеют равноправный статус и не исключают друг друга, позволяя части синтаксических конструкций редуцироваться к семантике одной альтернативы в определённых фрагментах теории и к семантике других альтернатив за их пределами. В настоящее время мы знаем, что аксиоматические теории множеств не являются полными, и это фундаментальное свойство всех импредикативных систем. Однако, как мы выясним среди таких систем существует те, которые допускают предикативное усиление, и те, предикативное усиление которых ведёт к противоречию. Именно термин Брауэра оказывается наиболее удачным для описания сущности явления, возникающего в случае неустранимой неполноты подобных теорий.

Мы проанализировали все основные доказательства несчётности множества вещественных чисел и установили, что они правомерны лишь в определённых контекстах и при соблюдении соответствующих методологических допущений. Все остальные опубликованные доказательства - по крайней мере те, которые удалось идентифицировать авторам - также подчиняются аналогичным ограничениям, разобранным в данной статье и статьях ''Несчётность множества континуум, шах'' и ''Несчётность множества континуум, ещё шах''. Поиск оказался несложным, благодаря тщательной работе Кристины Кнапп и Цезаря Сильва, собравших и систематизировавших варианты доказательств в своей публикации \cite{proof-list}. Так, в доказательстве основанном на игре, предложенной Гроссманом и Туретом, которое приводит в своей работе Бэйкер \cite{baker}, Элис и боб пытаются построить последовательности разделяющие множество, ровно такие же последовательности строил и Кантор. Рассуждения по существу соответствуют рассуждениям геометрических доказательств, Боб и Элис исходят из предположения, что существует внешняя алгебра позволяющая одновременно описывать последовательность ходов и множество точек аннотирующее непрерывный отрезок, такая, что завершающий ход Боб или Элес не принадлежит тому же классу эквивалентности, которому принадлежит множество точек, и что алгебра замкнута, то есть, игра когда либо закончится, и Боб и Элис в какой-то момент не будут способны продолжать ходы, как в случае с сюрреальными числами. Такая игра более наглядно демонстрирует конфликт импредикативности (замкнутости) континуума и предикативности (тотальности и универсальности) его упорядочения. Таким образом, правилами игры необходимо предусмотреть возможность определять победителя в случае различной аннотации бесконечного или в случае невозможности завершения игры за отведённое бесконечное время. При таких обстоятельствах, остаётся нерешённым весьма любопытный вопрос: если множество непрерывно, есть ли стратегия позволяющая одному из игроков заставить другого игрока сделать предельный ход на бесконечном шаге? В более общем случае, как различные аксиомы (они же правила игры), накладывающие ограничения на разделение множеств, в частности на возможность тех или иных сечений множеств, влияют на существование выигрышной стратегии для одного из игроков?

Доказательство, опубликованное Веннером \cite{wenner}, является, как признаёт сам автор, стандартным доказательством несчётности множества вещественных чисел, переформулированным в классы эквивалентности последовательностей Коши. Действительно, если мы принимаем аксиому непрерывности на множестве классов эквивалентности последовательностей Коши рациональных чисел, а также, как предложено в статье вводим алгебраические операции и отношение порядка стандартным образом, то из этого в силу теорема 2 из статьи ''Несчётность множества континуум, ещё шах'' следует, что последовательностей классов эквивалентности разделяющих множество не может существовать. Таким образом, само построение точно также как в классическом случае является неконсистентным и нуждается в предварительных предположениях о структуре бесконечности. Все дело в том, что не предполагая структуру бесконечности мы не можем утверждать, что одна бесконечность больше другой, но классическое определение фундаментальной последовательности не предполагает никакой структуры, здесь бесконечность - это сингулярность $\infty$ или $\aleph_0$, внутри которой каким-то образом происходит разрешение конфликта аксиомы непрерывности и последовательности вложенных интервалов. Математически анализ без теории множеств способен только указать на данные конфликт импредикативности непрерывного и предикативности дискретного, но он не гарантирует того, что разрешение этого конфликта находится в предикативной зоне и может быть выражено на языке теории. До того момента как мы предположили структуру бесконечности, у нас просто нет описания сингулярности, а следовательно нет и не может быть никакого противоречия. Для того, чтобы рассмотреть как именно происходит разрешение, необходима внешняя алгебра, которая позволит интерпретировать любую последовательность Коши как законченный упорядоченный объект типа $\omega$, для которого существуют ещё большие объекты. Только тогда, окажется, что совокупность всех последовательностей Коши расположенных друг за другом не могут формировать тот же класс эквивалентности $\omega$, что и одна единственная последовательность.

Доказательства основанные на понятии меры, например, опубликованное Стилуэллом \cite{continuum-measure}, которое он относит Гарнаку, в равной степени с доказательствами Кантора исходят из того, что бесконечность имеет предикативную структуру. В частности, мера множества считается равной нулю, если существует сжимающая последовательность интервалов порядкового типа $\omega$, совокупная длина которых меньше любого наперёд заданного сколь угодно маленького числа $\epsilon > 0$. В этом случае задача сводится к раскрытию неопределённости вида $0*\infty$, если $\infty_c$ - счётная бесконечность, то неопределённость $0*\infty_c$ разрешается в пользу нуля. Напротив, неопределённость вида $0*\infty_u$, где $\infty_u$ - несчётная бесконечность, может разрешаться в конечное число. По существу, данные рассуждения воспроизводят факт, очевидный в более предикативных алгебрах: выражение $\frac{n^{\dot{\omega}}}{\dot{\omega}^m}$ принадлежит аддитивному идеалу нуля для всех конечных $m$ и $n$, однако этот факт недоказуем в более импредикатиавных алгебрах. Кантор вновь демонстрирует глубокое опережение своего времени, приводя пример множества, носящего его имя, которое иллюстрирует, что троичный континуум по мере превосходит двоичный, то есть $\frac{2^{\dot{\omega}}}{3^{\dot{\omega}}}$ равно действительному числу 0 \cite{cantor-cantorset}. Проблема, однако, состоит в следующем. Если континуум описывается как точечное множество, то мера оказывается неразрывно связанной с его внешней структурой и неизбежно вступает в конфликт с замкнутостью внутренней алгебры. Отсюда вытекает общее правило раскрытия неопределённостей указанного вида для мер в пространствах любых мощностей. Данный пример наглядно показывает, что все подобные рассуждения опираются на неявное принятие структурного ограничения на последовательность интервалов определяющую нулевую меру. Таким образом, мы имеем дело не с превосходством мощности, а с взаимно исключающими требованиями импредикативности и предикативности. Если же определить меру на классе сюрреальных чисел, то окажется, что континууму вообще нельзя приписать какую-либо определённую мощность. В историческом контексте такая ситуация вполне объяснима. Понятие меры появляется в работах Жордана \cite{jordan}, Пеано \cite{peano} и Лебега \cite{lebesgue}, опубликованных существенно позднее трудов Кантораи в значительной степени опирающихся на его результаты, а также на результаты Дедекинда. Все эти авторы уже работают с континуумом, представленным как множество точек, и структурно ограниченными бесконечными последовательностями интервалов.

Любопытное аналитическое доказательство приводит Леви \cite{supremum}, позаимствованное из книги "Функция действительного переменного" Никола Бурбаки \cite{fonctions}. На первый взгляд оно использует лишь понятие точной верхней грани и из предположения, что множество действительных чисел счётно выводит противоречие, заключающееся в том, что некоторая точка $y + f(c)$ одновременно принадлежит множеству $A$ и превосходит точную верхнюю грань $c$ этого множества. В статье ''Несчётность множества континуум, ещё шах'' мы разобрали первое доказательство несчётности вещественных чисел Кантора. Мы показали, что без дополнительных предположений построение просто фиксирует парадокс, схожий с апориями Зенона. Дополнительные предположения о структуре бесконечности, способны как обеспечить неконсистентность данного доказательства, так атрибутировать конфликт предикативности и импредикативности разным уровням воспроизводства аддитивной группы. Ровно тоже самое происходит и в данном доказательстве, противоречие возникает вследствие того, что $y$ и $c$ не могут быть двумя смежными точками, поскольку сечения типа "скачок" запрещены предикативным порядком требуемым аксиомой непрерывности. В общем случае, построение парадоксально, но мы предполагая структуру бесконечности приходим к выводу, что конфликт замыкания упорядоченной группы и аксиомы непрерывности может быть описан через экспоненциальное различие сложности. Если мы рассмотрим то же самое определение $A$ на множестве рациональных чисел, предположим, что в $A$ существует максимальный элемент $c$, взяв $y = c - f(c)/2$, повторив рассуждения касающиеся точной верхней грани множества $A$ для точки $c$ приходим к противоречию, предположение, что дополнение к $A$ содержит минимальный элемент $c$, приводит к противоречию аналогичным образом. То есть, наше построение определяет множество $A$, как множество образованное сечением множества рациональных чисел типа ''пробел'', однако мы делаем из этого вывод о существовании другого вида чисел. Противоречие возникает из нашего стремление сделать континуум предикативным, в достаточной степени для определения непротиворечивого предельного аппарата, при этом коллизия оказывается структурно подобной экспоненциальному выходу во внешний контур. 

Таким образом, на сегодняшний день не существует ни одного доказательства несчётности континуума в смысле ёмкости. Все известные доказательства устанавливают лишь тот факт, что множество порядкового типа $2^{\omega}$ (не то же самое, что ординал $2^{\omega}$) невозможно переупорядочить изоморфно порядковому типу $\omega$. Различные доказательства содержат схожие неявные допущения, делающие их выводы контекстно-зависимыми и неоднозначными при попытке универсального применения. Тем не менее, проведённый анализ позволяет сформулировать следующий вывод: континуум допускает различные способы описания и представления. По-видимому, его сущность выходит за рамки предмета теории множеств. Иными словами, континуум не является множеством в строгом смысле этого слова - вывод, непосредственно следующий из первой доказанной Кантором теоремы. Следствием этого является принципиальная несводимость интегрального и дифференциального исчислений к суммам точечных элементов. В то же время теория множеств предоставляет мощный и точный язык для описания и исследования континуума. Хотя континуум сам по себе не является множеством, мы можем аннотировать его как множество с требуемой степенью детализации. Для этого в нашем распоряжении имеются, по крайней мере, три различные алгебры: кардинальная, ординальная (порядковая) и индексная. Рассмотрим некоторые из таких представлений. В статье ''Несчётность множества континуума, шах'' мы построили два множества вещественных чисел $C[0, 1]$ и $U[0, 1]$. По существу, $U[0, 1]$ представляет собой множество всех вещественных чисел в двоичной позиционной записи (теоремы 2 и 3) на отрезке $[0, 1]$. Легко видеть, что данное множество счётно по ёмкости и несчётно по мощности. При этом в кардинальной алгебре существует неупорядоченное взаимно однозначное соответствие между множеством позиций двоичной записи и множеством $U[0, 1]$. В порядковой алгебре его порядковый тип - $2^{\omega}$, который в рамках этой алгебры не допускает изоморфного отображения на порядковый тип $\omega$ множества позиций. В альтернативной ординальной алгебре его порядковый тип $\omega_1 = \aleph_1$. В индексной алгебре тип данного множества записывается как $2^{\dot{\omega}}$. Но что представляет из себя множество $C[0, 1]$? В рамках современных теорий это множество вещественных чисел с конечной позиционной записью, сколь угодно большой но конечной. Но проблема в том, что множество конечных позиций бесконечно, фактически это множество $\omega$, все варианты комбинаций этих позиций по смыслу представляют собой множество $U[0, 1]$. Аксиома выделения здесь нам не сильно помогает, да мы можем определить признак конечности, но если выбрать всё конечные числа из $\omega$ мы получим $\omega$, точно также множество задействованных позиций в множестве $C[0, 1]$ есть ничто иное как $\omega$. Тем не менее, интерпретация различающая синтаксическую редукцию выражения на уровне аддитивной структуры и синтаксическую редукцию того же выражения на первом уровне мультипликативного воспроизводства аддитивной структуры закрепляется в современных теориях, поскольку нам нужно каким-то образом обосновывать почему множество $U[0, 1]$ становится несчётным, в то время как множество $C[0, 1]$ является счётным. Несмотря на то, что в работах Кантора мы не найдём даже намёка на подобную интерпретацию. Кантор предлагает рассматривать $C[0, 1]$ как объект кардинальной алгебры, где отсутствие структуры бесконечности $\aleph_0$ воспроизводит классическую замкнутую арифметику, в которой совокупность натуральных чисел и совокупность натуральных степеней двойки сосуществуют без противоречия, поскольку отсутствие бесконечной структуры не допускает ни одного из построений способных разделить данные совокупности. Данное множеств внешне импредикативно, то есть всего-навсего не позволяет добавлять элементы в конец. В то время как $U[0, 1]$ - это объект порядковой алгебры, где множество разрядов $\omega$ вписано в универсальную структуру трансфинитных объектов, что позволяет чётко разнести порядковый тип множества строк $2^{\omega}$ и разрядов. Такое множеств предикативно внешне, и мы свободно можем продолжить его трансфинитное построение. Кантор действительно подразумевает, что мы можем вполне упорядочить любое множество, но такое упорядочение - это по сути аннотация множества, которая сама по себе уже является объектом другой алгебры. Иными словами, упорядочение - это морфизм переводящий объект одной алгебры в объект другой. Именно поэтому аксиомы выбора привносить столько странностей, ведь данная аксиома связывает две различных алгебраических структуры, которые мы пытаемся объединить в одну в рамках универсальной теории множеств. Если рассмотреть, например, свойства $p$-адических чисел, то не сложно видеть, что они сохраняют свои свойства при переходе из ординальной алгебры в порядковую, свойства изменяются в кардинальной алгебре, а в индексной алгебре они допускают естественное введение тотального порядка.

В принципе для описания континуума достаточно и кардинального представления с его бесструктурной бесконечностью. Рассмотрим, например, ранние представления Кантора о вещественных числах, как о классах эквивалентности фундаментальных последовательностей, где для сходимости используется критерий Коши \cite{cantor-real}.
\begin{align}
\begin{split}
\forall \varepsilon > 0 \; \exists N(\varepsilon): \forall n > N(\varepsilon) \; \forall m > 0 \; |a_{n+m} - a_n| < \varepsilon
\end{split}
\end{align}

Откуда непосредственно следуют правила сравнение вещественных чисел, в частности, две фундаментальные последовательность $\{a_n\}$ и $\{a_n\}$ определяют одно и тоже вещественное число тогда и только тогда, когда
\begin{align}
\begin{split}
\forall \varepsilon > 0 \; \exists N(\varepsilon): \forall n > N(\varepsilon) \; |a_n - b_n| < \varepsilon
\end{split}
\end{align}

Используя такой подход, мы можем получить замкнутую классическую модель континуума. Рассмотрим множество $S[0, 1]$ фундаментальных последовательностей, построенное схожим образом с множествами $U[0, 1]$ и $C[0, 1]$. Начинаем с последовательности $\{s^0_n\}$, $s^0_n = 0.0$, $\forall n \in \mathbb{N}$, добавляем его во множество, получаем $S_0[0, 1]$, состоящее из 1 элемента. Следующая последовательность $\{s^1_n\}$, $s^1_1 = 0.0$, $s^1_n = 0.1$, $\forall n \in \mathbb{N} \wedge n > 1$, добавляем её в наше множество, теперь $S_1[0, 1]$ состоит из двух элементов. На следующей итерации снова строим последовательности задействующие новый разряд вещественных чисел, имеем $\{s^2_n\}$, $s^2_1 = 0.0$, $s^2_2 = 0.0$, $s^2_n = 0.01$, $\forall n \in \mathbb{N} \wedge n > 2$ и $\{s^3_n\}$, $s^3_1 = 0.0$, $s^3_2 = 0.1$, $s^3_n = 0.11$, $\forall n \in \mathbb{N} \wedge n > 2$, включаем их во множество, получаем 4 элемента в $S_2[0, 1]$. Тоже самое делаем на $i$-той итерации, получаем $S_i[0, 1]$ из $2^i$ элементов и т.д. Следуя этой процедуре неограниченно долго получаем множество $S[0, 1] = \bigcup\limits_{i \in \mathbb{N}} S_i[0, 1]$:
\begin{center}
\begin{tabular}{c c c r r r r r r r c} 
$s^x$ & $=$ & $\{$ & $s^x_1$, & $s^x_2$, & $s^x_3$, & $s^x_4$, & $\dots$, & $s^x_n$, & $\dots$ & $\}$\\
$s^0$ & $=$  & $\{$ & 0.0, & 0.0, & 0.0, & 0.0, & $\dots$, & 0.0, & $\dots$ & $\}$\\
$s^1$ & $=$  & $\{$ & 0.0, & 0.1, & 0.1, & 0.1, & $\dots$, & 0.1, & $\dots$ & $\}$\\
$s^2$ & $=$  & $\{$ & 0.0, & 0.0, & 0.01, & 0.01, & $\dots$, & 0.01, & $\dots$ & $\}$\\
$s^3$ & $=$  & $\{$ & 0.0, & 0.1, & 0.11, & 0.11, & $\dots$, & 0.11, & $\dots$ & $\}$\\
$s^4$ & $=$  & $\{$ & 0.0, & 0.0, & 0.0, & 0.001, & $\dots$, & 0.001, & $\dots$ & $\}$\\
$s^5$ & $=$  & $\{$ & 0.0, & 0.1, & 0.1, & 0.101, & $\dots$, & 0.101, & $\dots$ & $\}$\\
$s^6$ & $=$  & $\{$ & 0.0, & 0.0, & 0.01, & 0.011, & $\dots$, & 0.011, & $\dots$ & $\}$\\
$s^7$ & $=$  & $\{$ & 0.0, & 0.1, & 0.11, & 0.111, & $\dots$, & 0.111, & $\dots$ & $\}$\\
$\dots$ & $=$  & $\{$ &  &  &  & \dots &  &  &  & $\}$
\end{tabular}
\end{center}

Все последовательности во множестве $S[0, 1]$ фундаментальны, в силу того, что предел каждой такой последовательности равен последнему элементу повторяемому неограниченное количество раз начиная с элемента, номер которого определяется двоичным логарифмом номера последовательности. 

Теорема 10: {\sl Для любой фундаментальной последовательности рациональных чисел $\{r_n\}$, определяющей вещественное число $r$, на отрезке $[0,1]$, во множестве $S[0, 1]$ найдётся фундаментальная последовательность рациональных чисел $\{s_n\}$, определяющая то же самое число $r$}

Развернув определения, можно сформулировать утверждение теоремы как: для любой фундаментальной последовательности рациональных чисел $\{r_n\}$ на отрезке $[0, 1]$, в построенном множестве $S[0, 1]$ существует фундаментальная последовательность рациональных чисел $\{s_n\}$, для которой для всякой сколь угодно маленькой величины $\varepsilon > 0$ существует порядковый номер такой, что для всякого $n > N(\varepsilon)$ выполняется неравенство $|r_n - s_n| < \varepsilon$.
\begin{align}
\begin{split}
\forall \{r_n\} \in [0, 1] \; \exists \{s_n\} \in S[0, 1]: \forall \varepsilon > 0 \; \exists N(\varepsilon): \forall n > N(\varepsilon) \; |r_n - s_n| < \varepsilon
\end{split}
\end{align}

Доказательство: 

Предположим противное, - существует некоторая фундаментальная последовательность $\{r^0_n\}$ на отрезке $[0, 1]$, обладающая тем интересным свойством, что для любой последовательности $\{s_n\} \in S[0, 1]$ можно указать $\varepsilon > 0$ такое, что существует некоторые порядковый номер $N(\varepsilon)$, начиная с которого все элементы этой последовательности с большими номерами $n$ удовлетворяют неравенству $|r_n - s_n| \geq \varepsilon$. Будем последовательно исключать из рассмотрения все элементы множества $S[0, 1]$, для которых $\varepsilon = 0.1, 0.01, 0.001, \dots$. Возможны всего два случая: 

\begin{enumerate}
    \item Нам никогда не удастся исключить всех последовательностей. В таком случае все оставшиеся фундаментальные последовательности определяют вещественное число $r$. Действительно, пусть $\{s^*_n\}$ какая-то из оставшихся последовательностей, предположим, что она не определяет число $r$. Тогда существует некоторая величина $\varepsilon > 0$ такая, что начиная с некоторого порядкового номера $N(\varepsilon)$ все элементы последовательностей $\{s^*_n\}$ и $\{r^0_n\}$ с номерами $n$ большими $N(\varepsilon)$ удовлетворяют неравенству $|r^0_n - s^*_n| \geq \varepsilon$. Однако, согласно принципу Архимеда \cite{matan-archimed}, каким бы большим ни было число $1/\varepsilon$, всегда найдётся число $k > 1/\varepsilon$, следовательно последовательность $\{s^*_n\}$ была бы исключена на шаге $k$, поскольку $2^{k} > k \Rightarrow \varepsilon < 1/k < 2^{-k}$, значит последовательность $\{s^*_n\}$ определяет $r$, что противоречит предположению.
    \item Мы исключим все последовательности на некотором конечном шаге $k$. В этом случае, величина $\varepsilon = 2^{-k}$ может быть указана абсолютно для любой последовательности множества $S[0, 1]$ и начиная с некоторого номера $N(2^{-k})$ элементы такой последовательности $\{s^x_n\}$ и элементы последовательности $\{r^0_n\}$ с номерами $n > N(2^{-k})$ будут удовлетворять неравенству $|r^0_n - s^x_n| \geq 2^{-k}$. Последовательность  $\{r^0_n\}$ фундаментальна, значит существует некоторый порядковый номер $N(2^{- (k + 1)})$ такой, что $\forall n > N(2^{- (k + 1)}) \; \forall m > 0 \; |r_{n + m} - r_n| < 2^{- (k + 1)}$. Из номеров $k + 1$ и $N(2^{- (k + 1)})$ выберем максимальный и обозначим его за $j$. Рассмотрим $S_j[0, 1]$, которое по построению содержит все возможные последовательности, которые определяют все вещественные числа с количеством значащих позиций не превосходящих $j$. Из всех последовательностей $S_j[0, 1]$ выберем одну, $j$-ый элемент которой совпадает с $j$-ым элементом последовательности $\{r^0_n\}$ во всех позициях до $j$ включительно, обозначим её как $\{s^*_n\}$, пусть $r^0_j - s^*_j = \delta$, тогда в силу $j \geq k + 1$ имеем $|\delta| < 2^{- (k + 1)}$. Заметим также, что по построению, начиная с номера определяемого двоичным логарифмом номера последовательности, элементы всех последовательностей повторяются, то есть во множестве $S_j[0, 1]$ для всех последовательностей справедливо: $s^i_i = s^i_j, \forall i > j$. Возьмём любой порядковый номер превосходящей $j + 1$ и $N(2^{-k}) + 1$, где $N(2^{-k})$ - тот самый порядковый номер, после которого все элементы всех последовательностей множества $S[0, 1]$ различаются с соответствующим элементом последовательности $\{r^0_n\}$ не менее чем на $2^{-k}$, запишем такой номер как $j + m$. Рассмотрим разность $r^0_{j + m} - s^*_{j + m}$, заметим что $s^*_{j + m} = s^*_j$, так как $j + m > j$, а $s^*_j = r^0_j - \delta$. Значит $|r^0_{j + m} - s^*_{j + m}| = |r^0_{j + m} - r^0_j + \delta| \geq 2^{-k}$, откуда раскрыв модуль получаем $|r^0_{j + m} - r^0_j| \geq 2^{-k} \pm \delta > 2^{-k} - 2^{- (k + 1)}$, то есть $|r^0_{j + m} - r^0_j| > 2^{- (k + 1)}$ и последовательность $\{r^0_n\}$ не может быть фундаментальной, ведь по определению $j > N(2^{- (k + 1)})$. Следовательно, данный случай также приводит к противоречию со сделанным предположением.
\end{enumerate}

Таким образом, множество $S[0, 1]$ действительно определяет все вещественные числа на отрезке $[0, 1]$. Теорема доказана. 

Множество $S[0, 1]$ является великолепным инструментом для математически строгого описания множества вещественных чисел в полностью замкнутой кардинальной алгебре, где мы можем гарантировать что исключение всех элементов никогда не произойдёт. Например, мы можем наблюдать, почему диагональный метод Кантора здесь не работает. Действительно, диагональный элемент на множестве $S[0, 1]$ является последовательностью, обозначим её $\{d_n\}$. Первый элемент её $d_1 = 0.d^1$, где $d^1$ цифра отличная от соответствующей цифры первого элемента первой последовательности, второй элемент $d_2 = 0.d^1d^2$ содержит ещё одну цифру отличную от второй цифры второго элемента второй последовательности, на шаге $n$ имеем элемент $d_n = 0.d^1d^2{\dots}d^n$, и так далее. Легко понять, что такая последовательность фундаментальна, поскольку каждый следующий элемент может отличаться от предыдущего только в следующем разряде, что делает её монотонно возрастающей, а так как она ограничена отрезком $[0, 1]$, то и имеющей предел, в силу теоремы Вейерштрасса о пределе монотонной последовательности \cite{matan-weierstrass}, обозначим её предел как $d = \lim_{n\to\infty} d_n$. Также легко видеть, что данный предел не может принадлежать последовательности $\{d_n\}$, ведь если он имеет некоторый номер, скажем $k$, то есть $d = d_k$, то именно столько значащих цифр имеет его двоичная запись, но это означает, что все последующие элементы также равны $d$, так как в противном случае для разности с первым отличающимся элементом не получится указать номер, начиная с которого все элементы попадают в окрестность данной разности, ведь последовательность монотонно возрастает, но в таком случае $S_k[0, 1]$ содержит последовательность идентичную $\{d_n\}$, чего быть не может, ввиду того, что их диагональные элементы должны различаться по построению. То есть, каким бы ни было вещественное число $d$ мы видим чередование его построения и достройки $S[0, 1]$ до размера содержащего построенный сегмент. Поскольку в кардинальной алгебре рефлексивности не ограничена экспоненциально две этих перемежающихся последовательности имеют один и тот же предел $\aleph_0$ и для всякого диагонального элемента в множестве $S[0, 1]$ уже есть эквивалентная последовательность.

Что можно также продемонстрировать следующим способом. Рассмотрим, подпоследовательность $S^d[0, 1]$ последовательности $S[0, 1]$, такую что всякая последовательность $s^d_n$ начинается с $n$ первых элементов последовательности $\{d_n\}$, после чего элемент $d^n$ повторяется неограниченное количество раз, то есть $s^d_n \in S_n[0, 1]$. По построению последовательность диагональных элементов подпоследовательности $S^d[0, 1]$ является ничем иным как диагональной последовательностью $\{d_n\}$, при этом каждому элементу $d^n$ во взаимно однозначное соответствие поставлена последовательность $s^d_n$, элементы которой совпадают с элементами $\{d_n\}$ вплоть до элемента $d^n$, иными словами номер элемента диагональной последовательности оказывается накрепко связан с номером последовательности из множества $S^d[0, 1]$. Следовательно, последовательность $S^d[0, 1]$ содержит последовательность, предел которой равен $d$, так как она содержит все последовательности совпадающие с $\{d_n\}$ вплоть до любого элемента с натуральным номером. Тогда, вопрос о номере этой последовательности во множестве $S^d[0, 1]$ равносилен вопросу о самом близком к пределу $d$ элементе последовательности $\{d_n\}$, который не имеет смысла поскольку мы устремили этот номер к бесконечности и тем самым устремили к бесконечности номер соответствующей последовательности во множестве $S^d[0, 1]$. Мы ведь не можем одновременно устремить номер к бесконечности и назвать его, точнее чтобы это сделать, нам требуется морфология бесконечности, которая недоступна. Раскрывая понятие предела, мы обнаруживаем, что для любого сколь угодно маленького $\varepsilon > 0$ существует такой натуральный номер $N(\varepsilon)$, что для всех $n \geq N(\varepsilon)$ справедливо неравенство $|s^d_{n, n} - r| < \varepsilon$, где $s^d_{n, n}$ - $n$-ый элемент последовательности $s^d_n$. Таким образом, связь оказывается настолько прочной, что разорвать её не лишившись понятия предела представляется крайне сложной задачей.

В кардинальной алгебре мы не можем зафиксировать множество натуральных чисел как внешне предикативный объект: существует своего рода принципиальное ограничение, не позволяющее вполне упорядочить все натуральные числа в единую последовательность, так чтобы нельзя было добавить следующее натуральное число. Именно это явление мы наблюдаем в новых основаниях Куайна, где отсутствие аксиомы выбора приводит к тому, что некоторые множества не допускают упорядочения. Аналогичная ситуация возникала и в предыдущих работах авторов, когда предположение об гомогенности бесконечности влекло за собой невозможность диагонального построения. Однако стоит перейти к внешней точке зрения, в которой бесконечности становятся сравнимыми между собой, как сразу обнаруживается, что линейная и экспоненциальная бесконечности оказываются несопоставимыми. Как было установлено ранее в порядковой алгебре Кантора, соотношение $\nu < 2^{\nu}$ выступает универсальным инвариантом в упорядоченном универсуме. Современная интерпретация $2^{\omega}$ как множества всех функций $f : \omega \to \{0, 1\}$ с конечным носителем представляется не только интуитивно малопонятной, но и формально трудноопределимой. Более того, она невозможна в рамках оригинальной теории Кантора. Позиционная конструкция обеспечивает качественно большее разнообразие, нежели линейное следование элементов, и это различие приобретает статус фундаментального закона. В качестве иллюстрации мощи и гибкости исходного подхода Кантора и проблемности стандартной интерпретации приведём ещё одну конструкцию. Рассмотрим множество $C[0, 1]$ всех вещественных чисел, позиционная запись которых является конечной. Ясно, что этим множеством мы пытаемся задать кардинальное (внутреннее) представление вещественных чисел, которое во внешней алгебре соответствовало бы чему-то вроде $2^{\omega - 1}$. Поскольку $\omega - 1 = \omega$, это выражение сливается с $2^{\omega}$, которое связано с бесконечной позиционной записью. Поэтому $C[0, 1]$ в теории Кантора можно ассоциировать только с кардинальной алгеброй $2^{\aleph_0}$. Отметим также, что в современной интерпретации множеству $C[0, 1]$ мы сопоставляем именно $2^{\omega}$. 

Лемма 1: {\sl Дополнение множества $\{1\} \cup C[0, 1]$ до множества всех вещественных чисел на отрезке $[0, 1]$ не превосходит по ёмкости $C[0, 1]$}

Доказательство: 

Рассмотрим множество $L[0, 1] = [0, 1] \setminus \{1\} \cup C[0, 1]$. Мы исключили из отрезка $[0, 1]$ счётное множество $\{1\} \cup C[0,1]$, в результате чего осталось бесконечное семейство попарно непересекающихся интервалов. С точки зрения современных представлений каждое из этих множеств является вырожденным отрезком вида $[x, x]$, то есть состоит из единственной точки. Покажем, что это действительно так. Предположим, напротив, что некоторый из оставшихся интервалов содержит по крайней мере две различные точки $a$ и $b$, пусть для определённости $a < b$. Согласно теореме 2 из работы ''Несчётность множества континуум, шах'', для любого вещественного числа $r \in [0, 1]$ множество $C[0, 1]$ содержит последовательность $\{c^r_n\}$, являющуюся нижним двоичным приближением числа $r$. Рассмотрим последовательность $\{c^b_n\}$ нижних двоичных приближений числа $b$, её предел равен $b$ \cite{matan-bottom-limit}. Следовательно в любой окрестности точки $b$ содержится бесконечно много элементов этой последовательности. Выберем $\varepsilon = b - a > 0$. Тогда в окрестности $(b - \varepsilon, b + \varepsilon)$ должно лежать бесконечно много членов последовательности $\{c^b_n\}$. Однако $b - \varepsilon = a$, поэтому указанная окрестность содержит точку $a$. Если бы в ней находились элементы $\{c^b_n\}$, то эти точки разделяли бы интервал между $a$ и $b$, что противоречило бы предположению о том, что $a$ и $b$ принадлежат одному и тому же оставшемуся интервалу. Полученное противоречие показывает, что каждый из оставшихся интервалов содержит не более одной точки. Далее рассмотрим индуктивное построение данного семейства интервалов, акцентируя внимание на соответствии каждого интервала его левой границе. Зададим функцию $f$, которая каждой точке $x \in C[0, 1]$ сопоставляет правый интервал: $f(0.00) = (0, 1)$, $f(0.10) = (0.1, 1)$, $f(0.01) = (0.01, 0.10)$, $f(0.11) = (0.11, 1)$ и так далее, то есть $f(x) = (x, x + 2^{ -\min \{ \; i \; \in \; \mathbb{N} \; \mid \; 2^{i}x \; \in \; \mathbb{N} \; \} })$. Таким образом, множество $C[0, 1]$ преобразуется во множество соответствующих интервалов. Определим теперь функцию $g(x) = \{ y \mid y \in f(x) \wedge \forall z \in C[0, 1] \; (z > x \Rightarrow y < z) \}$, которая каждой точке $x \in C[0, 1]$ ставит в соответствие её непосредственную правую окрестность (первый интервал справа от $x$, не содержащий более левых точек из $C[0, 1]$). Можно сразу записать функцию $h(x) = \{ y \mid y \in [0, 1] \wedge \forall z \in C[0, 1] \; (z > x \Rightarrow y < z) \}$, но подход с двумя функциями допускающий рекурсивную интерпретацию представляется более наглядным. Ёмкость множества таких окрестностей не превосходит ёмкости $C[0, 1]$. Поскольку оба описания - множество оставшихся интервалов и множество правых окрестностей - представляют один и тот же процесс разбиения, ёмкости этих двух семейств совпадают. Следовательно, ёмкость множества $L[0, 1]$ не превосходит ёмкости множества $C[0, 1]$. Лемма доказана.

Хотя функция $f$ и производная от неё функция $g$ существуют в рамках современных аксиоматик как объекты теории, между построенным таким образом множеством значений $g$ и множеством компонент дополнения $L[0, 1]$ нет биекции. Функция $g$ сопоставляет каждой точке $x \in C[0, 1]$ некоторую правую окрестность, однако в силу непрерывности линейного порядка на $\mathbb{R}$ эти окрестности на каждом конечном этапе остаются нетривиальными интервалами. Вырожденные интервалы возникают исключительно в пределе после завершения всех счётных шагов, причём подавляющее большинство из них не попадает в образ $g$. Таким образом, множество значений функции $g$ и множество вырожденных интервалов, составляющих $L[0, 1]$, - это два принципиально разных объекта. По существу, в данном построении $\omega$ трактуется как класс эквивалентности конечных ординалов (монада с внутренней структурой), иными словами, мы разделяем конечную арифметику Пеано внутри $\omega$ и само это множество, то есть используем разрыв предикативности. Одновременно формально рассматриваем $\omega$ как множество, состоящее исключительно из конечных множеств и интерпретационно допускаем наличие у него дополнительной внутренней структуры (сингулярности), свойства которой реализуют предельный переход. На каждом конечном шаге никакого образования интервалов не связанных с нашей функцией не происходит, однако в пределе они появляются. В рассматриваемом построении мы сделали тождество количества интервалов и количества точек разбиения инвариантом, справедливым для любого конечного случая. Однако при переходе к $\omega$ мы вынуждены переопределять принцип индукции, трактуя $\omega$ как точку разрыва предикативного описания, нуждающуюся в особо обращении, чтобы избежать противоречия. При этом данная интерпретация неявная, она не является частью теории, вынуждая нас снова принимать нечто вне теории, чтобы обосновать теорию. Доказательство строится в два независимых этапа: сначала устанавливается, что каждый оставшийся интервал содержит не более одной точки, а затем - что семейство таких интервалов счётно. Объединить эти два свойства в единую непротиворечивую картину не удаётся из-за непрерывности вещественной прямой. Данное наблюдение позволяет сформулировать теорему в рамках стандартной интерпретации аксиоматических теорий множеств.

Теорема 11: {\sl Не существует счётного множества $C[0, 1]$ вещественных чисел с конечной записью на отрезке $[0, 1]$}

Доказательство:

Предположим противное, что существует некоторое счётное множество $C[0, 1]$ вещественных чисел на отрезке $[0, 1]$, имеющих сколь угодно большую, но конечную позиционную запись. В контексте современных интерпретаций, это множество порядкового типа $2^{\omega}$ булевых функций с конечным носителем. Преобразуем данное множество функцией $h(x) = \{ y \mid y \in [0, 1] \wedge \forall z \in C[0, 1] \; (z > x \Rightarrow y < z) \}$, ставящей элементу множества $C[0, 1]$ его непосредственную правую окрестность не пересекающуюся ни с каким иным элементом $C[0, 1]$. Как показано в лемме 1, все такие окрестности либо пусты, либо имеют ровно одну точку. Это означает, что некоторому элементу $c_k \in C[0, 1]$ соответствует интервал $h(c_k)$ содержащий точку $\pi/4$, что противоречит аксиоме непрерывности. Следовательно наше предположение не верно. Теорема доказана. 

Коллизия в данном случае возникает именно из-за того, что в рамках современных представлений мы считаем множество $C[0, 1]$ порядкового типа $2^{\omega}$ счётным. Технически данное доказательство очень близко доказательству Кантора о несчётности множества вещественных чисел. В обоих случаях конфликт происходит с аксиомой непрерывности, которая требует экспоненциального скачка. Иными словами, неявное предположение современных теорий о допустимости распространения внутренней индукции на ординальное возведение в степень формально приводит к невозможности использовать ординальную арифметику в анализе. Мы не можем заключить из этого что данная интерпретация теории множеств противоречива, поскольку противоречие возникает из-за вспомогательной аксиомы анализа, а не внутри самой теории. Тем не менее, комбинация ординальной алгебры и анализа оказывается проблематичной. Современная интерпретация для объяснения данного феномена привлекает какой-то невообразимый арсенал дополнительных предположений, в конечном итоге сводящихся к прыжку веры. К счастью, конструкции с конечным носителем используются довольно редко, что сохраняет конструктивность большей части доказательств анализа в контексте современной интерпретации. В противоположность этому, интерпретация Кантора лишена указанной контринтуитивности. Множество $C[0, 1]$ существует здесь исключительно как кардинальное описание; следовательно, мы должны использовать соответствующее описание множества вещественных чисел, которое в кардинальной алгебре и есть $C[0, 1]$. Конфликта с аксиомой непрерывности не происходит, поскольку все интервалы в $L[0, 1]$ оказываются пустыми. Если же мы переходим в порядковую (ординальную) алгебру, где все вещественные числа отрезка $[0, 1]$ образуют множество $U[0,1]$, то, фиксируя множество всех разрядов как $\dot{\omega}$, получаем, что множество $C[0, 1]$ (вещественные числа, в которых задействованы не все разряды) имеет индексный тип $2^{\dot{\omega} - 1}$. Аналогичный индексный тип имеет и его дополнение $L[0, 1]$. При этом порядковый тип обоих множеств равен $2^{\omega}$, и ни одно из них не допускает изоморфного переупорядочения на $\omega$, как и сам отрезок. В результате, в кардинальной алгебре становится справедливой следующая теорема.

Теорема 12: {\sl Множество вещественных чисел на отрезке $[0, 1]$ имеет ёмкость $C[0, 1]$}

Доказательство:

Следует непосредственно из леммы 1, так как $[0, 1] = C[0, 1] \bigcup L[0, 1]$. Теорема доказана. 

Таким образом, современная интерпретации теории множеств сопряжена с рядом существенных недостатков. К недостаткам формального характера относятся архитектурные дефекты связанные с аксиомой бесконечности, нарушающей формальные требования к аксиомам. Онтологические - связаны с распространённой в математике ошибкой реификации - отождествление описания объекта с самим объектом. Существенную роль также играют методологические формальные ошибки преждевременного обобщения внутренних функций и индуктивной морфологии бесконечности как единственных в своём роде. Кроме того, отсутствует явная формулировка предположения о допустимости распространять внутреннюю кардинальную индукцию на ординальное возведение в степень. Необходимо подчеркнуть, что, выступая в настоящей работе в роли критиков общепринятой картины, авторы признают её продуктивность и концептуальную ясность. Проблема заключается не столько в наличии ошибок, которые невозможно полностью избежать, а идеологизации мейнстрима и нескрываемой цензуре альтернативных подходов. Между тем эти подходы открывают, более дифференцированную интерпретации, которая позволяет непротиворечиво объединить все накопленные факты в целостную систему. Авторы не ставят целью опровержение современной теории множеств. Их задача скромнее - показать, сколь небольшую часть многогранной природы множеств удалось охватить доминирующей ныне аксиоматической картиной. Многие выдающиеся математики, внёсшие фундаментальный вклад в развитие теории множеств (в частности, Г. Кантор, А. Френкель, К. Гёдель и другие), осознавали фундаментальную роль импредикативности в математике и видели необходимость в её проявлениях подобии и рефлексивности как свойствах отличающих конечное и бесконечное. Именно поэтому уже в формулировке аксиомы бесконечности неизбежно присутствует определённый порочный круг: иного способа ввести это свойство в формальную теорию на тот момент не существовало. Наша задача состоит лишь в том, чтобы устранить связанные с этим противоречия. Разделяя понятия подобия алгебраических структур и рефлексивности, мы не подрываем оснований теории множеств. Аксиоматические теории задают в первую очередь определённые алгебраические структуры, внутри которых внешнее подобие может восприниматься как внутренняя рефлексивность. Мы лишь демонстрируем, что такая структура не является единственно возможной: существует, по крайней мере, три легитимные алгебры, описывающие разные аспекты сущностей, называемых множествами. Современная теория множеств стремится свести все три алгебры к одной универсальной теории множеств. В отличие от этого, оригинальные работы Кантора оставляли вопрос об интерпретации открытым. Альтернативная интерпретация предполагает сохранение независимости трёх алгебр: кардинальной (с естественной формализацией в системах типа New Foundations), ординальной (формализуемой в ZFC и NBG) и внешней индексной алгебры (связанной с нестандартными натуральными числами, сюрреальными числами и гросс-объектами). Конфликт между современной и альтернативной интерпретациями наиболее ярко проявляется в теореме Кантора и, в особенности, в вопросе о границах её применимости. Теорема остаётся справедливой в обеих картинах, однако в одной из них она носит универсальный характер, тогда как в другой применима лишь к ординальной и индексной алгебрам и не доказуема в кардинальной. 

Здесь со всей очевидностью раскрывается сложность концептуального описания бесконечности. Современная математика поспешила отказаться от классического разделения на потенциальную и актуальную бесконечность, хотя именно это различение хорошо отражает разницу между неупорядоченным (кардинальным) и упорядоченным (ординальным) описаниями, между импредикативность и предикативностью. Критика современных теорий часто остаётся неэффективной именно потому, что язык современной математики изначально построен под общепринятую картину, делая понятие потенциальной бесконечности архаичным. Между тем, для Кантора это была одна из реальных форм бесконечности. Для адекватного описания возникающей проблематики необходим новый концептуальный язык. Именно такую попытку предпринимают авторы, акцентируя внимание на различных аспектах бесконечности в зависимости от используемой алгебры. В кардинальной алгебре бесконечность оказывается чрезвычайно близкой к потенциальной: мы не можем вполне упорядочить кардинально бесконечное множество, поскольку оно представляет собой не множество как таковое, а замкнутую область регулярности, внутренняя структура которой упорядочена, однако целое как таковое не допускает внешнего предикативного определения. В случае потенциальной бесконечности это проявляется как принципиальная незавершённость объекта, представляющего все свои элементы. Формально это связано со структурой отношения принадлежности; в частности, в системах типа NF это выражается в запрете самопринадлежности. В кардинальной арифметике Кантора данное свойство проявляется в коммутативности сложения и умножения, а также в том, что трансфинитные числа выступают неподвижными точками этих операций. В такой алгебре отсутствует морфология, способная поддержать теорему Кантора. Кардинальная алгебра абстрагируется от внутренней структуры множеств. Именно отсутствие глобальной структуры в бесконечности позволяет непротиворечиво замкнуть данную алгебру. Кардинальная алгебра обладает свойством универсальности, всякое множество представимо в кардинальной морфологии и сопоставимо с другими множествами. Но она лишена тотальности, поскольку отсутствует какая-либо структура способная обеспечить такую тотальность. Однако для сравнения объектов между собой требуются внешние структуры. Только помещая объекты в общий контекст, мы можем выявить различия между ними. Так, мера является внешней характеристикой: с точки зрения внутренней структуры два отрезка могут быть идентичны, но различаться по мере, когда они вписаны в общий контекст. Именно это делает Кантор в порядковой арифметике: он вводит внешние операции, создавая внешнее описание множеств, с которым связана актуальная бесконечность, позволяющая рассматривать множество как предикативный внешне объект. В этой картине теорема Кантора становится неопровержимой - как в ординальной, так и в индексной алгебре. Множества разбиваются на уровни, образуя иерархию, которую невозможно замкнуть, что приводит к парадоксу Бурали-Форти. Замыкание возможно лишь во внутренней (кардинальной) алгебре, однако глобальная упорядоченность несовместима с глобальной замкнутостью.

Потенциальная бесконечность оказывается тесно связанной с внутренней структурой множеств, кардинальной арифметикой и формальными теориями, накладывающими полный запрет на самореференцию. Напротив, актуальная бесконечность позволяет проводить сравнение множеств между собой, однако ценой принятия глобальной незамкнутой внешней структуры. Последняя находит своё воплощение в ординальной и индексной алгебрах, где формальные системы пытаются ограничить самореференцию. Особое внимание привлекает яркий концептуальный конфликт в интерпретациях теоремы Кантора. В классической интерпретации ординал $\omega^\omega$ считается счётным и соответствует, в частности, порядковому типу множества всех алгебраических чисел, счётность которого была установлена ещё самим Кантором. Аналогично, $2^\omega$ интерпретируется как порядковый тип множества всех вещественных чисел с конечной позиционной записью и также считается счётным. В рамках ZFC и NBG такие множества существуют как объекты ординальной алгебры благодаря аксиоме выделения; это так называемые наследственно конечные множества над бесконечным базовым множеством. Современные интерпретации не позволяют применить теорему Кантора к подобным объектам, поскольку бесконечная запись рассматривается как нечто принципиально новое, чего не может возникнуть в наследственно конечном множестве, построенном над другим бесконечным множеством. В альтернативной интерпретации множество алгебраических чисел, равно как и множество всех вещественных чисел с конечной позиционной записью, принадлежат кардинальной алгебре. В ней степени многочленов и позиции в позиционной записи не образуют законченной последовательности, что и обеспечивает счётность этих множеств, их замкнутость внутри алгебры и возможность взаимно однозначного соответствия. Однако как только степени многочленов или позиции в записи полностью упорядочиваются в множество $\omega$, мы немедленно получаем множество всех алгебраических чисел (включая многочлены с бесконечными степенями) и множество всех вещественных чисел (включая числа с бесконечной позиционной записью). Поскольку все элементы последовательности $\omega$ могут варьироваться независимо, они порождают полную комбинаторную мощность. В упорядоченной структуре множества всех конечных многочленов и всех чисел с конечной позиционной записью коллапсируют к одним и тем же порядковым типам, поскольку $\omega - 1 = \omega$ и $x^{\omega - 1} = x^\omega$. Теорема Кантора начинает действовать именно в тот момент, когда мы пытаемся сравнивать множества между собой. Для такого сравнения необходима внешняя структура. Как только она задана в виде целостной последовательности, диагональный аргумент становится неумолимым: $\omega < 2^\omega$. Все конечные многочлены и числа с конечной записью существуют только как объекты кардинальной алгебры. Но как только порядок формирует бесконечную последовательность, бесконечная последовательность степеней и позиций порождает бесконечные объекты - полноценные вещественные числа.

Обе интерпретации можно считать легитимными и даже консистентными в рамках своих собственных систем предположений. Классическая даёт единое упорядоченное описание всех множественных структур в единой алгебре с универсальным понятием числа, воплощая идеи Готлоба Фреге, тогда как альтернативная позволяет согласовать три независимые алгебры в едином концептуальном пространстве, где числа являются эмерджентными инвариантами структуры. Авторы оставляют открытым вопрос о том, насколько правомерно игнорировать теорему Кантора по отношению к множествам наследственно конечных множеств. Это, по-видимому, единственный серьёзный конфликт между двумя картинами. При этом следует отметить, что принятие аксиомы бесконечности в форме Сколема даёт более прозрачный аппарат, для построения ординальной арифметики Кантора и алгебры омнифических целых чисел. В альтернативной интерпретации три независимых структуры обладают различным пропорциями импредикативности и предикативности поддерживая друг друга. Классическая же интерпретация такой поддержки лишена. Кроме того, в классической интерпретации отсутствует строгое доказательство того, почему теорема Кантора не применима к множеству всех вещественных чисел с конечной позиционной записью, что вызывает проблемы в контексте анализа. Напротив, при рассмотрении упорядоченной структуры видно, что множество всех вещественных чисел с конечной записью для каждого натурального числа уже содержит все возможные комбинации позиций вплоть до минимальной трансфинитной. Этого достаточно для применения диагонального аргумента - достаточно лишь предположить, что все натуральные числа упорядочены в единую целостную последовательность. При построении диагонального элемента мы никогда не выйдем за пределы множества чисел с конечной записью. Следовательно, если множество всех конечных записей задано, диагональный элемент будет противоречить возможности взаимно однозначного соответствия. В кардинальной алгебре для существования взаимно однозначного соответствия требуется радикально устранить возможность самореференции, что позволяет сохранить непротиворечивость. В современных же аксиоматических системах эта уязвимость остаётся незакрытой, допуская самореференцию и потенциальные парадоксы.

Восстановив оригинальную онтологию Кантора, мы можем доказать следующую сильную теорему: не существует алгебры множеств (арифметики, как её называл сам Кантор), которая была бы одновременно замкнутой и допускала тотальный порядок всех своих объектов.

Теорема 13:  {\sl не существует множественной алгебры, удовлетворяющей одновременно следующим трём условиям:
\begin{enumerate}
    \item любое множество может быть аннотировано порядковым типом;
    \item существует универсальный тип множеств $\aleph$, для которого $\mathcal{P}(\aleph) \subseteq \aleph$;
    \item алгебра непротиворечива.
\end{enumerate}}

Доказательство:
Предположим, что такая алгебра существует. Тогда, в силу условия 2, имеется универсальный тип $\aleph$. Пользуясь возможностью аннотации порядковым типом (условие 1), мы можем построить трансфинитный диагональный элемент, отличный от всех элементов $\mathcal{P}(\aleph)$. То есть, доказать теорему Кантора $|\mathcal{P}(\aleph)| > |\aleph|$, из которой следует $\mathcal{P}(\aleph) \not\subseteq \aleph$ что противоречит условию 2. Следовательно, в любой непротиворечивой алгебре, допускающей универсальный тип множеств, не может одновременно выполняться полная форма теоремы Кантора и возможность тотальной аннотации порядковым типом. Теорема доказана.

В данном контексте множество всех подмножеств понимается в чисто комбинаторном смысле - через определение степени как множества всех покрытий одного множества другим, поскольку мы возвращаемся к структурному подходу Кантора. Ключевым для нас является то обстоятельство, что уже в рамках работ Кантора можно было установить фундаментальную несовместимость: теория с глобальным тотальным порядком не может быть замкнутой, а замкнутая теория не допускает глобального порядка. Данное наблюдение предвосхищает теорему Спекера и позволяет заключить, что интерпретация $2^\omega = \omega$ противоречива в рамках построенной Кантором теории. Мы не можем одновременно замкнут алгебру и обеспечить тотальный порядок в ней. Множество $\omega$ не может быть одновременно замкнутым и вполне упорядоченным. Если $\omega$ представляет собой универсум всех конечных чисел как вполне упорядоченный объект, то по теореме Кантора все возможные комбинации конечных чисел выходят за его пределы. Чтобы описывать $\omega$ как законченный объект, необходимо отказаться от возможности самореференции - например, через стратификацию, как это реализовано в NF. Однако в этом случае мы утрачиваем полную упорядоченность, поскольку отношение принадлежности $x \in y$ становится не всюду определённым. В общепринятой интерпретации это противоречие игнорируется. Поскольку мы работаем в формальной теории, ограниченной фиксированным набором аксиом, полноценное противоречие вывести не удаётся. Для его формулировки требуются дополнительные предположения, частично восстанавливающие онтологию Кантора. По своей сути, теоремы Гёделя о неполноте и Тарского о неопределимости истины представляют собой доказательство противоречивости выбранной интерпретации, в определённом смысле, хотя мы придаём им несколько иную трактовку. В результате возникает новая иерархия мета-уровней теорий. Антиномия Рассела - это парадокс совмещения интенсионального и экстенсионального описания: мы сопоставляем предикат и множество, сводя логику к операциям над множествами. Такое описание допускает опосредованную самореференцию - предикат может быть применён к самому себе через множество, представляющее данный предикат. Гёделева нумерация представляет собой в точности тот же механизм: мы пытаемся экстенсионально описать логику с помощью чисел и сталкиваемся с невозможностью такого описания вследствие диагонального аргумента, замыкающего выражение на его собственный код. Иными словами, в обоих случаях мы пытаемся интерпретировать логику теории объектами самой теории и в обоих случаях терпим неудачу. Решением становится введение иерархии типов - кумулятивной или мета-уровневой. Мы в буквальном смысле совершаем одну и ту же ошибку многократно. Из этого следует вывод, что наша интерпретация неполна. Теория не может доказать свою собственную непротиворечивость, то есть не способна полностью интерпретировать саму себя. Истина и доказуемость в ней - не одно и то же. Таким образом, мы по какой-то причине продолжаем верить в возможность одновременного описания замкнутого и вполне упорядоченного объекта сильнее, чем в эквивалентность истины и доказуемости.

На современном языке описанную выше теорему можно переформулировать следующим образом:

Теорема 13 (современная формулировка): {\sl не существует первопорядковой теории $T$, удовлетворяющей одновременно следующим условиям:
\begin{enumerate}
    \item в $T$ выполняется аксиома выбора (т. е. всякое множество вполне упорядочиваемо);
    \item в $T$ существует множество $V$, замкнутое относительно операции взятия степени: $\mathcal{P}(V) \subseteq V$;
    \item $T$ непротиворечива (консистентна).
\end{enumerate}}

При формальной записи такой теории $T$ требуется проявлять осторожность в точной формулировке её свойств, однако это не представляет принципиальной трудности. В частности, следует отметить, что аксиома выбора несовместима с существованием универсального множества $V = \{ x \mid x = x \}$ в стандартном смысле, поскольку уже само наличие такого $V$ приводит к противоречию с теоремой Кантора при наличии тотальной упорядоченности.

Рассмотрим, что формально происходит, когда мы игнорируем теорему Кантора в упорядоченном универсуме и принимаем равенство $n^{\omega} = \omega$. Такое соотношение позволяет нам считать, что домен формул можно взаимно однозначно сопоставить домену натуральных чисел, а значит мы можем использовать кодирование Гёделя как преобразование. Однако, теорема Гёделя в действительности не требует полноценного кодирования всех формул, для неё нужна довольно ограниченная группа формул. Поэтому давайте обобщим наше наблюдение. Для того, чтобы уметь кодировать хоть что-нибудь кроме объектов домена самой теории, нам нужно иметь хотя бы одно обозначение отличное от обозначений всех объектов домена. Диагональная конструкция Кантора обеспечивает такое различие посредством внешнего сложение для каждого элемента домена $\cdot + 1$, таким образом домен кодирования должен быть не меньше $\omega + 1 > \omega$ - включать число $\omega$ отличное от всех объектов домена. Отрицая теорему Кантора и внешние операции мы утверждаем рефлексивность нашего домена, его замкнутость относительно операции добавления элементов: $1 + \omega = \omega$, фактически предполагая, что для описания семантики домена, нам достаточно только его самого. По диагональной лемме, для любой формулы $\varphi(x)$ с одной свободной переменной существует замкнутая формула $G$ такая, что в теории доказуемо $G \Leftrightarrow \varphi(\ulcorner G \urcorner)$ \cite{diagonal-lemma}. Отсюда немедленно следует невозможность определения предиката истины внутри самой теории: если бы существовал предикат $P_{truth}$, то мы могли бы построить $G$ такое, что $G \leftrightarrow \neg P_{truth}(\ulcorner G \urcorner)$, что приводит к противоречию. Применительно к аксиоматической теории множеств $Th$ с фиксированной стандартной моделью $\mathcal{M}$ можно определить множество:
\begin{align}
\begin{split}
\omega_{truth} &= \{ n \in \omega \mid \exists \varphi \; (n = \ulcorner \varphi \urcorner \; \wedge \; \mathcal{M} \models \varphi) \}
\end{split}
\end{align}

Здесь возникает принципиальная развилка. Либо выражение $\omega_{truth} \in \mathcal{P}(\omega)$ справедливо внутри теории, то есть множество $\omega_{truth}$ являются полноценными объектами теории. Либо отношение принадлежности $\in$ внутри теории и отношение $\dot{\in}$ во внешней метатеории - это два различных отношения, и элемент множества всех подмножеств с внешней точки зрения не обязательно является таковым с внутренней. Это хорошо известное следствие теоремы Лёвенгейма–Скулема об относительности понятия множество всех подмножеств: отношение принадлежности представляет собой всего лишь совокупность пар и может различаться не только в зависимости от мощности модели, но и от роли (внутренней или внешней), которую играет теория по отношению к модели. В первом случае, $\omega_{truth}$ - объект внутренней теории, мы можем построить новую теорию, добавив константу, обозначающие это множество. Однако, в силу теорем Гёделя и Тарского, формула $x \in \omega_{truth}$ не может служить внутренним предикатом истины. Она остаётся таковой лишь во внешней теории, где $x \dot{\in} \omega_{truth}$ действительно выражает соответствующее понятие. То есть теория не может интерпретировать собственный предикат истины даже если он в ней есть. Таким образом, в любом случае внутреннее $\in$ и внешнее $\dot{\in}$ представления существенно различаются. Именно поэтому $\omega$-непротиворечивость столь отчётливо подтверждает теорему Кантора для упорядоченных множеств порядковых типов $\omega$ и $2^\omega$. Гёдель сделал невероятную вещь, по сути усилив теорему Кантора (поскольку диагональная лемма фиксирует любую попытку внешних семантических обобщений): для описания математической реальности нам недостаточно исключительно внутреннего описания. Однако, смешение ранних интуитивных взглядов Кантора с его зрелой, более последовательной позицией, чётко разделявшей неупорядоченные и упорядоченные системы, привело к тому, что мы восприняли теорему Гёделя как ограничение всей математики, а не только внутреннего описания. Показательно, что если применить классическую интуицию и определить, например:
\begin{align}
\begin{split}
P_{truth}(\ulcorner G \urcorner) = \exists i : (i = \ulcorner H \urcorner \wedge H \Rightarrow G)
\end{split}
\end{align}

то есть считать утверждение истинным, если существует число, кодирующее формулу, из которой оно следует, то мы немедленно приходим к выводу о противоречивости теории. Противоречие возникает именно из предположения смешивающего внутренний домен чисел и внешний домен формул, что нарушает обобщённую теорему Кантора - $\omega = \omega + 1 = 2^\omega$. Мы вынуждены заключить, что внутри теории невозможно говорить о доказуемости в полном смысле. При этом доказательство - это всего лишь конечная комбинаторная конструкция из объектов теории, цепочка символов, о существовании которой теория может высказываться. Таким образом, противоречие возникает в процессе наделения комбинаторных объектов семантикой. Гёдель показал, что экспоненциал в общем случае не требуется, что проблемна любая попытка включить в синтаксис объекты внешней семантики, будь то кванторы или предикат истинности. Иными словами, наделение комбинаторных объектов смыслом может приводить к противоречиям, или, что то же самое, не может существовать замкнутой семантики теории: мы не способны непротиворечиво придать смысл всем комбинаторным объектам теории одновременно. Чтобы избежать противоречия, необходимо явно разделить внутренний и внешний уровни и признать, что не существует взаимно однозначного соответствия связывающего $H$ как комбинаторный объект внутренней теории и $H$ как цепочку символов, доказывающую $G$ во внешней теории. Но тогда мы имеем дело с разными алгебрами - внутренней и внешней, - каждая из которых обладает своей собственной индукцией и своей структурой бесконечных объектов, в частности, домен формул - $\omega + 1$ или $n^{\omega}$, должен занимать внешнюю позицию по отношению к объекту $\omega$. Это в точности то разделение, которое изначально предложил Кантор.

В определённой степени это вполне очевидно. Если бы внешняя метатеория была достаточно сильна и описывала тот же объект $\omega$ как замкнутую алгебру (т. е. как арифметику Пеано с импредикативным отождествлением $\omega = 1 + \omega = 2^\omega$), то в ней были бы доказуемы теоремы Гёделя и Тарского относительно этого объекта $\omega$. Иными словами, в такой метатеории понятия истины и доказуемости для исходной теории также оказались бы неопределимыми. Однако в этом случае необходимая иерархия метатеорий не возникает, а значит, метатеория не могла бы обладать требуемой силой. Ей пришлось бы отказаться от некоторых ключевых предположений - в частности, от рефлексивности домена. Тем не менее, мы обычно считаем, что метатеории образуют иерархию в духе Тарского, тем самым неявно принимая точку зрения, согласно которой отношение принадлежности существенно зависит от того, в какой теории - внутренней или внешней - мы находимся. Формально, современная интерпретация рассматривает объект  $2^\omega$ как внутренний, а в качестве внешнего использует $2^{\aleph_0}$, который не тождественен $2^{\omega}$. Это позволяет считать, что арифметика Пеано первого порядка просто не видит бесконечных подмножеств натуральных чисел, которые могут выражать истинность и доказуемость. Тем не менее, уже в арифметике второго порядка или теории множеств, такой объект полностью легитимен, но теория оказывается неспособной интерпретировать его через отношение принадлежности. Получается, что аксиоматические теории в их современной интерпретации определяют некоторые объекты, однако эти объекты экстенсионально не универсальны и не могут считаться множествами в полном смысле слова. Такие теории, следовательно, не решают задачу математического описания множественных сущностей в строго экстенсиональной форме. Как только мы требуем универсальности экстенсионального описания - то есть гарантии, что некоторый объект является элементом другого объекта независимо от контекста, - теории становятся противоречивыми, воспроизводя парадоксы Рассела или Берри. Подчеркнём, что метатеория может быть полностью идентична описываемой теории в смысле мощности описываемого универсума. Тем не менее, между внутренним отношением принадлежности ($\in$) и внешним ($\dot{\in}$) всегда сохраняется различие: множества внешней теории структурно отличаются от множеств внутренней теории. Это не приводит к прямому противоречию, ведь у нас остаётся возможность интерпретировать различие как обусловленное ролью, которую играет теория по отношению к рассматриваемому множеству. Тем не менее, мы вновь приходим к разделению внутреннего и внешнего описания, однако это разделение значительно менее изящно и концептуально проработано, чем то, которое предлагал Кантор. Если мы отказываемся от предположения рефлексивности домена и используем теорему Кантора в рамках внешней упорядоченной алгебры, чтобы установить разницу собственно домена и домена кодирования, то становимся способны утверждать, что уже арифметика первого порядка не видит всех своих формул. Это не просто добавляет согласованности между арифметиками первого и второго порядков, но и снимает принципиальное ограничение на способность теории использовать механизмы естественного языка для обхода петель самореференции, в том числе семантических, используемых в теоремах Гёделя и Тарского. Мы будем называть такое ограничение классической интерпретации {\it устойчивой} неполнотой. Мы в буквальном смысле встаём на одни и те же грабли дважды, один раз для того, чтобы построить упорядоченный универсум как иерархию, второй раз, чтобы построить иерархию Тарского, даже не задумываясь, что механизм у этого расслоения один и связан он с внешними функциональными отношениями.

Вследствие теорем Гёделя и Тарского, мы приходим к выводу, что понятия истины и доказуемости становятся неопределимыми внутри самой теории, равно как и все производные от них понятия - предикат, противоречие, непротиворечивость и так далее. Иными словами, уже для корректной формулировки теории мы вынуждены предварительно построить иерархию метатеорий, которая позволяла бы строго формализовать понятия противоречия и предиката. Это создаёт порочный круг: ключевые семантические понятия (истина и доказуемость) зависят от теории, которая, в свою очередь, строится с их использованием. Более точно, порочный круг возникает вокруг понятий, связанных с бесконечностью. Мы не можем определить понятия $\omega$-истинности и $\omega$-доказуемости внутри теории: для осмысленного утверждения свойств бесконечных объектов нам неизбежно требуется внешняя метатеория. Эта ситуация весьма напоминает запрет на самопринадлежность в новых основаниях Куайна. Наши теории не могут осмысленно говорить об истинности и доказуемости применительно к самим этим понятиям и их производным, что приводит к аналогичной стратификации через иерархию мета-уровней. На практике это означает, что внутри теории истинность предикатов, относящихся к бесконечным объектам, остаётся неопределённой. В частности, утверждение $\omega \in \omega$ не может быть оценено как истинное или ложное средствами самой теории. В онтологии Кантора ситуация выглядит иначе. Выражение $\omega - 1 \in \omega$ не лишено смысла, хотя оба множества имеют порядковый тип $\omega$. Аналогичным образом обстоит дело в более детализированных структурных аннотациях, если имеется множество $X$ индексного типа $\dot{\omega}$, то ему соответствует {\it индексал} - множество канонических индексных множеств, определяемых аналогично ординалам. Удалив один элемент $x$ из $X$, мы получаем множество $X \setminus \{x\}$ с индексалом $\dot{\omega} - 1$. Таким образом, с точки зрения индексной алгебры $\dot{\omega} - 1 \in \dot{\omega}$ (или $\dot{\omega} - 1 < \dot{\omega}$), хотя порядковые типы обоих индексалов тождественны. Поскольку $\omega - 1 = \omega$, мы получаем честную неподвижную точку $\omega \in \omega$ в множественном (индексном) смысле, при этом арифметика порядковых чисел остаётся непротиворечивой: $\omega \nless \omega$. Современные аксиоматические теории множеств сохраняют рудименты этой исходной онтологии - в частности, равенство $\omega - 1 = \bigcup \omega = \omega$. С внешней точки зрения утверждение $\omega \notin \omega$ определено и истинно, тогда как внутри теории возникает запрет на присвоение истинностного значения выражению $\omega \notin \omega$. Это фактически воспроизводит механизм стратификации, известный по NF. Следствием является то, что аксиома выбора не может иметь полноценного смысла внутри теории (иначе мы смогли бы интерпретировать собственный объект истины): она осмыслена лишь как внешнее утверждение относительно внутренней теории. Именно поэтому аксиома выбора порождает такое количество парадоксальных и диссонантных результатов - она требует постоянного неявного переключения контекста между внутренней и внешней теорией.

Судя по всему, Гёдель и Тарский сделали всё возможное в рамках существовавшего тогда формального аппарата, чтобы продемонстрировать, что современная математика выбрала далеко не самую удачную интерпретацию идей Кантора, однако термин неполнота плохо передаёт глубинный смысл их открытия. В действительности неполные теории оказываются зависимыми от утверждений, которые невозможно ни доказать, ни опровергнуть внутри самой теории. Истина по своей природе процессуальна: даже если мы допускаем существование платоновских истин самих по себе, их интерпретация и принятие нами неизбежно носят процессуальный характер. Исходная идея аксиоматического метода состояла в том, чтобы принять на веру лишь те утверждения, которые непосредственно проверяемы - воспроизводимы в опыте, в реальности или в ясном воображении (допускают процедуру верификации, что уже порождает вопросы к аксиоме бесконечности). Все остальные истины должны были опираться на доказательства, то есть на процедуры подтверждения. Называя истины платоновскими и помещая их в мир идей, мы не решаем проблему, а лишь оправдываем акт принятия недоказуемых утверждений на веру с философских позиций. Таким образом, подобные теории остаются непротиворечивыми лишь при условии, что мы сознательно принимаем определённые недоказуемые утверждения. Если же мы откажемся от этого акта веры, теории становятся противоречивыми. Именно поэтому корректнее говорить об {\it условной непротиворечивости} или, используя терминологию Лёйтзена Брауэра, о {\it недостоверности} таких теорий. С классической точки зрения отсутствие доказательства означает независимость утверждения от аксиом теории. Однако если утверждение независимо, оно может быть как истинным, так и ложным. В случае его ложности теория немедленно становится противоречивой, следовательно, мы не можем утверждать консистентность теории. В силу закона исключённого третьего из этого факта мы вынуждены сделать заключение о её неконсистентности. Теория не может быть одновременно консистентной и неконсистентной в рамках булевой логики. Чтобы сохранить её консистентность, необходимо либо включить в теорию дополнительные аксиомы, позволяющие считать внешние утверждения истинными, либо признать, что булевой логики недостаточно для описания самих математических теорий и требуется перейти к трёхзначной логике. Кроме того, возникает путаница между интуитивным пониманием неполноты (возможностью сформулировать на языке теории утверждения, независимые от неё) и её формальным определением. Полнота в системах функций, логических операций или метрических пространствах означает достаточность системы для описания всех выражений или достижения всех её элементов. В частности, часто в качестве иллюстрации приводят теорему Гудстейна \cite{goodstein}, независимую от аксиоматики Пеано. Однако этот случай относится именно к интуитивному пониманию неполноты, поскольку ложность теоремы Гудстейна не влечёт противоречивости арифметики Пеано. Это, пожалуй, наиболее яркая демонстрация правоты позднего Витгенштейна: даже математика оказывается утончённой языковой игрой. Термин неполнота, маскирующий условную непротиворечивость, превращается в нейтральное и вполне допустимое понятие, которое служит защитой для определённой формы веры, фактически отрицающей один из ключевых критериев научности - критерий фальсифицируемости. Что позволяет закрыть глаза на то, что классическая интерпретация является не просто недостоверной, а устойчиво недостоверной.

В свете изложенных выше фактов становится очевидным, почему именно идеи Кантора - пусть и в несколько искажённой форме - легли в основу теории, ставшей фундаментом современной математики. Кантору удалось избежать всех основных парадоксов, возникавших впоследствии: от парадокса Бурали-Форти до парадокса Рассела. Более того, его подход фактически предотвращает возникновение этих парадоксов. Парадокс Бурали-Форти, в частности, не возникает в онтологии Кантора. Это демонстрирует следующая теорема, которую по сути практически сформулировал сам Кантор:

Теорема 14: {\sl Совокупность всех порядковых чисел $\Omega$ незамкнута (потенциально бесконечна).}

Доказательство: 
Предположим противное: совокупность $\Omega$ замкнута, то есть является множеством. Тогда $\Omega$ вполне упорядочена и обладает наименьшим элементом. Как множество, она должна иметь некоторый порядковый тип $\gamma$. Этот тип $\gamma$ превосходил бы все порядковые числа, входящие в $\Omega$. Но поскольку $\gamma$ само является порядковым числом, оно должно принадлежать $\Omega$, откуда следовало бы $\gamma < \gamma$. Полученное противоречие показывает, что предположение о замкнутости $\Omega$ ложно. Таким образом, совокупность всех порядковых чисел потенциально бесконечна. Теорема доказана.

Иными словами, в онтологии Кантора парадокс Бурали-Форти принципиально невозможен, поскольку Кантор работал исключительно с конструируемыми объектами и не предполагал схемы неограниченного усложнения, которая появилась у Фреге. Аналогичным образом теорема Кантора решает проблему парадокса Рассела, естественным образом порождая иерархию при внешнем упорядоченном рассмотрении множеств. Ключевой идеей Кантора было то, что числа представляют собой фундаментальные эмерджентные свойства алгебраических структур, именно эту концепцию он противопоставлял господствовавшей традиции об универсальной природе чисел, которую развивал Фреге. И именно такой подход позволяет обойти все противоречия связанные с самореферентными конструкциями. Кантор сосредоточил внимание на одном из наиболее сложных объектов - порядковых числах, в которых удивительным образом сплетаются внутренняя алгебра кардинальных чисел и внешняя алгебра индексных чисел. Он гениально разделил арифметику на левую и правую, подчёркивая, что $1 + x \neq x + 1$ и $2 * x \neq x * 2$. Его последователи, к сожалению, часто упускали смысл этого фундаментального различия. Левая арифметика рефлексивна: $\omega$ является неподвижной точкой. Правая арифметика таковой не является. Кантор сумел совместить внутреннюю и внешнюю алгебры в единой системе порядковых чисел. При этом сам объект $\omega$ возникает только во внешней алгебре; внутренняя алгебра не способна его породить. Для внутренней алгебры $\omega$ выступает внешней неподвижной точкой - подобно тому, как предел $0$ возникает для функции $f(n) = 1/2^n$, хотя $0 = 0/2$. Внутренняя алгебра порождает лишь $\aleph$, лишённый собственной морфологии. Следовательно, внутренняя алгебра не годится для доказательства утверждений относительно $\omega$, как объекта. Для неё действительно выполняется $-1 + \omega = \omega$, однако во внешней алгебре мы имеем $\omega - 1 < \omega$ и $\omega - 1 + 1 = \omega$. Правая алгебра представляет собой полноценную алгебру нормальной формы Кантора-Хессенберга, которая коммутативна благодаря внешней симметрии операций. Левая же алгебра просто не способна осмысленно говорить о свойствах $\omega$, поскольку для неё этот объект является внешним и в ней не возникает. Данный факт также естественным образом согласуется с тем, что $\omega$ представляет собой порядковое число в нормальной форме, которая в свою очередь непосредственно связана с внешней индексной алгеброй через нормальные операции Хессенберга.

К сожалению, Кантор не успел завершить свою работу. Обнаружение парадоксов в наивной теории множеств привело к поспешной аксиоматизации, которая сосредоточилась лишь на узком фрагменте его видения. Теория самого Кантора не только допускала сосуществование трёх независимых алгебр, но и делала их взаимодополняющими: каждая алгебра освещала другие с разных сторон. Это принципиально отличается от последующей конфронтации между различными аксиоматическими системами - ZFC, NBG, NF и альтернативными конструкциями, такими как сюрреальные числа и гросс-объекты. Наиболее детализированную аннотацию бесконечных множеств дают индексные числа (нестандартные натуральные числа, сюрреальные числа Конвея, гросс-объекты Сергеева и близкие к ним конструкции). Их операции строго внешние: любое нетривиальное объединение или другая операция порождает объект, выходящий за пределы исходных операндов. Именно поэтому индексные числа кристаллизуют трансфинитные объекты, наиболее близкие к естественному интуитивному пониманию числа, хотя и остаются классами эквивалентности множеств и не дают окончательного ответа на вопросы типа делится ли бесконечность нацело пополам. Промежуточное положение занимают порядковые числа Кантора. Их операции носят гибридный характер: левый операнд трактуется внутренне, а правый - внешне. Эта двойственность особенно ярко проявляется в экспоненциации: внешний предел задаётся через бесконечную сумму, тогда как внутренний - через супремум. В результате бесконечно большие и бесконечно малые величины, такие как $\omega - 1$, $\frac{\omega}{2}$ или $\log_2\omega$, оказываются спрятанными внутри классов эквивалентности и не выражаются напрямую. Наконец, наиболее абстрактной и универсальной является алгебра кардинальных чисел. Она разделяет бесконечность на слои регулярности. Конечная кардинальность завершается первым бесконечным кардиналом $\aleph_0$ - классом эквивалентности, который вмещает всё, что может быть получено из порядкового типа $\omega$, индексного типа $\dot{\omega}$, их экспоненциальных башен $\omega^\omega$, $\dot{\omega}\uparrow\uparrow\dot{\omega}$ и так далее, пока результат остаётся счётным. Как только появляется объект, выходящий за эти пределы - первый несчётный порядковый тип $\omega_1$ или соответствующее индексное число $\dot{\omega}_1$, - возникает следующий слой мощности $\aleph_1$. Каждая из трёх алгебр замкнута относительно своих собственных операций и обладает собственной, отличной от остальных формой трансфинитной индукции. Мы действительно можем варьировать, что именно для нас первично - бесконечность или индукция: определив свойства одного объекта, свойства другого становятся выводимыми. Однако оба объекта всегда остаются замкнутыми внутри одной алгебры. Смешение алгебр приводит логическим коллизиям и устойчивости условной непротиворечивости теории. Схематично взаимоотношение трёх алгебр изображено на Рис. \ref{fig:three-algebras}.

\begin{figure}[htbp]
\centering
\begin{tikzpicture}[
    node distance=1.8cm,
    every node/.style={font=\sffamily},
    darrow/.style={double, double distance=2pt, -{Latex[open]}, dashed, thick},
    arrow/.style={-{Latex}, thick},
    sing/.style={red!60!gray, dashed, ultra thick}
]

\node[font=\Large\bfseries] at (1,7.5) {Алгебраические структуры бесконечности};

\node (fincard) at (-4,5.8) [blue!70!gray, align=center] {0, 1, 2, };
\node (finord) at (-4,3.8) [teal!70!green, align=center] {0, 1, 2, };
\node (finind) at (-4,1.8) [violet!70!lightgray, align=center] {0, 1, 2, };

\draw[sing] (-2.1,1) -- (-2.1,6.6);
\node[red!60!gray, font=\bfseries] at (-2,6.9) {Сингулярность};

\node[blue!70!gray, font=\bfseries] (card) at (2,6.4) {Кардинальная алгебра};
\draw[darrow, blue!80!gray] (-3.3,5.8) -- (-2.1,5.8);
\draw[arrow, blue!80!gray] (2.5,5.8) -- (-2.1,5.8);
\node[blue!70!gray] (o2) at (2.8,5.8) {$\aleph_0$};
\draw[arrow, blue!80!gray] (3.05,5.8) -- (7.65,5.8);
\node[blue!70!gray] (c1) at (8,5.8) {$,\; \aleph_1$};

\node[teal!70!green, font=\bfseries] (ord) at (2,4.4) {Ординальная алгебра};
\draw[darrow, teal!80!green] (-3.3,3.8) -- (-2.1,3.8);
\draw[arrow, teal!80!green, dashed] (-0.5,5) -- (-2.1,5);
\node[teal!70!green, font=\small\bfseries] (o1) at (1.25,5) {изоморфизм с $\omega$};
\draw[arrow, teal!80!green, dashed] (3,5) -- (4.5,5);
\draw[arrow, teal!80!green] (-0.1,3.8) -- (-2.1,3.8);
\node[teal!70!green] (o2) at (0.15,3.8) {$\omega$};
\draw[arrow, teal!80!green] (0.4,3.8) -- (2.3,3.8);
\node[teal!70!green] (o3) at (5.32,3.8) {$,\; \omega+1,\; \dots,\; 2^\omega,\; \dots,\; \omega^\omega,\; \dots,\; \omega_1$};

\node[violet!70!lightgray, font=\bfseries] (ind) at (2,2.4) {Индексная алгебра};
\draw[darrow, violet!80!lightgray] (-3.3,1.8) -- (-2.1,1.8);
\draw[darrow, violet!80!lightgray] (-0.9,1.8) -- (-2.1,1.8);
\node[violet!70!lightgray] (i1) at (3.8,1.8) {$, \dot{\omega} - 2,\; \dot{\omega} - 1,\; \dot{\omega},\; \dot{\omega} + 1,\; \dots,\; 2^{\dot{\omega}},\; \dots,\; \dot{\omega}^{\dot{\omega}},\; \dots,\; \dot{\omega}_1$};

\draw[thick, gray!70] (5.1,5.3) -- (5.1,4.2);
\draw[thick, gray!70] (5.1,3.3) -- (5.1,2.2);

\node[align=center, font=\small\itshape, gray!80] at (6.8,4.7) {Интерпретации\\бесконечности};
\node[align=center, font=\small\itshape, gray!80] at (6.8,2.7) {Классы\\эквивалентности};

\node[rotate=90, blue!70!gray, font=\scriptsize\itshape] at (-5.2,5.8) {Кардиналы};
\node[rotate=90, teal!70!green, font=\scriptsize\itshape] at (-5.2,3.8) {Ординалы};
\node[rotate=90, violet!70!lightgray, font=\scriptsize\itshape] at (-5.2,1.8) {Индексалы};

\end{tikzpicture}
\caption{Взаимосвязи между тремя алгебрами описания бесконечности в онтологии Кантора}
\label{fig:three-algebras}
\end{figure}

Настоящая работа, посвящёна критическому переосмыслению границ применимости теоремы Кантора о несчётности континуума и связанных с ней фундаментальных вопросов теории множеств. Проведённый анализ показывает, что классическое диагональное доказательство (в его современном обобщённом виде, включая версию Хессенберга) является логически безупречным и абсолютно верным в контексте упорядоченных универсумов, в частности, внутри аксиоматических систем ZFC и NBG. Оно надёжно устанавливает, что никакое множество не может быть равномощным множеству всех своих подмножеств при условии, что мы работаем в рамках вполне упорядоченной кумулятивной иерархии и принимаем стандартные аксиомы, обеспечивающие предикативное противоречие экспоненциальной самореференции. Однако теорема Кантора не является универсальным инструментом исследования самой структуры бесконечности. Её доказательство требует предварительного принятия довольно сильных утверждений о природе бесконечного, в частности, универсальности позволяющей одновременно сопоставлять множество объектов и тотальности, гарантирующей способность теории описывать все элементы структуры. При изменении этих исходных предположений (например, при отказе от экспоненциального ограничения рефлексивности, постулировании свойств сингулярности между областями регулярности или переходе к чистой кардинальной алгебре Кантора) построение теряет консистентность или перестаёт приводить к тому же выводу. Таким образом, теорема описывает не природу бесконечности как таковой, а свойства конкретной алгебраической структуры, в которой она применяется. Одной из основных методологических трудностей современной теории множеств остаётся ошибка реификации - отождествление алгебраической аннотации (порядковых типов, ординалов как канонических представителей) с самими исследуемыми объектами. Хотя эта ошибка была осознана ещё Джоном фон Нейманом и не является критической для внутренней непротиворечивости ZFC и NBG, она существенно затрудняет ясный анализ расхождений между поздней онтологией Кантора и современной картиной. В архитектуре современных теорий множеств были выявлены существенные дефекты, связанные с аксиомой бесконечности, которая не удовлетворяет основным формальным требованиям к аксиомам. Она не является атомарной, чистой, однозначной и независимой, что ещё сильнее затрудняет точное концептуальное моделирование теории. Это приводит сразу к двум методологическим формальным ошибкам преждевременного обобщения внутреннего функционального описания и индуктивной морфологии бесконечности на весь универсум теории (логическая ошибка ложной дихотомии). В действительности, внутреннее функциональное описание представляет собой лишь фрагмент теории, связанный с рефлексивностью, а индуктивная морфология бесконечности - фрагмент, обеспечивающий возможность представления всех объектов теории в морфологии вполне упорядочения. Эти ошибки приводят к неоднозначности определения отношения принадлежности, которое очень точно было уловлено интуиционистыми. Дело в том, что потенциальные альтернативы не исключают друг друга, пока мы не сделаем одну из них реальностью. Это прямое нарушение закона исключённого третьего, на которое указывал Брауэр. Игнорирование данного факта, приводит к рассинхронизации синтаксиса и семантики, противоречие может быть синтаксически выводим в теории, но семантически бессмысленно поскольку вывод утверждения и его отрицания упирается в разные семантические корни, формирующие параллельные непересекающиеся поддеревья вывода внутри теории. Принятие методологического предположения Френкеля (что в универсуме существуют лишь объекты, непосредственно необходимые по аксиомам, делая альтернативные множества нелегитимными) способно привести теорию к консенсусу мейнстрима, но имеет сложности с аргументированным обоснованием. С формальной точки зрения аксиому бесконечности лучше всего заменить на несколько атомарных и независимых утверждений, чтобы иметь возможность чётко прослеживать границы применимости доказательств.

Тем не менее, даже внутри ZFC и NBG в рамках современной интерпретации отождествление $\omega = 2^\omega$, опровергающем теорему Кантора, остаётся недостаточно обоснованным. В рамках этой интерпретации неявное соглашение, универсиализирующее единую функциональную структуру и индуктивную морфологию, корректно редуцирует синтаксическое выражение внутреннего содержания бесконечного объекта $\omega$ к семантике непосредственно самого этого объекта: $\bigcup \omega = \omega$. Однако некритически отказывает в такой редукции на мультипликативных уровнях воспроизводства аддитивной структуры: $\omega^{\bigcup \omega} \neq \omega^{\omega}$. Это приводит к формированию интерпретационного порочного круга связанного с концепцией наследственно конечных подмножеств степенных множеств (конечных областей определения или конечных носителей). Кардинальные свойства этой концепции следуют из теорему о счётности счётного объединения счётных множеств, доказательство которой в свою очередь опирается на данную концепцию. Такое некритическое решение требует дополнительного методологического допущения - распространение внутренней индукции (арифметики Пеано) на $\omega$ как на внешний объект. Факт того, что $\omega$ является точкой разрыва предикативности предоставляет нам возможность выбора какие именно инварианты мы хотим сделать транзитивными при переходе через сингулярность разделяющую индуктивную область регулярности и объект $\omega$. Это оказывается относительно консистентно внутри самой теории, однако при попытке описывать математический анализ в данной рамке возникают явно выраженные коллизии таких промежуточных интерпретационных объектов с аксиомой непрерывности. Теории такого рода оказываются фундаментально слабыми, не допуская усиления предикативности, что обуславливает их устойчивую условную непротиворечивость в фундаментальном смысле: их консистентность зависит от внешних утверждений, которые не могут быть ни доказаны, ни опровергнуты средствами самой теории, при этом эффект семантических самореферентных петель оказывается практически не устраним. Неявные допущения формальных аксиоматических систем образуют главную линию раздела между современной мейнстрим-теорией и альтернативными взглядами (включая поздние построения Кантора, системы типа NF, сюрреальные числа Конвея и гросс-объекты). Простая формализация данных предположений позволяет сделать всю математику целостной и внутренне согласованной: мейнстрим сохраняет свою мощь и удобство работы с упорядоченными структурами унифицирующими понятие числа (в духе Фреге), а альтернативные направления рассматривающие числа как эмерджентные свойства алгебраических структур (в духе Кантора) получают законное место без взаимного отрицания. Признание теоремы Кантора правомерной именно для упорядоченных объектов, в сочетании с легитимностью альтернативных объектов и явным разделением внутренней (кардинальной, потенциальной) и внешней (ординальной/индексной, актуальной) алгебр, открывает двери к более богатой и гибкой картине. Что естественным образом ведёт в мир нумерости, сюрреальных чисел и гросс-множеств, где бесконечно большие и бесконечно малые величины интегрируются значительно более гармонично, чем в нестандартном анализе. В свою очередь блокировка обобщённой теоремой Кантора синтаксических сомореферентных конструкций открывает двери для расширений теории, способных избавить её от условной непротиворечивости, переводя её в разряд строгих, хотя и контекстно-зависимых математических построений. Таким образом, кризис оснований математики, переживаемый на протяжении более века, оказывается не следствием внутренней ограниченности математики, а результатом недостаточно аккуратного обращения с границами применимости ключевых понятий и методологических допущений. Переосмысление этих границ не разрушает существующее здание математики, а лишь уточняет его фундамент, делая его более прочным и открытым для дальнейшего развития. Теория множеств перестаёт претендовать на универсальное описание всех аспектов бесконечности и становится тем, чем она и должна быть - мощным, но специализированным инструментом, требующим ясного указания контекста и алгебраической структуры, в рамках которой он применяется.

\printbibliography

 \end{document}